% concatenated from Singular_fractional_p-Laplacian_final.tex
\documentclass[11pt,reqno]{amsart}

\usepackage{amsmath,amssymb,amsthm,mathtools,mathrsfs,esint}
\usepackage[margin=1.10in]{geometry}
\usepackage{hyperref}

\hypersetup{
  %hidelinks,
  bookmarksnumbered=true,
  bookmarksopen=true,
  pdfauthor={Chao Zhang},
  pdftitle={Gradient Hölder regularity for singular fractional p-Laplace equations},
  pdfsubject={Singular fractional p-Laplacian and gradient regularity}
}

\numberwithin{equation}{section}

\newcommand{\R}{\mathbb R}
\newcommand{\Jp}{\mathcal J_p}
\newcommand{\dd}{\,d}
\newcommand{\Lip}{\operatorname{Lip}}
\newcommand{\osc}{\operatorname*{osc}}

\newtheorem{theorem}{Theorem}[section]
\newtheorem{proposition}[theorem]{Proposition}
\newtheorem{lemma}[theorem]{Lemma}
\newtheorem{corollary}[theorem]{Corollary}
\theoremstyle{remark}
\newtheorem{remark}[theorem]{Remark}

\makeatletter
\@namedef{subjclassname@2020}{%
	\textup{2020} Mathematics Subject Classification}
\makeatother


\title[Gradient regularity for singular fractional equations]
{Gradient H\"older regularity for singular fractional
\(p\)-Laplace equations}

\author[C. Zhang]{Chao Zhang}
\address{Chao Zhang  \hfill\break School of Mathematics and Institute for Advanced Study in
Mathematics, Harbin Institute of Technology, Harbin 150001, China}
\email{czhangmath@hit.edu.cn}

\subjclass[2020]{35R11, 35B65, 35J70}
\keywords{fractional \(p\)-Laplacian, singular equations, gradient
regularity, Liouville theorem, improvement of flatness, nonlocal tails}

\begin{document}
\raggedbottom

\begin{abstract}
Let \(n\ge2\), \(1<p<2\), \(0<s<1\), and \(sp>p-1\).  We prove that
every globally bounded fractional \(p\)-harmonic function is locally
\(C^{1,\alpha}\) for some \(\alpha=\alpha(n,p,s)>0\).  This settles the
open problem of interior gradient H\"older regularity in the singular
range throughout the natural first-order regime \(sp>p-1\).  The proof
combines an affine-invariant improvement-of-flatness argument with a Liouville
theorem for globally Lipschitz entire solutions.  In the large-slope
regime, the shifted Bregman energies converge to an anisotropic stable
form of order \(sp-p+2>1\).  In the bounded-slope regime, the Liouville
theorem follows from rigidity of extremal secants, a recurrent blow-down
argument, and a directional Morrey--Kato estimate for the singular
linearized kernel.  An affine Campanato argument controls the variation
of the best affine approximations across scales.  These estimates yield
a scale-invariant decay of the affine excess and hence the local
\(C^{1,\alpha}\) estimate.
\end{abstract}

\maketitle

\section{Introduction and main result}
\label{sec:introduction}

Let \(n\ge2\), \(1<p<2\), and \(0<s<1\).  We write
\(\Jp(t):=|t|^{p-2}t\) for \(t\in\R\), and consider weak solutions of
\begin{equation}
 (-\Delta_p)^su=0\qquad\text{in }\Omega,
 \label{eq:introduction-equation}
\end{equation}
where, up to a positive normalization constant,
\[
 (-\Delta_p)^su(x)
 =\operatorname{P.V.}\int_{\R^n}
 \frac{\Jp(u(x)-u(y))}
 {|x-y|^{n+sp}}\dd y.
\]
We use the tail space
\[
 L^{p-1}_{sp}(\R^n)
 :=\left\{v\in L^{p-1}_{\rm loc}(\R^n):
 \int_{\R^n}\frac{|v(y)|^{p-1}}
 {(1+|y|)^{n+sp}}\dd y<\infty\right\}.
\]
Thus \(u\in W^{s,p}_{\rm loc}(\Omega)\cap L^{p-1}_{sp}(\R^n)\) is a
weak solution if
\begin{equation*}
 \iint_{\R^n\times\R^n}
 \frac{\Jp(u(x)-u(y))(\varphi(x)-\varphi(y))}
 {|x-y|^{n+sp}}\dd x\!\dd y=0
\end{equation*}
for every \(\varphi\in C_c^\infty(\Omega)\).

The regularity theory for fractional \(p\)-Laplace equations includes
local boundedness, Harnack inequalities, and interior and boundary
H\"older estimates
\cite{DiCastroKuusiPalatucci,DiCastroKuusiPalatucciHarnack,
IannizzottoMosconiSquassina,LindgrenViscosity}.  Fine boundary regularity
in the singular range was established in \cite{IannizzottoMosconi}.
Higher H\"older, Sobolev, and differentiability estimates were developed
in the superquadratic case in
\cite{BrascoLindgren,BrascoLindgrenSchikorra,
BoegeleinDuzaarLiaoMolicaBisciServadeiJFA,DieningNowak} and in the
singular case in \cite{GarainLindgren,DieningKimLeeNowak,
BoegeleinDuzaarLiaoMolicaBisciServadei}; the local Lipschitz estimate of
\cite{BiswasTopp} will be used below.  Nonlinear potential
estimates and self-improving properties were obtained in
\cite{KuusiMingioneSireMeasure,KuusiMingioneSireSelf}; see also
\cite{KuusiMingioneSireSurvey}.

At the linear level, gradient regularity and first-order potential
estimates for nonlocal equations of order larger than one were proved in
\cite{KuusiNowakSire}.  A nonlinear gradient-level potential theory is
available for operators with a globally Lipschitz and uniformly strongly
monotone constitutive function \cite{DieningKimLeeNowakGradient}.  That
structural class does not contain the singular flux
\(t\mapsto |t|^{p-2}t\) when \(1<p<2\).  Gradient H\"older estimates are
also known for mixed local--nonlocal problems with a leading local part
\cite{DeFilippisMingione}.

The interior H\"older continuity of the gradient for fractional
\(p\)-harmonic functions had remained a central open problem in the
regularity theory.  It was explicitly singled out in
\cite{BoegeleinDuzaarLiaoMolicaBisciServadeiHigher,
DieningKimLeeNowak,DieningNowak}; see also the discussion surrounding the
higher differentiability results in
\cite{BoegeleinDuzaarLiaoMolicaBisciServadei}.  Giovagnoli, Jesus, and
Silvestre \cite{GJS} proved interior \(C^{1,\alpha}\)
regularity in the range
\[
 2\le p<\frac{2}{1-s}.
\]
In particular, their result covers the entire
superquadratic part of the natural first-order regime \(sp>p-1\).

The singular case \(1<p<2\), however, is not covered by that argument;
indeed, \cite[p.~3]{GJS} explicitly leaves the range \(p<2\) open.
The available results in this range give weak differentiability,
arbitrarily high integrability and fractional differentiability of the
gradient, as well as almost Lipschitz continuity of the solution, but not H\"older
continuity of the gradient.  The obstruction is structural: after
subtraction of an affine function, the linearized coefficient is
singular at vanishing increments.  The purpose of the present paper is
to resolve precisely this remaining singular problem.  The theorem
below establishes interior gradient H\"older regularity throughout the
natural first-order regime \(sp>p-1\).

An earlier version of the present work obtained a partial result under
an additional restriction on the fractional order
\cite{ZhangPartial}.  The argument developed here is different at the
two points where that restriction entered: the large-slope limit is
treated at fixed order through a shift-uniform Orlicz compactness
estimate, while the bounded-slope Liouville theorem is closed by a
directional Morrey--Kato estimate for the crossed singularity.  These
two ingredients remove the additional restriction and yield the full
range \(sp>p-1\).  Thus the present paper supersedes the partial result
of \cite{ZhangPartial} at the level of the main theorem.

The main result is the following.

\begin{theorem}[Gradient H\"older regularity]
\label{thm:main}
Let \(n\ge2\), \(1<p<2\), and \(0<s<1\) satisfy
\begin{equation}
 sp>p-1.
 \label{eq:main-structural-range}
\end{equation}
There are \(\alpha_0=\alpha_0(n,p,s)>0\) and
\(C=C(n,p,s)<\infty\) such that every globally bounded weak solution of
\eqref{eq:introduction-equation} in \(B_2\) belongs to
\(C^{1,\alpha_0}(B_{1/2})\) and satisfies
\begin{equation}
 \|\nabla u\|_{L^\infty(B_{1/2})}
 +[\nabla u]_{C^{\alpha_0}(B_{1/2})}
 \le C\|u\|_{L^\infty(\R^n)}.
 \label{eq:fixed-order-gradient-holder-estimate}
\end{equation}
\end{theorem}

The condition \eqref{eq:main-structural-range} has a direct first-order
meaning.  The flux of a linearly growing function is integrable at
infinity exactly in this range.  In the large-slope compactification the
limiting stable operator has order
\[
 \tau_s=sp-p+2,
\]
and \eqref{eq:main-structural-range} is equivalent to \(\tau_s>1\).
The available gradient exponent in that regime is every number below
\(\tau_s-1=sp-p+1\).  The exponent in Theorem~\ref{thm:main} may be
smaller, because the bounded-slope estimate is obtained by a qualitative
Liouville compactness argument.

We outline the proof.  After subtracting an affine function, we measure
the singular energy by the corresponding Bregman remainder and distinguish
between bounded and unbounded normalized slopes.  In the latter case,
the normalized energies converge to an anisotropic quadratic stable
form, which yields a fixed-scale affine improvement.  The bounded-slope
case is reduced to a Bernstein-type Liouville theorem for globally
Lipschitz entire solutions.  The proof of this theorem combines rigidity
of an attained extremal secant, a recurrent dilation argument, a
nondegeneracy estimate for the rescaled amplitude, and a strong minimum
principle for the linearized equation.  A directional coarea estimate
places the exterior part of the singular linearized coefficient in a
Morrey--Kato class of codimension \(2-p\); the resulting deficit relative
to the regional order is \(sp\).

The Liouville theorem yields a compactness improvement in the
bounded-slope regime.  Possible growth of the best affine slopes is
controlled by an affine Liouville theorem in Campanato classes.  If these
slopes are unbounded, the large-slope estimate can be iterated from a
diverging scale to a fixed scale.  Consequently, both slope regimes
satisfy the same intrinsic excess estimate and the same affine-invariant
relative-tail bound.  The corresponding iteration is
\[
 r_{k+1}=\rho r_k,\qquad
 H_{k+1}=\rho^\beta H_k.
\]
After fixing the constants and the contraction factor, one chooses
\(0<\beta<sp-p+1\).  It follows that \(H_k\le Cr_k^\beta\) without an
annular renormalization constant.

The paper is organized as follows.
Section~\ref{sec:setting-contact} establishes the pointwise contact
lemma and the extremal-secant rigidity theorem.
Section~\ref{sec:large-slope-and-blowdown} develops the fixed-order
large-slope compactness and proves the affine blow-down theorem.
Section~\ref{sec:dilation-liouville} then proves the entire-solution
Liouville theorem.  Section~\ref{sec:common-recursion} derives the
bounded-slope compactness estimate, combines it with the large-slope
alternative in a single Campanato iteration, and proves the main theorem.

For clarity, the logical dependence is as follows.  The large-slope
compactness, large-slope improvement, noncritical expansion, and affine
blow-down principle in Section~\ref{sec:large-slope-and-blowdown}
depend only on the shifted energy estimates and the regularity theory
for linear stable operators.  The nonlinear Liouville theorem of
Section~\ref{sec:dilation-liouville} enters only the bounded-slope
compactness argument in Section~\ref{sec:common-recursion}.  The two
finite-scale estimates are then combined there, with one amplitude and
one relative-tail state.

\medskip
\noindent\textit{Notation and conventions.}
For a measurable set \(E\) of positive measure, we write
\(\fint_E g:=|E|^{-1}\int_E g\).  Unless a center is displayed,
\(B_r:=B_r(0)\).  For a function \(w\), we set
\(\delta w(x,y):=w(x)-w(y)\); the same convention is used for the
rescaled variables \(X,Y\).  The letter \(C\) may change from line to
line and depends only on the parameters indicated in the relevant
statement.  We write \(A\lesssim B\) if \(A\le CB\), and
\(A\asymp B\) if \(cB\le A\le CB\), for positive structural constants
\(c,C\).

\section{Pointwise contact and extremal-secant rigidity}
\label{sec:setting-contact}

\subsection{The pointwise contact lemma}
\label{subsec:pointwise-contact}

We consider the whole-space equation
\begin{equation}
 \operatorname{P.V.}\int_{\R^n}
 \frac{\Jp(U(x)-U(y))}{|x-y|^{n+sp}}\dd y=0
 \qquad\text{in }\R^n.
 \label{eq:equation}
\end{equation}
An \emph{entire weak solution} of \eqref{eq:equation} is a function
\[
 U\in W_{\rm loc}^{s,p}(\R^n)\cap L_{sp}^{p-1}(\R^n)
\]
that satisfies the weak formulation stated in
Section~\ref{sec:introduction} with \(\Omega=\R^n\).  When no ambiguity
can arise, we use \emph{entire solution} as shorthand for
\emph{entire weak solution}.

From now on we work under the standing structural assumption
\eqref{eq:main-structural-range}.  For globally Lipschitz functions,
this is exactly the condition ensuring absolute integrability at
infinity of the flux generated by linear growth.

We first establish the local contact lemma used in the
extremal-secant argument.  The equation is used directly in that
argument only through this lemma.

\begin{lemma}[Pointwise realization at a nonzero quadratic contact]
\label{lem:pointwise-contact}
Under the standing assumption \eqref{eq:main-structural-range}, let
$U$ be a globally Lipschitz entire weak solution of
\eqref{eq:equation}.  If, at $x_0$,
\[
 U(x_0+z)=U(x_0)+q\cdot z+O(|z|^2),
 \qquad q\ne0,
\]
then the principal value in \eqref{eq:equation} exists at $x_0$ and is
zero.
\end{lemma}

\begin{proof}
Write
\[
 U(x_0+z)=U(x_0)+q\cdot z+R(z),
 \qquad |R(z)|\le K|z|^2
\]
for \(|z|<r_0\), after decreasing \(r_0\) if necessary.  We first prove
that the flux differs from its affine part by an absolutely integrable
function.  For \(z=r\omega\), set
\[
 a=-r q\cdot\omega,\qquad b=-R(r\omega),\qquad
 \Gamma_r:=\{\omega\in\mathbb S^{n-1}:|q\cdot\omega|\le2Kr\}.
\]
Since \(q\ne0\), for all sufficiently small \(r\),
\[
 |\Gamma_r|\le C r.
\]
On \(\mathbb S^{n-1}\setminus\Gamma_r\) one has
\(|b|\le |a|/2\), and the mean-value estimate for \(\Jp\) gives
\[
 |\Jp(a+b)-\Jp(a)|
 \le C|a|^{p-2}|b|
 \le Cr^p|q\cdot\omega|^{p-2}.
\]
The angular factor is integrable because \(p-2>-1\).  On \(\Gamma_r\),
the global \((p-1)\)-H\"older continuity of \(\Jp\) yields
\[
 |\Jp(a+b)-\Jp(a)|
 \le C|b|^{p-1}\le Cr^{2(p-1)}.
\]
Consequently,
\[
 \int_{\mathbb S^{n-1}}
 |\Jp(-q\cdot r\omega-R(r\omega))
       -\Jp(-q\cdot r\omega)|\dd\omega
 \le C\bigl(r^p+r^{2p-1}\bigr)\le Cr^p.
\]
Polar integration therefore gives
\[
 \int_{B_{r_0}}
 \frac{|\Jp(U(x_0)-U(x_0+z))-\Jp(-q\cdot z)|}
 {|z|^{n+sp}}\dd z
 \le C\int_0^{r_0}r^{p-sp-1}\dd r<\infty,
\]
where the last inequality uses \(s<1\), equivalently \(sp<p\).

The far field is also absolutely integrable.  Indeed, if \(L\) is the
global Lipschitz constant of \(U\), then
\[
 \int_{\R^n\setminus B_{r_0}}
 \frac{|\Jp(U(x_0)-U(x_0+z))|}{|z|^{n+sp}}\dd z
 \le C+C L^{p-1}\int_1^\infty r^{p-sp-2}\dd r<\infty,
\]
because \(sp>p-1\).  Finally,
\[
 \int_{\varepsilon<|z|<r_0}
 \frac{\Jp(-q\cdot z)}{|z|^{n+sp}}\dd z=0
\]
by oddness.  The preceding estimates thus show that the principal value
at \(x_0\) exists; denote it by \(I\).

It remains to identify \(I\) with the weak equation.  The same far-field
estimate shows that \(U\in L^{p-1}_{sp}(\R^n)\).  Hence the
weak--viscosity equivalence applies to the continuous weak solution
\(U\); see \cite[Theorem~1.2]{KKL}.  Put
\[
 P_\pm(x_0+z):=U(x_0)+q\cdot z\pm K|z|^2.
\]
After enlarging \(K\), \(P_+\) and \(P_-\) touch \(U\) from above and
below, respectively, in \(B_{r_0}(x_0)\).  For
\(0<\rho<r_0\), let \(U^\pm_\rho=P_\pm\) in \(B_\rho(x_0)\) and
\(U^\pm_\rho=U\) outside that ball.  The pointwise contact formulation
\cite[Proposition~3.1]{KKL} gives, with the sign convention in
\eqref{eq:equation},
\[
 \operatorname{P.V.}\int_{\R^n}
 \frac{\Jp(U^+_\rho(x_0)-U^+_\rho(y))}
 {|x_0-y|^{n+sp}}\dd y\le0,
 \qquad
 \operatorname{P.V.}\int_{\R^n}
 \frac{\Jp(U^-_\rho(x_0)-U^-_\rho(y))}
 {|x_0-y|^{n+sp}}\dd y\ge0.
\]
Each patched function has the same nonzero affine part at \(x_0\).
Applying the local remainder estimate above to \(U\) and to
\(P_\pm\) shows that both displayed principal values converge to \(I\)
as \(\rho\downarrow0\).  Hence \(I\le0\) and \(I\ge0\), so \(I=0\).
\end{proof}

Thus the quadratic contact and the affine tail use the two strict
inequalities $sp<p$ and $sp>p-1$, respectively.

We also record the following consequence of the contact lemma.

\begin{lemma}[Attained translation maximum]
\label{lem:attained-translation}
Let $h\ne0$ and suppose
\[
 M_h:=\sup_{x\in\R^n}\{U(x+h)-U(x)\}
\]
is attained at $x_0$.  Assume that $U$ has quadratic expansions at
$x_0$ and $x_0+h$, with the same nonzero first-order slope.  Then
\[
 U(x+h)-U(x)=M_h\qquad(x\in\R^n).
\]
The analogous conclusion holds for an attained translation minimum.
\end{lemma}

\begin{proof}
Set $V(x)=U(x+h)-M_h$.  Translation invariance and invariance under
addition of constants show that \(V\) is again an entire weak solution.
By the definition of \(M_h\),
\(V\le U\) in \(\R^n\) and \(V(x_0)=U(x_0)\).  The assumed expansions
at \(x_0\) and \(x_0+h\) show that \(V\) and \(U\) have quadratic
expansions at \(x_0\) with the same nonzero affine part.  Thus
Lemma~\ref{lem:pointwise-contact} applies to both functions there.
Moreover,
\[
 V(x_0)-V(y)\ge U(x_0)-U(y).
\]
Subtracting the two pointwise equations gives
\[
 0=\int_{\R^n}
 \frac{\Jp(V(x_0)-V(y))-\Jp(U(x_0)-U(y))}
 {|x_0-y|^{n+sp}}\dd y.
\]
This integral is absolutely convergent: near \(x_0\), both fluxes have
the same odd affine part and their remainders satisfy the estimate in
the proof of Lemma~\ref{lem:pointwise-contact}; at infinity, absolute
integrability follows from \(sp>p-1\) and the global Lipschitz bounds.
The integrand is nonnegative by monotonicity of \(\Jp\).  It therefore
vanishes almost everywhere, and strict monotonicity gives
\[
 V(x_0)-V(y)=U(x_0)-U(y)
 \qquad\text{for almost every }y\in\R^n.
\]
Since \(V(x_0)=U(x_0)\), one has \(V=U\) almost everywhere, and then
everywhere by continuity.
\end{proof}

\subsection{Rigidity of an extremal secant}
\label{sec:extremal-secant}

The following rigidity theorem requires no a priori second-order
regularity.  The extremal segment itself yields the quadratic expansions
needed in Lemma~\ref{lem:attained-translation}.

\begin{theorem}[Extremal-secant rigidity]
\label{thm:extremal-secant-rigidity}
Assume $1<p<2$ and $sp>p-1$.  Let $U$ be a globally Lipschitz entire
weak solution of \eqref{eq:equation}, and write
\[
 L:=[U]_{C^{0,1}(\R^n)}.
\]
If there are distinct points $a,b\in\R^n$ such that
\begin{equation*}
 |U(b)-U(a)|=L|b-a|,
\end{equation*}
then $U$ is affine.
\end{theorem}

\begin{proof}
If $L=0$, the conclusion is immediate.  Replacing $U$ by $-U$ if
necessary, put
\[
 d:=|b-a|,\qquad e:=\frac{b-a}{d},
\]
and assume
\[
 U(b)-U(a)=Ld.
\]
For $0\le t\le d$, the Lipschitz inequalities
\[
 U(a+te)-U(a)\le Lt,
 \qquad
 U(b)-U(a+te)\le L(d-t)
\]
have a sum equal to $Ld$.  Equality therefore holds in both, and
\begin{equation}
 U(a+te)=U(a)+Lt
 \qquad(0\le t\le d).
 \label{eq:linearity-on-extremal-segment}
\end{equation}

Fix $t\in(0,d)$ and let $z$ be small.  The Lipschitz cones with vertices
at $a$ and $b$ give
\begin{align*}
 U(a+te+z)
 &\le U(a)+L|te+z|,\\
 U(a+te+z)
 &\ge U(b)-L|(d-t)e-z|.
\end{align*}
Taylor expansion of the two Euclidean norms yields
\begin{align*}
 |te+z|&=t+e\cdot z+O_t(|z|^2),\\
 |(d-t)e-z|&=(d-t)-e\cdot z+O_{d-t}(|z|^2).
\end{align*}
Together with \eqref{eq:linearity-on-extremal-segment}, these estimates
show that
\begin{equation}
 U(a+te+z)=U(a+te)+Le\cdot z+O_{t,d}(|z|^2).
 \label{eq:automatic-quadratic-expansion}
\end{equation}
Thus every interior point of the extremal segment has a quadratic
expansion with the same nonzero slope $Le$.

Let $0<q<d$.  Choose $t\in(0,d-q)$ and set
$x_0=a+te$, $h=qe$.  By the global Lipschitz bound,
\[
 U(x+h)-U(x)\le Lq\qquad(x\in\R^n),
\]
whereas \eqref{eq:linearity-on-extremal-segment} gives equality at
$x_0$.  Both $x_0$ and $x_0+h$ satisfy
\eqref{eq:automatic-quadratic-expansion}.  Lemma
\ref{lem:attained-translation} therefore gives
\begin{equation}
 U(x+qe)-U(x)=Lq
 \qquad(x\in\R^n)
 \label{eq:all-small-extremal-translations}
\end{equation}
for every $q\in(0,d)$.  Iteration, followed by continuity and reversal
of the translation, extends \eqref{eq:all-small-extremal-translations}
to
\begin{equation}
 U(x+te)=U(x)+Lt
 \qquad(x\in\R^n,\ t\in\R).
 \label{eq:global-extremal-direction}
\end{equation}

It remains to exclude transverse dependence.  Given $z\perp e$, put
$c_z=U(z)-U(0)$.  From \eqref{eq:global-extremal-direction} and the
global Lipschitz bound,
\[
 |c_z+Lt|
 =|U(z+te)-U(0)|
 \le L\sqrt{|z|^2+t^2}
 \qquad(t\in\R).
\]
After squaring and letting $t\to+\infty$ and $t\to-\infty$, we obtain
$c_z=0$.  Given $x\in\R^n$, write $x=z+te$, where $z\perp e$.
Equation \eqref{eq:global-extremal-direction} now gives
\[
 U(x)=U(0)+Le\cdot x,
\]
as asserted.
\end{proof}

\begin{remark}[The singular range]
The proof does not differentiate the equation.  At a generic maximum of
a difference quotient one would need more regularity than is presently
available at the effective order $sp-p+2$.  Here the equality case of
the Lipschitz inequality yields the required quadratic contact through
a geometric two-cone argument, which remains valid in the singular
range.
\end{remark}

\section{Fixed-order large-slope compactness and affine blow-down rigidity}
\label{sec:large-slope-and-blowdown}

We first develop the fixed-order estimates used in the Liouville
argument of Section~\ref{sec:dilation-liouville} and in the iteration
below.  The entire section is independent of the nonlinear Liouville
theorem.  This order of presentation makes every analytic ingredient of the
rigidity argument available before it is invoked.

The argument is organized in three blocks.  Lemmas
\ref{lem:angularly-averaged-shifted-modular}--%
\ref{lem:uniform-shifted-bregman-gluing}
give estimates that are uniform with respect to the affine tilt.
Propositions~\ref{prop:fixed-order-infinite-slope-compactness} and
\ref{prop:large-slope-affine-improvement} identify the stable limit and
transfer its affine decay back to large finite slopes.  Proposition
\ref{prop:noncritical-pointwise-expansion} and Proposition
\ref{prop:affine-blowdown-rigidity} then give the two consequences used
in the nonlinear rigidity argument.

Put
\[
 \Phi(t):=\frac{|t|^p}{p},\qquad
 D_\Phi(a;b):=\Phi(a+b)-\Phi(a)-\Jp(a)b,
 \qquad \lambda_Q:=1+|Q|.
\]
Notice that \(D_\Phi(-a;-b)=D_\Phi(a;b)\).  Thus the corresponding
double-integral density is symmetric under interchange of \(X\) and
\(Y\), even though the Bregman remainder is written with an oriented
increment.
For a tilt \(Q\in\R^n\) and a function \(v\), write
\begin{equation*}
 \mathcal D_Qv(X,Y):=
 \Jp\bigl(Q\cdot(X-Y)+\delta v(X,Y)\bigr)
 -\Jp\bigl(Q\cdot(X-Y)\bigr).
\end{equation*}
The shifted equation is obtained from \eqref{eq:equation} by writing
\(U(X)=Q\cdot X+v(X)\) and subtracting the odd affine flux.  More
precisely, \(v\) solves the equation shifted by \(Q\) in a domain \(D\)
if
\begin{equation*}
 \iint_{\R^n\times\R^n}
 \frac{\mathcal D_Qv(X,Y)(\varphi(X)-\varphi(Y))}
 {|X-Y|^{n+sp}}\dd X\!\dd Y=0
\end{equation*}
for every \(\varphi\in C_c^\infty(D)\).  This is equivalent to saying
that \(Q\cdot X+v\) solves \eqref{eq:equation} in \(D\).  The normalized
interior energy and pointwise exterior flux are
\begin{align}
 \mathcal E_{s,Q}(v;B)
 &:=(1-s)\lambda_Q^{2-p}
 \iint_{B\times B}
 \frac{D_\Phi(Q\cdot(X-Y);\delta v(X,Y))}
 {|X-Y|^{n+sp}}\dd X\!\dd Y,
 \notag\\
 \mathcal T^{\rm pt}_{s,Q}(v)
 &:=\max\left\{
 {\operatorname*{ess\,sup}}_{X\in B_{5/4}}
 \int_{\R^n\setminus B_{4/3}}
 \frac{\lambda_Q^{2-p}|\mathcal D_Qv(X,Y)|}
 {|X-Y|^{n+sp}}\dd Y,
 \right.\notag\\[-2mm]
 &\hspace{31mm}\left.
 {\operatorname*{ess\,sup}}_{X\in B_{3/2}}
 \int_{\R^n\setminus B_2}
 \frac{\lambda_Q^{2-p}|\mathcal D_Qv(X,Y)|}
 {|X-Y|^{n+sp}}\dd Y\right\}.
 \label{eq:large-slope-pointwise-tail}
\end{align}
The first component is used to identify the equation on \(B_{5/4}\);
the second is the separated tail needed by cutoffs supported in
\(B_{3/2}\).  Keeping both components in the definition avoids any
implicit enlargement of the \(X\)-region in the local estimates.
If \(m(X)=c+b\cdot X\) is affine and \(0<\rho\le1\), define
\begin{align}
 \mathfrak Z(v,m;\rho)
 &:={\fint}_{B_{2\rho}}
 \left|\frac{v-m}{\rho}\right|^p\dd X,
 \label{eq:block-zero-modular}\\
 \mathfrak B_{s,Q}(v,m;\rho)
 &:=(1-s)\rho^{sp-p-n}
 \iint_{B_{2\rho}\times B_{2\rho}}
 \frac{D_\Phi\bigl((Q+b)\cdot(X-Y);
              \delta(v-m)(X,Y)\bigr)}
 {|X-Y|^{n+sp}}\dd X\!\dd Y.\notag
\end{align}
We let \(\mathfrak a_{s,Q}(v,m;\rho)\) be the infimum of all \(a>0\)
such that
\begin{equation}
 \mathfrak Z(v,m;\rho)\le a^p,\qquad
 \mathfrak B_{s,Q}(v,m;\rho)
 \le (|Q+b|+a)^{p-2}a^2.
 \label{eq:block-intrinsic-amplitude}
\end{equation}
The two inequalities are kept separate because the energy is quadratic
relative to a large slope, whereas it has \(p\)-growth at zero slope.

We isolate the fixed-gain step used below.  This is the only place where
the two growth regimes of the shifted energy have to be interpolated.
Throughout this section,
\begin{equation*}
 \Psi_p(t):=\frac{t^2}{(1+t)^{2-p}},\qquad t\ge0.
\end{equation*}

For the cutoff and gluing arguments it is convenient to symmetrize the
shifted density.  Put
\[
 \mathcal H_{Q,Z}(t):=\lambda_Q^{2-p}
 \bigl[D_\Phi(Q\cdot Z;t)+D_\Phi(Q\cdot Z;-t)\bigr].
\]

\begin{lemma}[Angularly averaged shifted modular]
\label{lem:angularly-averaged-shifted-modular}
There is \(C=C(n,p)\) such that, for \(|Z|\le4\),
\begin{align}
 c\Psi_p(|t|)&\le \mathcal H_{Q,Z}(t),
 \label{eq:shifted-modular-lower}\\
 \mathcal H_{Q,Z}(t_1+t_2)
 &\le C\bigl[\mathcal H_{Q,Z}(t_1)
              +\mathcal H_{Q,Z}(t_2)\bigr].
 \label{eq:shifted-modular-delta-two}
\end{align}
If \(Q=qe\), \(q\ge0\), \(e\in\mathbb S^{n-1}\), then, for
\(0<r\le4\),
\begin{equation}
 r^{-p}\int_{\mathbb S^{n-1}}
 \mathcal H_{Q,r\omega}(rt)\dd\omega
 \le C\bigl(|t|^p+|t|^2\bigr).
 \label{eq:shifted-modular-angular-upper}
\end{equation}
The same estimates hold if \(D_\Phi(a;t)\) is replaced by the scalar
product
\([\Jp(a+t)-\Jp(a)]t\).
\end{lemma}

\begin{proof}
Taylor's formula, first on the segment from \(a\) to \(a+t\) and then
on the segment from \(a\) to \(a-t\), gives
\begin{equation}
 D_\Phi(a;t)+D_\Phi(a;-t)
 \asymp (|a|+|t|)^{p-2}|t|^2.
 \label{eq:symmetric-shifted-density}
\end{equation}
The constants depend only on \(p\).  The lower estimate follows by
separating \(|t|\le\lambda_Q\) and \(|t|>\lambda_Q\), while the
uniform \(\Delta_2\) estimate follows from
\eqref{eq:symmetric-shifted-density} and the triangle inequality.

It remains to prove the angular upper bound.  The case \(q\le2(1+|t|)\)
follows from the global \(p\)-growth estimate and
\(\lambda_Q^{2-p}\le C(1+|t|)^{2-p}\).  If
\(q>2(1+|t|)\), split the sphere into
\[
 E_t:=\{|e\cdot\omega|\le2|t|/q\}
 \quad\hbox{and}\quad \mathbb S^{n-1}\setminus E_t.
\]
On the complement use the quadratic part of
\eqref{eq:symmetric-shifted-density}; integration of
\(|e\cdot\omega|^{p-2}\) gives \(Cr^p|t|^2\).  On \(E_t\), whose
measure is at most \(C|t|/q\), use the global \(p\)-growth estimate.
The resulting contribution is bounded by
\(Cr^p(|t|^p+|t|^2)\); the quadratic term is needed when \(|t|>1\).
This proves
\eqref{eq:shifted-modular-angular-upper}.  The product version follows
from the same Taylor formula without the factor \(1-t\).
\end{proof}

\begin{lemma}[Fixed-gain Orlicz--Hardy iteration]
\label{lem:fixed-gain-orlicz-hardy}
Let \(g\) be measurable in \(B_2\), let \(F\ge0\), and put, for
\(j\ge1\),
\[
 r_j=\frac43+\frac16\,2^{-j},\qquad
 k_j=A(1-2^{-j}),\qquad d_j=k_{j+1}-k_j=A2^{-j-1},
\]
where \(A\ge1\).  Define
\[
 E_j:=\{g>k_j\}\cap B_{r_j},\qquad
 u_j:=\min\left\{\frac{(g-k_j)_+}{d_j},1\right\},
 \qquad Y_j:=\frac{|E_j|}{|B_2|}.
\]
Suppose that, for cutoffs \(\eta_j\) which equal one
on \(B_{r_{j+1}}\), vanish outside \(B_{r_j}\), and satisfy
\(|\nabla\eta_j|\le C2^j\),
\begin{equation}
 (1-s)\iint_{B_2\times B_2}
 \frac{\Psi_p(|\delta(\eta_ju_j)(X,Y)|)}
 {|X-Y|^{n+sp}}\dd X\!\dd Y
 \le C_0b_0^j\left(1+\frac{F}{A^{p-1}}\right)|E_j|.
 \label{eq:orlicz-hardy-caccioppoli-estimate}
\end{equation}
Then there are \(b>1\) and \(\theta>0\), depending only on
\(n,p,s,C_0,b_0\), such that
\begin{equation}
 Y_{j+1}\le Cb^j
 \left(1+\frac{F}{A^{p-1}}\right)^{1+\theta}
 Y_j^{1+\theta}\qquad(j\ge1).
 \label{eq:orlicz-hardy-fixed-gain}
\end{equation}
Consequently, if \(Y_1\) is sufficiently small, with the smallness
depending only on the displayed coefficient, then \(Y_j\to0\).
The same assertion holds for the lower sets
\(\{-g>k_j\}\) and the corresponding capped functions.
\end{lemma}

\begin{proof}
The function \(\Psi_p\) is an \(N\)-function, satisfies the uniform
\(\Delta_2\) condition, and
\begin{equation}
 \Psi_p(t)\asymp\min\{t^2,t^p\},
 \qquad
 p\le\frac{t\Psi_p'(t)}{\Psi_p(t)}\le2.
 \label{eq:orlicz-two-indices}
\end{equation}
Put \(z_j=\eta_ju_j\).  Since \(0\le z_j\le1\), one has
\(|\delta z_j|\le1\), and therefore
\[
 \Psi_p(|\delta z_j|)\ge c|\delta z_j|^2.
\]
Set \(\sigma=sp/2\).  Since \(sp<p<2\le n\), one has
\(0<\sigma<1\) and \(2\sigma=sp<n\).  Extend \(z_j\) by zero outside
\(B_2\).  Its support is contained in \(B_{r_j}\Subset B_{3/2}\), so
the interactions meeting \(\R^n\setminus B_2\) are separated from the
support and are bounded by \(C\|z_j\|_2^2\).  The fractional Sobolev
inequality, with
\[
 \chi:=\frac{n}{n-sp}>1,
\]
then gives
\begin{align*}
 \left(\int_{B_2}|z_j|^{2\chi}\dd X\right)^{1/\chi}
 &\le C\left\{(1-s)\iint_{B_2^2}
 \frac{|\delta z_j(X,Y)|^2}{|X-Y|^{n+sp}}\dd X\!\dd Y
 +\int_{B_2}|z_j|^2\dd X\right\}\\
 &\le Cb_0^j\left(1+\frac{F}{A^{p-1}}\right)|E_j|.
\end{align*}
Here the harmless factor \((1-s)^{-1}\) is included in
\(C=C(n,p,s)\).  This is precisely the bounded, quadratic branch of
the fractional Orlicz--Sobolev embedding; compare
\cite[Theorems~5.1 and~6.1]{AlbericoCianchiPickSlavikova}.  In the last
line we used \eqref{eq:orlicz-hardy-caccioppoli-estimate} and the fact
that \(z_j\) vanishes outside \(E_j\).
This Sobolev step supplies the fixed Hardy gain in the capped regime.

On \(E_{j+1}\), one has \(u_j=1\) and \(\eta_j=1\), hence \(z_j=1\).
Consequently,
\[
 |E_{j+1}|^{1/\chi}
 \le Cb_0^j\left(1+\frac{F}{A^{p-1}}\right)|E_j|.
\]
After division by \(|B_2|\) and raising to the power \(\chi\), this is
\eqref{eq:orlicz-hardy-fixed-gain}, with
\(\theta=\chi-1>0\) and a larger geometric constant \(b\).
The standard nonlinear sequence lemma now gives \(Y_j\to0\) whenever
\(Y_1\) is sufficiently small.  In particular, the gain is fixed
before \(F\) and \(A\) are chosen.  Replacing \(g\) by \(-g\) gives
the lower-level assertion.
\end{proof}

\begin{lemma}[Uniform shifted level estimate]
\label{lem:uniform-shifted-level-estimate}
Let \(g\) solve the equation shifted by \(P\in\R^n\) in \(B_2\), assume
that \(\mathcal E_{s,P}(g;B_2)<\infty\), fix \(A\ge1\), and put, for
\(j\ge1\),
\begin{equation*}
 r_j=\frac43+\frac16\,2^{-j},\qquad
 k_j=A(1-2^{-j}),\qquad d_j=A2^{-j-1},
\end{equation*}
and
\begin{equation*}
 E_j:=\{g>k_j\}\cap B_{r_j},\qquad
 u_j:=\min\left\{\frac{(g-k_j)_+}{d_j},1\right\}.
\end{equation*}
Let \(\eta_j\in C_c^\infty(B_{r_j})\) satisfy
\(0\le\eta_j\le1\), \(\eta_j=1\) on \(B_{r_{j+1}}\), and
\(|\nabla\eta_j|\le C2^j\).  If the second component in
\eqref{eq:large-slope-pointwise-tail}, with \((Q,v)\) replaced by
\((P,g)\), is at most \(F\), then
\begin{equation}
 (1-s)\iint_{B_2\times B_2}
 \frac{\Psi_p(|\delta(\eta_ju_j)(X,Y)|)}
 {|X-Y|^{n+sp}}\dd X\!\dd Y
 \le Cb_0^j\left(1+\frac{F}{A^{p-1}}\right)|E_j|,
 \label{eq:uniform-shifted-level-estimate}
\end{equation}
where \(C=C(n,p,s)\) and \(b_0>1\) are independent of \(P\).
The same assertion holds for the lower capped levels of \(-g\).
\end{lemma}

\begin{proof}
Fix \(j\), abbreviate \(d=d_j\), and rescale both the function and the
shift:
\[
 \widetilde g:=\frac gd,\qquad \widetilde P:=\frac Pd,\qquad
 u:=u_j=\min\{(\widetilde g-k_j/d)_+,1\}.
\]
By homogeneity, \(\widetilde g\) solves the equation shifted by
\(\widetilde P\).  If \(F_j\) denotes the second tail component of
\((\widetilde P,\widetilde g)\), then
\begin{equation}
 F_j\le F
 \begin{cases}
 d^{1-p},&d\ge1,\\
 d^{-1},&0<d<1,
 \end{cases}
 \le C2^j\frac{F}{A^{p-1}}.
 \label{eq:normalized-capped-tail}
\end{equation}
Indeed,
\(\mathcal D_{\widetilde P}\widetilde g
=d^{1-p}\mathcal D_Pg\).  If \(d\ge1\), use
\(\lambda_{P/d}\le\lambda_P\); if \(d<1\), use
\(\lambda_{P/d}\le d^{-1}\lambda_P\).

Test the rescaled equation with \(\eta_j^2u\), first with the kernel
truncated away from the diagonal.  The scalar monotonicity inequality,
the product version of \eqref{eq:shifted-modular-lower}, and the
\(\varepsilon\)-Young inequality for the shifted
modular give the following two-point bound after the two orientations
of a pair are added:
\begin{align*}
 &\mathcal H_{\widetilde P,X-Y}
       \bigl(\eta_j(X)u(X)-\eta_j(Y)u(Y)\bigr)\\
 &\quad\le C\lambda_{\widetilde P}^{2-p}
 \bigl[\Jp\bigl(\widetilde P\cdot(X-Y)
       +\delta\widetilde g(X,Y)\bigr)
       -\Jp\bigl(\widetilde P\cdot(X-Y)\bigr)\bigr]\\
 &\qquad\qquad\times
 \bigl[\eta_j(X)^2u(X)-\eta_j(Y)^2u(Y)\bigr]\\
 &\qquad+C\mathcal H_{\widetilde P,X-Y}
 \bigl((u(X)+u(Y))\delta\eta_j(X,Y)\bigr).
\end{align*}
Integration of this inequality gives
\begin{align}
 &(1-s)\iint_{B_2^2}
 \frac{\Psi_p(|\delta(\eta_ju)(X,Y)|)}
 {|X-Y|^{n+sp}}\dd X\!\dd Y\notag\\
 &\quad\le C(1-s)\iint_{B_2^2}
 \frac{\mathcal H_{\widetilde P,X-Y}
 \bigl((u(X)+u(Y))\delta\eta_j(X,Y)\bigr)}
 {|X-Y|^{n+sp}}\dd X\!\dd Y
 +CF_j\int_{B_{r_j}}u\dd X.
 \label{eq:unit-cap-caccioppoli}
\end{align}
This is the usual Caccioppoli splitting, but the bounded cap is
essential here.  To check the cutoff term uniformly in
\(\widetilde P\), write \(X-Y=r\omega\) and
\[
 (u(X)+u(Y))\delta\eta_j(X,Y)=r\,t_{X,r}(\omega).
\]
Since \(0\le u\le1\) and \(|\delta\eta_j|\le C2^jr\), one has
\(|t_{X,r}(\omega)|\le C2^j\).  If \(|\widetilde P|<1\), the global
\(p\)-growth estimate gives the required bound.  If
\(\widetilde P=qe\), \(q\ge1\), the pointwise quadratic estimate
\[
 r^{-p}\mathcal H_{\widetilde P,r\omega}
       (r t_{X,r}(\omega))
 \le C|e\cdot\omega|^{p-2}|t_{X,r}(\omega)|^2
\]
and \(\int_{\mathbb S^{n-1}}|e\cdot\omega|^{p-2}\dd\omega<\infty\)
show that the spherical integral is bounded by \(Cb_0^jr^p\).
The cutoff term vanishes unless at least one endpoint belongs to
\(E_j\).  Hence
\[
 (1-s)\int_0^4 r^{p-sp-1}\dd r\le C(p)
\]
shows that the first term on the right of
\eqref{eq:unit-cap-caccioppoli} is at most \(Cb_0^j|E_j|\).
The exterior term is bounded by \(F_j\int u\le F_j|E_j|\).
Now use \eqref{eq:normalized-capped-tail} and absorb \(2^j\) into
\(b_0^j\).  Removing the kernel truncation by monotone convergence
proves \eqref{eq:uniform-shifted-level-estimate}.  The argument for
\(-g\) is identical.
\end{proof}

\begin{lemma}[Shift-uniform local boundedness]
\label{lem:shift-uniform-local-boundedness}
Let \(v\) solve the shifted equation in \(B_2\).  If
\[
 \fint_{B_2}|v|^p\dd X\le1,\qquad
 \mathcal E_{s,Q}(v;B_2)\le1,\qquad
 \mathcal T^{\rm pt}_{s,Q}(v)\le L,
\]
then
\begin{equation}
 \|v\|_{L^\infty(B_{4/3})}
 \le C(n,p,s)\bigl(1+L^{1/(p-1)}\bigr).
 \label{eq:shift-uniform-local-boundedness}
\end{equation}
The constant is independent of \(Q\).
\end{lemma}

\begin{proof}
We first make explicit the coercivity used below.  If
\(|X-Y|\le4\), then
\begin{equation}
 \lambda_Q^{2-p}D_\Phi\bigl(Q\cdot(X-Y);t\bigr)
 \ge c\Psi_p(|t|),
 \label{eq:shift-uniform-orlicz-coercivity}
\end{equation}
with \(c=c(p)>0\), independently of \(Q\).  Indeed, if
\(|t|\le\lambda_Q\), use
\(|Q\cdot(X-Y)|+|t|\le C\lambda_Q\) in
\eqref{eq:large-slope-bregman-equivalence}; if
\(|t|>\lambda_Q\), the same equivalence gives a lower bound by
\(\lambda_Q^{2-p}|t|^p\ge|t|^p\).

Choose \(A=C_0(1+L^{1/(p-1)})\), where \(C_0\ge1\) will be fixed
below.  Apply Lemma~\ref{lem:uniform-shifted-level-estimate} to \(v\),
with this level amplitude \(A\), and set
\begin{equation*}
 r_j=\frac43+\frac16\,2^{-j},\qquad
 k_j=A(1-2^{-j}),\qquad
 E_j^+=\{v>k_j\}\cap B_{r_j},\qquad
 Y_j^+=\frac{|E_j^+|}{|B_2|}.
\end{equation*}
Its conclusion is precisely
\eqref{eq:orlicz-hardy-caccioppoli-estimate}, with \(F=L\).
Lemma~\ref{lem:fixed-gain-orlicz-hardy} consequently yields
\begin{equation*}
 Y_{j+1}^+\le Cb^j
 \left(1+\frac{L}{A^{p-1}}\right)^{1+\theta}
 (Y_j^+)^{1+\theta}\qquad(j\ge1),
\end{equation*}
where \(b>1\) and \(\theta>0\) are structural.  The coefficient in
parentheses is bounded, while Chebyshev's inequality gives
\[
 Y_1^+\le C A^{-p}\fint_{B_2}|v|^p\dd X\le CA^{-p}.
\]
Taking \(C_0\) sufficiently large in the standard iteration lemma gives
\(Y_j^+\to0\), hence \(v\le A\) almost everywhere in \(B_{4/3}\).
Applying the same argument to \(-v\) gives \(v\ge-A\) and proves
\eqref{eq:shift-uniform-local-boundedness}.
\end{proof}

\begin{lemma}[Large-slope scalar and angular bounds]
\label{lem:large-slope-scalar-angular}
There is \(C=C(p)\) such that, for \(a,b\in\R\),
\begin{align}
 D_\Phi(a;b)&\asymp (|a|+|b|)^{p-2}|b|^2,
 \label{eq:large-slope-bregman-equivalence}\\
 |\Jp(a+b)-\Jp(a)|
 &\le C\min\{|b|^{p-1},(|a|+|b|)^{p-2}|b|\}.
 \label{eq:large-slope-flux-bound}
\end{align}
Moreover, if \(Q=q e\), \(q\ge1\), \(e\in\mathbb S^{n-1}\), and
\(Z=r\omega\), then
\begin{equation}
 \lambda_Q^{2-p}
 |\Jp(Q\cdot Z+t)-\Jp(Q\cdot Z)|
 \le Cr^{p-2}|e\cdot\omega|^{p-2}|t|,
 \label{eq:large-slope-angular-flux}
\end{equation}
with the right-hand side interpreted as \(+\infty\) on the equator.
For \(a>1\),
\begin{equation}
 \sup_{e\in\mathbb S^{n-1}}
 \int_{\mathbb S^{n-1}}|e\cdot\omega|^{(p-2)a}\dd\omega<\infty
 \quad\Longleftrightarrow\quad (p-2)a>-1.
 \label{eq:large-slope-angular-integrability}
\end{equation}
\end{lemma}

\begin{proof}
Taylor's formula with integral remainder gives
\[
 D_\Phi(a;b)
 =(p-1)b^2\int_0^1(1-t)|a+tb|^{p-2}\dd t.
\]
Splitting into \(|a|\ge2|b|\) and \(|a|<2|b|\) proves
\eqref{eq:large-slope-bregman-equivalence}.  The global
\((p-1)\)-H\"older continuity of \(\Jp\), combined with the mean-value
formula away from zero, proves \eqref{eq:large-slope-flux-bound}.
The same formula with \(a=Q\cdot Z\), followed by
\(\lambda_Q\asymp q\), gives \eqref{eq:large-slope-angular-flux}.
Finally, in coordinates with \(e=e_1\), the only singularity in the
spherical integral is \(|\omega_1|^{(p-2)a}\) at \(\omega_1=0\);
one-dimensional integration gives \eqref{eq:large-slope-angular-integrability}.
\end{proof}

\begin{lemma}[Relative gluing for shifted Bregman energies]
\label{lem:uniform-shifted-bregman-gluing}
Let \(\Omega\) be bounded with \(\operatorname{diam}\Omega\le4\), and let
\(\eta\in W^{1,\infty}(\Omega)\) satisfy \(0\le\eta\le1\).  Put
\[
 I_\eta:=\{\eta=1\},\qquad O_\eta:=\{\eta=0\},\qquad
 C_\eta:=\{0<\eta<1\},
\]
\[
 \Gamma_\eta:=(C_\eta\times\Omega)\cup(\Omega\times C_\eta),
 \qquad
 \Xi_\eta:=(I_\eta\times O_\eta)\cup(O_\eta\times I_\eta),
\]
and assume that \(C_\eta\Subset\Omega\) and
\(d_\eta:=\operatorname{dist}(I_\eta,O_\eta)>0\).
For \(E\subset\Omega^2\), set
\[
 \mathscr G_{s,Q}(u;E):=(1-s)\lambda_Q^{2-p}
 \iint_E\frac{D_\Phi(Q\cdot(x-y);\delta u(x,y))}
 {|x-y|^{n+sp}}\dd x\!\dd y.
\]
Assume that \(\mathscr G_{s,Q}(v;\Omega^2)\),
\(\mathscr G_{s,Q}(z;I_\eta^2)\), and
\(\mathscr G_{s,Q}(z;\Gamma_\eta)\) are finite and that
\(z-v\in L^p(\Omega)\cap L^m(\Omega)\), where
\[
 a>1,\qquad (p-2)a>-1,\qquad m:=\frac{2a}{a-1}>2.
\]
If \(w:=\eta z+(1-\eta)v\), then, for every
\(\varepsilon\in(0,1)\),
\begin{align}
 &\mathscr G_{s,Q}(w;\Omega^2)-\mathscr G_{s,Q}(v;\Omega^2)
 \notag\\
 &\quad\le \mathscr G_{s,Q}(z;I_\eta^2)
       -\mathscr G_{s,Q}(v;I_\eta^2)
       +\varepsilon\mathscr G_{s,Q}(v;\Xi_\eta)\notag\\
 &\qquad+C\bigl[
 \mathscr G_{s,Q}(z;\Gamma_\eta)
 +\mathscr G_{s,Q}(v;\Gamma_\eta)\bigr]\notag\\
 &\qquad+C_{\varepsilon,\eta,\Omega}
 \bigl(\|z-v\|_{L^p(\Omega)}^p
       +\|z-v\|_{L^m(\Omega)}^2\bigr).
 \label{eq:uniform-shifted-bregman-gluing}
\end{align}
Here \(C=C(p)\).  The constants are independent of \(Q\) and are
locally uniform for \(s\) in compact subsets of \((0,1)\).  The last
constant may depend on \(d_\eta^{-1}\) and
\(\|\nabla\eta\|_\infty\); accordingly, in applications the
approximating-sequence limit is taken before the transition width tends
to zero.
\end{lemma}

\begin{proof}
Set
\[
 \mathcal B_{Q,Z}(t):=\lambda_Q^{2-p}
 D_\Phi(Q\cdot Z;t).
\]
The oriented density must be retained here: interchange of the two
endpoints replaces \((Q\cdot Z,t)\) by \((-Q\cdot Z,-t)\), and
\(D_\Phi(-a;-t)=D_\Phi(a;t)\).  In particular, that interchange does
not replace \(t\) by \(-t\) while keeping the base point fixed.
Convexity, the fact that \(\mathcal B_{Q,Z}(0)=0\), and the uniform
two-index bounds following from
\eqref{eq:large-slope-bregman-equivalence} give
\begin{align}
 \mathcal B_{Q,Z}(t+r)
 &\le(1+\varepsilon)\mathcal B_{Q,Z}(t)
      +C_\varepsilon\mathcal B_{Q,Z}(r),
 \label{eq:relative-gluing-epsilon-delta-two}\\
 \mathcal B_{Q,Z}(\theta t+(1-\theta)r)
 &\le\theta\mathcal B_{Q,Z}(t)
      +(1-\theta)\mathcal B_{Q,Z}(r),
 \qquad 0\le\theta\le1,
 \label{eq:relative-gluing-convex-combination}
\end{align}
with constants independent of \(Q\) and \(Z\).  For the first estimate,
the one-dimensional formula for \(D_\Phi\), together with
\eqref{eq:large-slope-bregman-equivalence}, gives an exponent
\(\kappa_p<\infty\), depending only on \(p\), such that
\[
 0\le t\,\partial_t\mathcal B_{Q,Z}(t)
 \le\kappa_p\mathcal B_{Q,Z}(t)\qquad(t\ne0).
\]
Integration along each ray gives the constant-one estimate
\[
 \mathcal B_{Q,Z}(\lambda t)
 \le\lambda^{\kappa_p}\mathcal B_{Q,Z}(t),
 \qquad \lambda\ge1.
\]
Write \(t+r=(1-\delta)t/(1-\delta)+\delta r/\delta\) and use convexity.
The first term is at most
\((1-\delta)^{1-\kappa_p}\mathcal B_{Q,Z}(t)\); choose
\(\delta=\delta(p,\varepsilon)\) so that this coefficient is at most
\(1+\varepsilon\).
We also use the crude consequence
\[
 \mathcal B_{Q,Z}(t+r)
 \le C_p\bigl[\mathcal B_{Q,Z}(t)+\mathcal B_{Q,Z}(r)\bigr].
\]
Set \(d=z-v\) and
write
\begin{equation}
 \delta w=\eta(x)\delta z+(1-\eta(x))\delta v
 +d(y)\delta\eta.
 \label{eq:gluing-increment-decomposition}
\end{equation}
On \(I_\eta^2\) one has \(w=z\), while on \(O_\eta^2\) one has
\(w=v\).  On pairs meeting \(C_\eta\), the preceding crude inequality,
\eqref{eq:relative-gluing-convex-combination}, and
\eqref{eq:gluing-increment-decomposition} give the two energy terms over
\(\Gamma_\eta\), together with the error containing
\(d(y)\delta\eta\).  On \(I_\eta\times O_\eta\), with \(x\in I_\eta\),
one instead has
\[
 \delta w=\delta v+d(x).
\]
Estimate \eqref{eq:relative-gluing-epsilon-delta-two} leaves
\(\varepsilon\mathscr G_{s,Q}(v;\Xi_\eta)\) and a separated-pair error
containing \(d(x)\).  Notice that the baseline interaction between the
pure inner and outer regions has not been discarded; this is precisely
what permits its cancellation in a local minimality comparison.

It remains to bound the two errors uniformly in the shift.  In polar
coordinates about \(x\), use
\(|\delta\eta|\le\|\nabla\eta\|_\infty r\).  If \(|Q|\le1\), global
\(p\)-growth gives the required \(L^p\)-bound.  If \(Q=qe\), \(q>1\),
then, for almost every \(\omega\),
\[
 r^{-p}\mathcal B_{Q,r\omega}(rt)
 \le C|e\cdot\omega|^{p-2}|t|^2.
\]
H\"older's inequality on the sphere, with exponents \(a\) and
\(a/(a-1)\), applies because \((p-2)a>-1\) and controls the variable
angular amplitude by its \(L^m\)-norm.  Jensen's inequality in \(x\),
followed by translation invariance, gives
\begin{align*}
 &C(1-s)\int_0^{\operatorname{diam}\Omega}r^{p-sp-1}\dd r\\
 &\qquad\times\left(
 \|\nabla\eta\|_\infty^p\|d\|_p^p
 +\|\nabla\eta\|_\infty^2\|d\|_m^2\right)
\end{align*}
for the cutoff error.  For the separated pairs, the same calculation, with
\(|d|/r\) in place of \(\|\nabla\eta\|_\infty|d|\), bounds the
separated-pair error.  The lower bound \(r\ge d_\eta\) is included
in \(C_{\varepsilon,\eta,\Omega}\).  Finally,
\((1-s)\int_0^{\operatorname{diam}\Omega}r^{p-sp-1}\dd r\) is locally
bounded.  This proves \eqref{eq:uniform-shifted-bregman-gluing}.
\end{proof}

Set
\begin{equation*}
 \tau_s:=sp-p+2.
\end{equation*}
Thus \(sp>p-1\) is equivalent to \(1<\tau_s<2\).

\begin{lemma}[Fixed-level clearing and compactness upgrade]
\label{lem:shift-uniform-fixed-level-clearing}
Let \(I\Subset(0,1)\), let
\(D'\Subset D_0\Subset D\subset\R^n\), with \(D\) bounded, and let
\(M\ge1\).  Suppose that \(s_j\in I\), \(Q_j\in\R^n\), and
\(\mathscr G_{s_j,Q_j}(v_j;D^2)<\infty\).  Assume that
\(f_j\in L^\infty(D_0)\),
\begin{equation}
 \begin{split}
 &(1-s_j)\lambda_{Q_j}^{2-p}
 \iint_{D\times D}
 \frac{\mathcal D_{Q_j}v_j(x,y)\delta\varphi(x,y)}
 {|x-y|^{n+s_jp}}\dd x\!\dd y
 =\int_{D_0}f_j\varphi\dd x
 \end{split}
 \label{eq:shifted-regional-source}
\end{equation}
for every \(\varphi\in C_c^\infty(D_0)\), and
\begin{equation}
 \sup_j\left(
 \|v_j\|_{L^\infty(D)}+\|f_j\|_{L^\infty(D_0)}\right)\le M.
 \label{eq:shifted-regional-source-bounds}
\end{equation}
If \(v_j\to v\) strongly in \(L^p_{\rm loc}(D_0)\) and \(v\) is
continuous in \(D_0\), then
\begin{equation}
 \|v_j-v\|_{L^\infty(D')}\longrightarrow0.
 \label{eq:shift-uniform-local-uniform-upgrade}
\end{equation}
Here and in the proof the \(L^\infty\)-norm is the essential-supremum
norm, and the conclusion is uniform with respect to the shifts \(Q_j\).
In the application in Proposition
\ref{prop:fixed-order-infinite-slope-compactness}, each
\(U_j(x):=Q_j\cdot x+v_j(x)\) is a weak solution of the full fractional
\(p\)-Laplace equation.  After discarding finitely many indices,
\(s_jp>p-1\), and
\cite[Theorem~1.1]{BiswasTopp} gives a locally Lipschitz, hence
continuous, representative of \(U_j\), and therefore of \(v_j\).
Thus the convergence in that application is uniform in the ordinary
pointwise sense.
\end{lemma}

\begin{proof}
We first isolate the fixed-level clearing estimate used in the
compactness argument.  Fix
\(B_{2r}(x_0)\Subset D_0\) and an amplitude \(A>0\).  There is
\(\delta_A>0\), depending on \(A,M,I\) and the fixed geometry, but not
on \(Q_j\), such that
\begin{equation}
 \begin{split}
 &\left|\{v_j>k+A/2\}\cap B_{3r/2}(x_0)\right|
 \le\delta_A|B_{2r}|\\
 &\hspace{30mm}\Longrightarrow\qquad
 \operatorname*{ess\,sup}_{B_{4r/3}(x_0)}v_j\le k+A.
 \end{split}
 \label{eq:fixed-level-clearing}
\end{equation}
The corresponding lower-level statement is
\[
 \left|\{v_j<k-A/2\}\cap B_{3r/2}(x_0)\right|
 \le\delta_A|B_{2r}|
 \quad\Longrightarrow\quad
 \operatorname*{ess\,inf}_{B_{4r/3}(x_0)}v_j\ge k-A.
\]

To prove it, set
\[
 g(X):=\frac{v_j(x_0+rX)-k}{A},
 \qquad P:=\frac{rQ_j}{A}.
\]
The normalized source is
\[
 f_{r,A}(X)
 :=r^{s_jp}A^{1-p}
 \left(\frac{\lambda_P}{\lambda_{Q_j}}\right)^{2-p}
 f_j(x_0+rX).
\]
Since
\[
 \frac{\lambda_P}{\lambda_{Q_j}}\le1+\frac rA,
\]
one has
\begin{equation}
 \|f_{r,A}\|_\infty
 \le C M\left(r^{s_jp}A^{1-p}
              +r^{s_jp-p+2}A^{-1}\right).
 \label{eq:fixed-level-scaled-source}
\end{equation}
For fixed \(r,A\), this is bounded uniformly in \(j\) and \(Q_j\).

When the regional equation is tested in \(B_2\), the pairs whose second
endpoint belongs to the rescaled set
\((D-x_0)/r\setminus B_2\) form a separated exterior term.  Since
\(|v_j|\le M\), the angular flux estimate in
Lemma~\ref{lem:large-slope-scalar-angular}, integrated over the
separated radial range, gives
\begin{equation}
 \operatorname*{ess\,sup}_{X\in B_{3/2}}
 \int_{(D-x_0)/r\setminus B_2}
 \frac{\lambda_P^{2-p}|\mathcal D_Pg(X,Y)|}
 {|X-Y|^{n+s_jp}}\dd Y
 \le C_{A,M,I,r,D}.
 \label{eq:fixed-level-regional-tail}
\end{equation}
Indeed, with \(T=2M/A\) and \(R_*=C(D,r)\), the global H\"older
branch for \(|P|<1\) and the angular branch for \(|P|\ge1\) are
bounded together by
\[
 C\left[
 T^{p-1}\int_{1/2}^{R_*}\rho^{-1-s_jp}\dd\rho
 +T\int_{1/2}^{R_*}\rho^{p-3-s_jp}\dd\rho\right].
\]
Both integrals are uniformly finite because their lower endpoint is
separated from zero and \(s_j\in I\Subset(0,1)\).
This is the only point in the level estimate where the shift-specific
angular bound is used; both the source and tail bounds are independent
of \(Q_j\).

Apply the proof of Lemma~\ref{lem:uniform-shifted-level-estimate} to
the levels of \(g\) between \(1/2\) and \(1\).  More precisely, take
the level sequence, capped functions, sets and cutoffs in Lemma
\ref{lem:fixed-gain-orlicz-hardy}, with \(A=1\), and denote them by
\(k_\ell,u_\ell,E_\ell,\eta_\ell\), respectively.  The source term
\eqref{eq:fixed-level-scaled-source} is placed on the right-hand side,
and \eqref{eq:fixed-level-regional-tail} controls the crossed pairs.
Indeed, the test satisfies \(0\le\eta_\ell^2u_\ell\le1\) and is
supported in \(E_\ell\); hence the source and crossed contributions are
bounded, respectively, by
\(\|f_{r,A}\|_\infty|E_\ell|\) and
\(C_{A,M,I,r,D}|E_\ell|\).
Thus
\[
 (1-s_j)\iint_{B_2^2}
 \frac{\Psi_p(|\delta(\eta_\ell u_\ell)(X,Y)|)}
 {|X-Y|^{n+s_jp}}\dd X\!\dd Y
 \le C_Ab^\ell|E_\ell|,
\]
with \(C_A,b\) independent of \(j,Q_j\).  Lemma
\ref{lem:fixed-gain-orlicz-hardy} now gives an initial measure threshold
\(\delta_A>0\) for which the upper levels disappear in \(B_{4/3}\).
This is \eqref{eq:fixed-level-clearing}.  Applying the same argument to
\(-v_j\), with shift \(-Q_j\) and source \(-f_j\), gives the lower
assertion.  Notice that \(A\) is fixed throughout this argument; no
vanishing cap is renormalized.

It remains to use the clearing estimate.  Fix \(x_0\in D_0\) and
\(\varepsilon>0\).  By continuity of \(v\), choose \(r>0\) such that
\(B_{2r}(x_0)\Subset D_0\) and
\[
 \osc_{B_{2r}(x_0)}v\le\frac{\varepsilon}{4}.
\]
In \eqref{eq:fixed-level-clearing} take
\[
 A=\frac{\varepsilon}{2},
 \qquad k=v(x_0)+\frac{\varepsilon}{4}.
\]
If \(v_j>k+A/2=v(x_0)+\varepsilon/2\), then
\(|v_j-v|>\varepsilon/4\) on \(B_{2r}(x_0)\).  Strong \(L^p\)
convergence makes the measure of this set smaller than the fixed
threshold \(\delta_A|B_{2r}|\) for all sufficiently large \(j\).
Therefore
\[
 \operatorname*{ess\,sup}_{B_{4r/3}(x_0)}v_j
 \le v(x_0)+\frac{3\varepsilon}{4}.
\]
The lower clearing estimate, applied with
\(k=v(x_0)-\varepsilon/4\), gives the corresponding lower bound.
Since \(v\) oscillates by at most \(\varepsilon/4\) on this ball,
\(\|v_j-v\|_{L^\infty(B_{4r/3}(x_0))}\le\varepsilon\) for all large
\(j\).  A finite covering of \(D'\) proves
\eqref{eq:shift-uniform-local-uniform-upgrade}.
\end{proof}

\begin{proposition}[Fixed-order infinite-slope compactness]
\label{prop:fixed-order-infinite-slope-compactness}
Let \(s_j\to s\in(0,1)\), \(Q_j=q_je_j\),
\(q_j\to\infty\), and \(e_j\to e\in\mathbb S^{n-1}\).
Suppose that \(v_j\) solves the shifted equation in \(B_2\) and, for
some \(L<\infty\),
\[
 \|v_j\|_{L^\infty(B_{4/3})}
 +\mathcal E_{s_j,Q_j}(v_j;B_2)
 +\mathcal T^{\rm pt}_{s_j,Q_j}(v_j)\le L.
\]
Then, after a subsequence, \(v_j\to v\) locally uniformly in
\(B_{5/4}\) and strongly in \(L^p_{\rm loc}(B_{4/3})\).  Put
\(D=B_{4/3}\).  In \(B_{5/4}\),
the limit solves the regional equation
\begin{equation}
 L^D_{s,e}v=f,\qquad f\in L^\infty(B_{5/4}),
 \label{eq:large-slope-limit-equation}
\end{equation}
whose weak formulation is
\begin{equation}
 (p-1)(1-s)
 \iint_{D\times D}
 \frac{|e\cdot(X-Y)|^{p-2}\delta v(X,Y)\delta\varphi(X,Y)}
 {|X-Y|^{n+sp}}\dd X\!\dd Y
 =\int_Df\varphi\dd X,
 \label{eq:large-slope-limit-form}
\end{equation}
for every \(\varphi\in C_c^\infty(B_{5/4})\).
Its order is \(\tau_s\).  On every \(B\Subset B_{5/4}\), the normalized
energies converge.  More precisely, if \(m=c+b\cdot X\) is fixed, then
\begin{equation}
 \begin{split}
 &\mathcal E_{s_j,Q_j+b}(v_j-m;B)\\
 &\quad\longrightarrow
 \frac{(p-1)(1-s)}2
 \iint_{B\times B}
 \frac{|e\cdot(X-Y)|^{p-2}
       |\delta(v-m)(X,Y)|^2}
 {|X-Y|^{n+sp}}\dd X\!\dd Y.
 \end{split}
 \label{eq:infinite-slope-recentered-energy-convergence}
\end{equation}
\end{proposition}

\begin{proof}
We divide the argument into five steps.

\emph{Step 1: equicoercivity.}
For \(|X-Y|\le4\), \eqref{eq:shift-uniform-orlicz-coercivity} and
\eqref{eq:orlicz-two-indices} give a uniform bound for
\[
 (1-s_j)\iint_{B_2^2}
 \frac{\Psi_p(|\delta v_j(X,Y)|)}
 {|X-Y|^{n+s_jp}}\dd X\!\dd Y.
\]
Fix \(B\Subset B_{4/3}\) and
\(0<t<\frac12\inf_js_jp\), after discarding finitely many indices.
On \(\{|\delta v_j|>1\}\), the preceding modular directly controls the
\(W^{t,p}\) integrand.  On \(\{|\delta v_j|\le1\}\), H\"older's
inequality with exponents \(2/p\) and \(2/(2-p)\) gives
\begin{equation*}
 [v_j]_{W^{t,p}(B)}^p
 \le C\left(1+
 \iint_{B_2^2}
 \frac{\Psi_p(|\delta v_j|)}{|X-Y|^{n+s_jp}}
 \dd X\!\dd Y\right).
\end{equation*}
The integrability of the remaining radial factor is exactly
\(2t<s_jp\).  Rellich compactness and the uniform local bound therefore
give, after a subsequence,
\begin{equation}
 v_j\to v\quad\hbox{strongly in }L^p(B)
 \quad\hbox{and almost everywhere in }B.
 \label{eq:infinite-slope-strong-compactness}
\end{equation}
For \(X\in B_{5/4}\), introduce already at this stage the crossed
interaction
\begin{equation}
 f_j(X):=-2(1-s_j)\int_{\R^n\setminus B_{4/3}}
 \frac{\lambda_{Q_j}^{2-p}\mathcal D_{Q_j}v_j(X,Y)}
 {|X-Y|^{n+s_jp}}\dd Y.
 \label{eq:large-slope-crossed-source}
\end{equation}
Splitting the full weak equation into \(B_{4/3}^2\) and crossed pairs
gives
\begin{equation*}
 (1-s_j)\lambda_{Q_j}^{2-p}
 \iint_{B_{4/3}^2}
 \frac{\mathcal D_{Q_j}v_j(X,Y)\delta\varphi(X,Y)}
 {|X-Y|^{n+s_jp}}\dd X\!\dd Y
 =\int_{B_{5/4}}f_j\varphi\dd X
\end{equation*}
for every \(\varphi\in C_c^\infty(B_{5/4})\).  The separation of
\(B_{5/4}\) and \(\R^n\setminus B_{4/3}\), together with the first
component of \eqref{eq:large-slope-pointwise-tail}, gives
\begin{equation}
 \|f_j\|_{L^\infty(B_{5/4})}\le CL.
 \label{eq:large-slope-crossed-source-bound}
\end{equation}
After discarding finitely many indices, the numbers \(s_j\) belong to a
fixed compact subset of \((0,1)\).  We defer the upgrade from strong
\(L^p_{\rm loc}\) to local uniform convergence until Step~4, where the
continuity of the limiting stable solution is available.

\emph{Step 2: convergence of the regional forms.}
For fixed \(h\ne0\) and \(t\in\R\), Taylor's formula gives, for almost
every \(h\),
\begin{equation}
 \lambda_{Q_j}^{2-p}D_\Phi(Q_j\cdot h;t)
 \longrightarrow\frac{p-1}{2}|e\cdot h|^{p-2}t^2.
 \label{eq:infinite-slope-pointwise-density}
\end{equation}
Choose \(a>1\) such that \((p-2)a>-1\).  On each annulus
\(\varepsilon<|h|<R\), Lemma
\ref{lem:large-slope-scalar-angular} and
\eqref{eq:large-slope-angular-integrability} give an \(L^a\) angular
majorant, uniformly in \(j\).  This bound will be used for the recovery
sequence and for the separated-pair passage below.  The local
\(L^\infty\) bound upgrades
\eqref{eq:infinite-slope-strong-compactness} to strong \(L^2\)
convergence on compact subsets and, with
\(m=2a/(a-1)\), also to strong \(L^m\) convergence.  This is the
integrability dual to the angular \(L^a\)-bound.

For the lower bound, take a subsequence realizing the liminf in
\eqref{eq:infinite-slope-form-liminf}.  After a further subsequence,
\eqref{eq:infinite-slope-strong-compactness} gives
\(v_j\to v\) almost everywhere in \(B\), and hence
\(\delta v_j(X,Y)\to\delta v(X,Y)\) for almost every
\((X,Y)\in B^2\).  For almost every such pair one also has
\(e\cdot(X-Y)\ne0\); Taylor's formula in
\eqref{eq:infinite-slope-pointwise-density} remains valid with
\(t=\delta v_j(X,Y)\).  All Bregman densities are nonnegative, so
Fatou's lemma on \(B^2\cap\{|X-Y|>\varepsilon\}\), followed by monotone
convergence as \(\varepsilon\downarrow0\), yields
\begin{align}
 &\frac{(p-1)(1-s)}2
 \iint_{B^2}
 \frac{|e\cdot(X-Y)|^{p-2}|\delta v|^2}
 {|X-Y|^{n+sp}}\dd X\!\dd Y\notag\\
 &\qquad\le\liminf_{j\to\infty}
 (1-s_j)\lambda_{Q_j}^{2-p}
 \iint_{B^2}
 \frac{D_\Phi(Q_j\cdot(X-Y);\delta v_j)}
 {|X-Y|^{n+s_jp}}\dd X\!\dd Y.
 \label{eq:infinite-slope-form-liminf}
\end{align}
To justify the recovery part, observe that
\[
 \frac{|e\cdot h|^{p-2}}{|h|^{n+sp}}
 =\frac{|e\cdot\omega|^{p-2}}{|h|^{n+\tau_s}}.
\]
The angular density is bounded below by one and has finite spherical
mass.  More precisely, the Fourier symbol of the full quadratic form is
\[
 \mathfrak m_e(\xi)
 =c_{n,\tau_s}\int_{\mathbb S^{n-1}}
 |\xi\cdot\omega|^{\tau_s}|e\cdot\omega|^{p-2}\dd\omega,
\]
and the lower bound and finite mass of the spectral density give
\(c|\xi|^{\tau_s}\le\mathfrak m_e(\xi)\le C|\xi|^{\tau_s}\).
After localization and extension from a ball, the associated energy is
therefore equivalent to the standard \(H^{\tau_s/2}\) energy.  Smooth
functions are dense in this space.  For a smooth function, pointwise
convergence away from the limiting equator and the uniform angular
\(L^a\)-bound in \eqref{eq:large-slope-angular-integrability} give
angular convergence by Vitali's theorem.  Radial dominated convergence
then applies; the radial majorant is integrable precisely because
\(\tau_s<2\).  A
diagonal smooth approximation supplies a recovery sequence.  Thus the
regional functionals have the strong-\(L^2\) \(\Gamma\)-liminf and
recovery properties on every \(B\Subset B_{4/3}\).  These are the only
variational convergence properties used below.

\emph{Step 3: the exterior functional.}
For the crossed interaction \(f_j\) introduced in
\eqref{eq:large-slope-crossed-source}, estimate
\eqref{eq:large-slope-crossed-source-bound} gives weak-star compactness.
After a subsequence,
\(f_j\stackrel{*}{\rightharpoonup}f\) in \(L^\infty(B_{5/4})\).
When convenient, \(f_j\) and \(f\) are extended by zero to \(D\); all
tests below are supported in \(B_{5/4}\).
Consequently, for every \(\varphi\in C_c^\infty(B_{5/4})\),
\begin{equation*}
 \int f_j\varphi\dd X\longrightarrow\int f\varphi\dd X.
\end{equation*}
No convergence of the exterior values of \(v_j\) is needed here; their
entire contribution is encoded in the bounded functions \(f_j\).

\emph{Step 4: identification of the equation and exclusion of an energy
defect.}
Let \(D=B_{4/3}\).  After multiplication of the shifted equation by
\((1-s_j)\lambda_{Q_j}^{2-p}\), its regional part satisfies
\begin{equation}
 \delta\mathscr G_j(v_j)[\varphi]=\int_Df_j\varphi\dd X,
 \qquad \varphi\in C_c^\infty(B_{5/4}),
 \label{eq:regional-euler-large-slope}
\end{equation}
where
\begin{equation*}
 \mathscr G_j(w):=(1-s_j)\lambda_{Q_j}^{2-p}
 \iint_{D\times D}
 \frac{D_\Phi(Q_j\cdot(X-Y);\delta w(X,Y))}
 {|X-Y|^{n+s_jp}}\dd X\!\dd Y.
\end{equation*}
For \(E\subset D^2\), we use
\(\mathscr G_j(u;E):=\mathscr G_{s_j,Q_j}(u;E)\) for the restriction
of this integral to \(E\), and set
\[
 \mathscr G_\infty(u;E):=
 \frac{(p-1)(1-s)}2
 \iint_E\frac{|e\cdot(X-Y)|^{p-2}|\delta u(X,Y)|^2}
 {|X-Y|^{n+sp}}\dd X\!\dd Y.
\]
Thus \(v_j\) locally minimizes
\(\mathscr G_j(w)-\int_Df_jw\) against variations supported in
\(B_{5/4}\).  The sign in the definition of \(f_j\) was chosen so that
\eqref{eq:regional-euler-large-slope} has this form.

For completeness, we spell out the localization which rules out a
diagonal energy defect.  Fix
\(0<r<r_0<r_1<r'<r''<5/4\).  The radii \(r_1,r'\) may be chosen from
the full-measure set on which the localized recovery energies
converge.  We use a bounded interior regularization rather than a
Markov truncation, since the latter need not decrease a Bregman energy
with an oriented base point.  Choose a cutoff which is one on a
neighborhood of \(\overline B_{r'}\) and is compactly supported in
\(B_{r''}\), extend the resulting bounded function by zero, and
mollify it.  This gives smooth functions \(v^{(k)}\) on \(B_{r'}\)
such that
\[
 \|v^{(k)}\|_{L^\infty(B_{r'})}\le\|v\|_{L^\infty(B_{r''})},\qquad
 v^{(k)}\longrightarrow v
 \quad\hbox{in }H^{\tau_s/2}(B_{r'})
        \cap L^q(B_{r'})
\]
for every finite \(q\).  For each fixed \(k\), pointwise convergence
away from the limiting equator, the uniform angular \(L^a\)-bound, and
Vitali's theorem, followed by radial dominated convergence, give
convergence of the prelimit energies of \(v^{(k)}\) on both
\(B_{r_1}^2\) and \(B_{r'}^2\).  Taking first \(j\to\infty\) at fixed
\(k\), and then a diagonal \(k=k(j)\to\infty\), we obtain a bounded
recovery sequence \(z_j:=v^{(k(j))}\) in \(B_{r'}\) such that
\begin{align}
 z_j-v_j&\longrightarrow0
 \quad\hbox{strongly in }L^p(B_{r'})\cap L^m(B_{r'}),\notag\\
 \sup_j\bigl(\|z_j\|_{L^\infty(B_{r'})}
       +\|v_j\|_{L^\infty(B_{r'})}\bigr)&<\infty,
 \label{eq:bounded-recovery-strong-lp-lm}\\
 \sup_j\mathscr G_j(z_j;B_{r'}^2)&<\infty,\notag\\
 \mathscr G_j(z_j;B_{r_1}^2)
 &\longrightarrow\mathscr G_\infty(v;B_{r_1}^2).
 \label{eq:localized-recovery-energy}
\end{align}
Here \(m>2\) is the exponent in
Lemma~\ref{lem:uniform-shifted-bregman-gluing}.  The strong
\(L^m\)-convergence follows from convergence in measure and the uniform
bound.  Outside \(B_{r'}\) set \(z_j=v_j\).  This extension may have a
jump across \(\partial B_{r'}\), and no bound for
\(\mathscr G_j(z_j;D^2)\) is asserted.  Only
\(\mathscr G_j(z_j;I_\eta^2)\) and
\(\mathscr G_j(z_j;\Gamma_\eta)\) enter the gluing argument.  The first
is contained in \(B_{r'}^2\), while the part of the second meeting
\(D\setminus B_{r'}\) is separated from the transition collar and is
estimated below.

For a fixed \(N\ge1\), put
\[
 \rho_i=r_0+\frac{i}{N}(r_1-r_0),\qquad
 C_i=B_{\rho_i}\setminus\overline{B_{\rho_{i-1}}},
 \qquad 1\le i\le N,
\]
and
\[
 \Gamma_i:=(C_i\times D)\cup(D\times C_i).
\]
Since a pair has at most two endpoints in the disjoint annuli,
\[
 \sum_{i=1}^N\mathscr G_j(v_j;\Gamma_i)
 \le2\mathscr G_j(v_j;D^2)\le C.
\]
The analogous sum for \(z_j\) is uniformly bounded as well.  Indeed,
the part with both endpoints in \(B_{r'}\) is bounded by the recovery
energy, while \(C_i\subset B_{r_1}\) is separated from
\(D\setminus B_{r'}\).  On the latter pairs, write
\[
 z_j(x)-v_j(y)=v_j(x)-v_j(y)+z_j(x)-v_j(x)
\]
and apply the uniform \(\Delta_2\) estimate and the separated-pair
bound from Lemma~\ref{lem:uniform-shifted-bregman-gluing}.  Consequently,
\begin{equation*}
 \sum_{i=1}^N\bigl[
 \mathscr G_j(z_j;\Gamma_i)+\mathscr G_j(v_j;\Gamma_i)\bigr]\le C,
\end{equation*}
with \(C\) independent of \(N\) and \(j\).

Start with a subsequence realizing
\(\limsup_j\mathscr G_j(v_j;B_r^2)\).  The pigeonhole principle and a
further subsequence, on which the selected index is constant, give
\begin{equation}
 \mathscr G_j(z_j;\Gamma_i)+\mathscr G_j(v_j;\Gamma_i)
 \le\frac CN
 \label{eq:small-energy-subcollar}
\end{equation}
for some \(i=i(N)\).  Choose a radial Lipschitz cutoff satisfying
\[
 \eta=1\quad\hbox{on }B_{\rho_{i-1}},\qquad
 \eta=0\quad\hbox{on }D\setminus B_{\rho_i},\qquad
 0<\eta<1\quad\hbox{on }C_i,
\]
and set \(w_j=\eta z_j+(1-\eta)v_j\).  Then \(w_j-v_j\) is supported
compactly in \(B_{5/4}\) and is an admissible variation.

Apply Lemma~\ref{lem:uniform-shifted-bregman-gluing} with \(v_j\) as
the reference function.  The energy on \(O_\eta^2\) cancels exactly,
and the complete interaction between \(I_\eta\) and \(O_\eta\) is
retained with coefficient \(\varepsilon\).  Local minimality and
\eqref{eq:small-energy-subcollar} yield, for fixed \(N\) and
\(\varepsilon\),
\begin{align*}
 \mathscr G_j(v_j;B_r^2)
 &\le \mathscr G_j(v_j;I_\eta^2)\\
 &\le \mathscr G_j(z_j;I_\eta^2)
 +\frac CN+\varepsilon\mathscr G_j(v_j;\Xi_\eta)
 +o_{j,N,\varepsilon}(1)\\
 &\le \mathscr G_j(z_j;B_{r_1}^2)
 +\frac CN+C\varepsilon+o_{j,N,\varepsilon}(1).
\end{align*}
The error in the last line of
\eqref{eq:uniform-shifted-bregman-gluing} tends to zero for each fixed
\(N\) by \eqref{eq:bounded-recovery-strong-lp-lm}; its constant is
allowed to depend on the subannulus width.  The linear term is also
included in the remainder, since
\[
 \left|\int_Df_j(w_j-v_j)\dd X\right|
 \le C\|z_j-v_j\|_{L^1(B_{r'})}\longrightarrow0.
\]
Using \eqref{eq:localized-recovery-energy} and taking, in this order,
\[
 j\to\infty,\qquad N\to\infty,\qquad
 \varepsilon\downarrow0,\qquad r_1\downarrow r,
\]
we obtain
\begin{equation}
 \limsup_{j\to\infty}\mathscr G_j(v_j;B_r^2)
 \le\mathscr G_\infty(v;B_r^2).
 \label{eq:localized-relative-gluing-limsup}
\end{equation}
The last passage follows by monotone convergence through good radii.
Together with \eqref{eq:infinite-slope-form-liminf}, this proves
convergence of the regional energies on every ball compactly contained
in \(B_{5/4}\), with no diagonal defect.  A Lipschitz exhaustion gives
the corresponding statement for general compactly contained
subdomains.

We finally pass local minimality, rather than merely the formal Euler
identity, to the limit.  Fix
\[
 K\Subset K_0\Subset K_1\Subset B_{5/4},\qquad
 \psi\in C_c^\infty(K),
\]
and put \(u=v+\psi\).  We describe a simultaneous recovery which keeps
the complete interaction difference.  Apply the bounded interior
regularization to \(v\) in \(K_1\), and use the same diagonal sequence
for \(v\) and \(v+\psi\); at the smooth level these are
\(v^{(k)}\) and \(v^{(k)}+\psi\).  For fixed \(k\), angular Vitali
convergence and radial domination, followed by the diagonal choice,
give recovery of both
limiting energies on every good subdomain of \(K_1\).

Partition the collar \(K_1\setminus K_0\) into \(N\) layers by a
Lipschitz exhaustion.  The sum of the \(v_j\)- and recovery energies
over these layers is uniformly bounded, so one layer has total energy
at most \(C/N\).  Glue the recovery of \(v\) to \(v_j\) across that
layer and call the result \(\widehat v_j\).  On the transition layer,
the two functions differ by a quantity tending to zero in
\(L^p\cap L^m\).  Lemma
\ref{lem:uniform-shifted-bregman-gluing} gives, for fixed
\(N\) and \(\varepsilon>0\),
\[
 \limsup_{j\to\infty}
 \bigl[\mathscr G_j(\widehat v_j)-\mathscr G_j(v_j)\bigr]
 \le\frac CN+C\varepsilon.
\]
Conversely, local minimality with \(\widehat v_j\) gives
\[
 \liminf_{j\to\infty}
 \bigl[\mathscr G_j(\widehat v_j)-\mathscr G_j(v_j)\bigr]\ge0,
\]
because
\(\|f_j\|_\infty\|\widehat v_j-v_j\|_{L^1(D)}\to0\).
Taking \(j\to\infty\), \(N\to\infty\), and
\(\varepsilon\downarrow0\), in that order, and then taking a diagonal
sequence, gives
\begin{equation}
 \widehat v_j-v_j\longrightarrow0
 \quad\hbox{in }L^p(D)\cap L^m(D),\qquad
 \mathscr G_j(\widehat v_j)-\mathscr G_j(v_j)\longrightarrow0.
 \label{eq:simultaneous-local-recovery-base}
\end{equation}

Now set \(w_j=\widehat v_j+\psi\).  Then \(w_j-v_j\) is compactly
supported in \(K_1\).  The uniform \(L^\infty(D)\) bound and local strong
\(L^p\)-convergence imply \(v_j\to v\) strongly in \(L^q(D)\) for every
finite \(q\).  Indeed, one first passes to the limit on \(B_r\),
\(r<4/3\), and then lets \(r\uparrow4/3\); the contribution of
\(D\setminus B_r\) to the \(L^q\)-norm is bounded by
\(C|D\setminus B_r|^{1/q}\).  Together with
\eqref{eq:simultaneous-local-recovery-base}, this gives
\(\widehat v_j\to v\) strongly in \(L^m(D)\).

We now keep track of the oriented density occurring in the energy.  For
\(Z\in\R^n\), set
\begin{equation*}
 \mathcal B_{j,Z}(t):=\lambda_{Q_j}^{2-p}
 D_\Phi(Q_j\cdot Z;t),\qquad
 \Delta_{j,Z}(t,r):=\mathcal B_{j,Z}(t+r)-\mathcal B_{j,Z}(t).
\end{equation*}
Apply \eqref{eq:relative-gluing-epsilon-delta-two} first to \((t,r)\)
and then to \((t+r,-r)\), and use the crude \(\Delta_2\) estimate from
the proof of Lemma~\ref{lem:uniform-shifted-bregman-gluing}.  After
decreasing the small parameter by a structural factor, one obtains, for
every \(\theta\in(0,1)\),
\begin{equation}
 |\Delta_{j,Z}(t,r)|
 \le\theta\mathcal B_{j,Z}(t)
 +C_\theta\bigl[\mathcal B_{j,Z}(r)+\mathcal B_{j,Z}(-r)\bigr]
 =\theta\mathcal B_{j,Z}(t)+C_\theta\mathcal H_{Q_j,Z}(r).
 \label{eq:oriented-bregman-perturbation-bound}
\end{equation}
No symmetrization of the base increment is made in this estimate.  With
\(t=\delta\widehat v_j(X,Y)\) and
\(r=\delta\psi(X,Y)\), the first term is bounded by
\(\theta\mathscr G_j(\widehat v_j;D^2)\), uniformly in \(j\), while
the pointwise estimates behind
\eqref{eq:shifted-modular-angular-upper} give the required bound even
though the angular amplitude is not fixed.  Indeed, write
\(X-Y=r\omega\) and
\[
 \delta\psi(X,Y)=r\,t_{X,r}(\omega),
 \qquad |t_{X,r}(\omega)|\le\|\nabla\psi\|_\infty.
\]
If \(|Q_j|<1\), the global \(p\)-growth estimate gives
\(\mathcal H_{Q_j,r\omega}(rt_{X,r})\le C_\psi r^p\).  If
\(Q_j=q_je_j\), \(q_j\ge1\), then
\[
 r^{-p}\mathcal H_{Q_j,r\omega}(rt_{X,r}(\omega))
 \le C|e_j\cdot\omega|^{p-2}|t_{X,r}(\omega)|^2.
\]
Since \(p-2>-1\), the last angular weight is integrable, uniformly in
the direction \(e_j\).  Thus the spherical integral is at most
\(C_\psi r^p\), pointwise in \((X,r)\), and radial integration gives
\[
 (1-s_j)\iint_{\substack{D^2\\|X-Y|<\delta}}
 \frac{\mathcal H_{Q_j,X-Y}(\delta\psi(X,Y))}
 {|X-Y|^{n+s_jp}}\dd X\!\dd Y
 \le C\delta^{p-s_jp}.
\]
Since \(s_j\to s<1\), the right-hand side tends to zero uniformly for
large \(j\).  Thus the prelimit near-diagonal perturbation is made small
by first choosing \(\theta\), and then \(\delta\).

The limiting quadratic form has the same absolute continuity at the
diagonal.  Indeed, Cauchy--Schwarz gives
\begin{align*}
 &\iint_{\substack{D^2\\|X-Y|<\delta}}
 \frac{|e\cdot(X-Y)|^{p-2}
 \bigl|2\,\delta v\,\delta\psi+|\delta\psi|^2\bigr|}
 {|X-Y|^{n+sp}}\dd X\!\dd Y\\
 &\quad\le C
 \mathscr G_\infty(v;\{|X-Y|<\delta\})^{1/2}
 \mathscr G_\infty(\psi;\{|X-Y|<\delta\})^{1/2}
 +C\mathscr G_\infty(\psi;\{|X-Y|<\delta\}).
\end{align*}
The right-hand side tends to zero: the first energy is absolutely
continuous and the smooth-function energy is
\(O(\delta^{p-sp})\).

On \(|X-Y|\ge\delta\), the angular \(L^a\)-bound, the strong
\(L^m(D)\)-convergence, and Vitali's theorem give convergence of the
oriented energy differences on every separated compact region.  The
crossed region between \(K_0\) and \(D\setminus K_1\) has positive
separation and is covered by the same estimate.  Letting first
\(j\to\infty\), then
\(\delta\downarrow0\), yields
\[
 \mathscr G_j(\widehat v_j+\psi)-\mathscr G_j(\widehat v_j)
 \longrightarrow
 \mathscr G_\infty^D(v+\psi)-\mathscr G_\infty^D(v).
\]
Together with \eqref{eq:simultaneous-local-recovery-base}, this yields
\begin{align}
 w_j-v_j&\longrightarrow\psi\quad\hbox{in }L^1(D),\notag\\
 \limsup_{j\to\infty}
 \bigl[\mathscr G_j(w_j)-\mathscr G_j(v_j)\bigr]
 &\le
 \mathscr G_\infty^D(v+\psi)-\mathscr G_\infty^D(v).
 \label{eq:localized-minimality-recovery}
\end{align}
Here the right-hand side denotes the finite energy difference over
pairs having at least one endpoint in \(K_1\); outside that set the two
functions agree.

Local minimality of \(v_j\), weak-star convergence of \(f_j\), and
\eqref{eq:localized-minimality-recovery} give
\[
 \int_Df_j(w_j-v_j)\dd X
 =\int_Df_j\bigl[(w_j-v_j)-\psi\bigr]\dd X
  +\int_Df_j\psi\dd X
 \longrightarrow\int_Df\psi\dd X.
\]
Consequently,
\[
 \mathscr G_\infty^D(v)-\int_Dfv\dd X
 \le
 \mathscr G_\infty^D(v+\psi)-\int_Df(v+\psi)\dd X.
\]
Apply this inequality to \(t\psi\), \(t\in\R\), divide by \(t\), and
let \(t\to0^\pm\).  Since the limiting functional is quadratic, the
two one-sided inequalities coincide and give
\eqref{eq:large-slope-limit-form}, hence also
\eqref{eq:large-slope-limit-equation}.

The limiting regional kernel has order \(\tau_s>1\), finite spectral
mass, and a uniform spectral lower bound.  After a bounded extension
outside \(D\), its regional equation becomes a full stable equation
with a bounded right-hand side on compact subsets of \(B_{5/4}\).
The interior stable estimate \cite{RosOtonSerraStable} therefore shows
that \(v\) is continuous there.  Apply Lemma
\ref{lem:shift-uniform-fixed-level-clearing}, for every
\(B\Subset B_{5/4}\), with
\(D'=B\), \(D=B_{4/3}\), \(D_0=B_{5/4}\), the source bound
\eqref{eq:large-slope-crossed-source-bound}, and
\eqref{eq:infinite-slope-strong-compactness}.  We conclude that
\begin{equation}
 v_j\longrightarrow v\qquad\hbox{locally uniformly in }B_{5/4}.
 \label{eq:infinite-slope-local-uniform-convergence}
\end{equation}

\emph{Step 5: affine recentering.}
Let \(m=c+b\cdot X\) be fixed.  Then \(v_j-m\) solves the equation with
tilt \(Q_j+b\).  Since
\[
 \frac{|Q_j+b|}{|Q_j|}\longrightarrow1,
 \qquad \frac{Q_j+b}{|Q_j+b|}\longrightarrow e,
\]
the normalized factors for \(Q_j+b\) and \(Q_j\) are comparable and
their ratio tends to one.  We also verify the two bounds used in
Steps~1--4.  The Bregman fundamental estimate above, with the fixed
affine increment \(b\cdot(X-Y)\), preserves the local energy bound.  For
the exterior flux, add and subtract \(\Jp(Q_j\cdot(X-Y))\).  The original
term is controlled by \(\mathcal T^{\rm pt}_{s_j,Q_j}(v_j)\), while
Lemma~\ref{lem:large-slope-scalar-angular} bounds the affine correction
by
\[
 C|b|\,|X-Y|^{p-1}|e_j\cdot\omega|^{p-2},
 \qquad \omega=\frac{X-Y}{|X-Y|}.
\]
Its angular integral is finite and its radial integral at infinity is
\(\int_1^\infty r^{p-2-s_jp}\dd r\), uniformly finite for all large
\(j\) because \(sp>p-1\).  Thus Steps~1--4 apply with unchanged
structural constants.  This proves the stated energy convergence after
subtracting \(m\) and completes the proof.
\end{proof}

\begin{proposition}[Fixed-order large-slope affine improvement]
\label{prop:large-slope-affine-improvement}
Fix \(1<p<2\), \(s\in(0,1)\) with \(sp>p-1\), and
\[
 0<\alpha<sp-p+1,\qquad L\ge1.
\]
There are \(\rho_\infty\in(0,1/8)\), \(Q_0<\infty\), and \(C<\infty\)
such that the following holds.  If \(v\) solves the shifted equation in
\(B_2\), \(|Q|\ge Q_0\), and
\[
 \fint_{B_2}|v|^p\dd X\le1,\qquad
 \mathcal E_{s,Q}(v;B_2)\le1,\qquad
 \mathcal T^{\rm pt}_{s,Q}(v)\le L,
\]
then there is an affine map \(m\), with \(|m(0)|+|\nabla m|\le C\),
such that
\begin{equation}
 \mathfrak a_{s,Q}(v,m;\rho_\infty)\le\rho_\infty^\alpha.
 \label{eq:large-slope-affine-improvement}
\end{equation}
More precisely, one may fix
\begin{equation}
 \alpha<\bar\alpha<sp-p+1
 \label{eq:large-slope-intermediate-exponent}
\end{equation}
and choose the same affine map so that
\begin{equation}
 \mathfrak a_{s,Q}(v,m;r)
 +r^{-1}\|v-m\|_{L^\infty(B_{2r})}
 \le C_{\rm L}r^{\bar\alpha}
 \quad(\rho_\infty\le r\le1/2),
 \qquad
 C_{\rm L}\rho_\infty^{\bar\alpha-\alpha}\le1.
 \label{eq:large-slope-intermediate-block}
\end{equation}
Here the constants may depend on \(L\), but not on \(Q\).
\end{proposition}

\begin{proof}
Lemma~\ref{lem:shift-uniform-local-boundedness} first gives a uniform
\(L^\infty(B_{4/3})\) bound.  If the conclusion failed, Proposition
\ref{prop:fixed-order-infinite-slope-compactness} would give a limit
solving \eqref{eq:large-slope-limit-equation}.  In polar coordinates the
kernel is
\[
 \frac{|e\cdot\omega|^{p-2}}{r^{n+\tau_s}}.
\]
Its spectral mass is bounded above by
\eqref{eq:large-slope-angular-integrability}; since \(p-2<0\) and
\(|e\cdot\omega|\le1\), it also satisfies the uniform spectral lower
ellipticity condition.  Extend \(v\) boundedly from \(B_{4/3}\) to
\(\R^n\).  On \(B_{3/4}\), the difference between the regional operator
in \eqref{eq:large-slope-limit-form} and the full stable operator is a
bounded function, because the two sets are separated.  Thus the usual
interior estimate for the full stable operator applies after this
localization.  Choose \(\bar\alpha\) as in
\eqref{eq:large-slope-intermediate-exponent}.  The interior estimate for
symmetric stable operators \cite{RosOtonSerraStable} gives
\[
 \|v\|_{C^{1,\bar\alpha}(B_{3/4})}\le C.
\]
Let \(m\) be the first-order Taylor polynomial of \(v\) at the origin.
The zero-order remainder is \(O(r^{1+\bar\alpha})\), while direct
integration in \eqref{eq:large-slope-limit-form} gives the normalized
quadratic energy remainder \(O(r^{2\bar\alpha})\).  Consequently
\begin{equation}
 \mathfrak a_{s,e}^{\rm lin}(v,m;r)
 \le C_*r^{\bar\alpha}
 \qquad(0<r\le1/2),
 \label{eq:large-slope-limit-full-block}
\end{equation}
Here \(\mathfrak a_{s,e}^{\rm lin}\) is defined by the zero-order
condition in \eqref{eq:block-zero-modular} and by
\[
 (1-s)r^{sp-p-n}
 \iint_{B_{2r}^2}
 \frac{|e\cdot(X-Y)|^{p-2}
       |\delta(v-m)(X,Y)|^2}
 {|X-Y|^{n+sp}}\dd X\!\dd Y\le a^2.
\]

We record why the estimate transfers uniformly on a compact interval of
radii.  Otherwise there would be a contradicting sequence
\(Q_j\to\infty\), affine maps \(m_j\) converging to \(m\), and radii
\(r_j\in[\rho,1/2]\) for which the estimate with constant \(2C_*\)
fails.  After a subsequence, \(r_j\to r\in[\rho,1/2]\).  Strong local
\(L^p\) convergence transfers the zero-order modular, and the recentered
no-defect statement in Proposition
\ref{prop:fixed-order-infinite-slope-compactness} transfers the energy
modular.  This contradicts \eqref{eq:large-slope-limit-full-block}.
Hence, for every fixed \(\rho>0\), the estimate is uniform on
\([\rho,1/2]\) once \(|Q|\) is sufficiently large.

Take \(C_{\rm L}=2C_*\), decrease \(\rho_\infty\) until
\(C_{\rm L}\rho_\infty^{\bar\alpha-\alpha}\le1\), and then choose
\(Q_0\) by the preceding compactness argument.  This proves
\eqref{eq:large-slope-intermediate-block}; its endpoint is
\eqref{eq:large-slope-affine-improvement}.
\end{proof}

\begin{corollary}[Large-slope improvement at a prescribed scale]
\label{cor:prescribed-large-slope-block-main}
Fix
\[
 0<\alpha<\bar\alpha<sp-p+1,
 \qquad L\ge1.
\]
There is \(C_{\rm L}=C_{\rm L}(n,p,s,\alpha,\bar\alpha,L)\)
such that, for every prescribed \(0<\rho<1/8\) satisfying
\begin{equation}
 C_{\rm L}\rho^{\bar\alpha-\alpha}\le1,
 \label{eq:prescribed-large-slope-radius-main}
\end{equation}
there is \(Q_{\rm th}<\infty\) for which
\begin{equation}
 \mathfrak a_{s,Q}(v,m;t)
 +t^{-1}\|v-m\|_{L^\infty(B_{2t})}
 \le C_{\rm L}t^{\bar\alpha}
 \qquad(\rho\le t\le1/2)
 \label{eq:prescribed-large-slope-full-block-main}
\end{equation}
holds whenever \(v\) solves the shifted equation in \(B_2\),
\(|Q|>Q_{\rm th}\), and
\[
 \fint_{B_2}|v|^p\dd X\le1,\qquad
 \mathcal E_{s,Q}(v;B_2)\le1,\qquad
 \mathcal T^{\rm pt}_{s,Q}(v)\le L.
\]
In particular,
\[
 \mathfrak a_{s,Q}(v,m;\rho)\le\rho^\alpha,
 \qquad |m(0)|+|\nabla m|\le C,
\]
where \(C\) is independent of \(Q_{\rm th},Q,\rho\).
\end{corollary}

\begin{proof}
The stable estimate and therefore \(C_{\rm L}\) are obtained before
the radius is chosen.  For a fixed radius satisfying
\eqref{eq:prescribed-large-slope-radius-main}, the compactness transfer
on \(t\in[\rho,1/2]\), proved in Proposition
\ref{prop:large-slope-affine-improvement}, supplies \(Q_{\rm th}\).
The endpoint follows from
\eqref{eq:prescribed-large-slope-radius-main}.  The unit \(L^p(B_2)\)
bound for \(v\), the estimate at \(t=1/2\), and norm equivalence for
affine maps on \(B_1\) give the coefficient bound.
\end{proof}

\begin{proposition}[Noncritical pointwise expansion]
\label{prop:noncritical-pointwise-expansion}
Let \(U\) be a globally Lipschitz entire solution and let \(x_0\) be a
differentiability point with \(q=\nabla U(x_0)\ne0\).  For every
\(0<\alpha<sp-p+1\), there are \(r_0>0\) and \(C_{x_0}<\infty\) such that
\begin{equation}
 \left(\fint_{B_r(x_0)}
 |U-U(x_0)-q\cdot(\,\cdot-x_0)|^p\dd x\right)^{1/p}
 \le C_{x_0}r^{1+\alpha}
 \qquad(0<r<r_0).
 \label{eq:large-slope-noncritical-expansion}
\end{equation}
\end{proposition}

\begin{proof}
Write
\[
 \varepsilon(r):=\frac1r
 \|U-U(x_0)-q\cdot(\,\cdot-x_0)\|_{L^\infty(B_r(x_0))}.
\]
Differentiability gives \(\varepsilon(r)\to0\).  Replace it by the
monotone modulus
\[
 \varepsilon_*(r):=\sup_{0<t\le r}\varepsilon(t),
\]
which also tends to zero.  Fix first a sufficiently large structural
number \(C_*\).  Choose \(r_*>0\) so small that
\[
 H_*:=C_*\bigl(\varepsilon_*(r_*)+r_*^{sp-p+1}\bigr)
 \quad\hbox{satisfies}\quad
 \frac{|q|}{H_*}\ge4Q_0
\]
and \(CH_*\le |q|/4\), where \(Q_0\) and \(C\) are the constants in
Proposition~\ref{prop:large-slope-affine-improvement}.  Next choose
\(0<r_0<r_*/8\) so small that the normalized zero-order and Caccioppoli
energies at scale \(r_0\), with amplitude \(H_*\), are at most one and
\[
 \frac{r_0^{sp-p+1}}{H_*}
 r_*^{p-1-sp}\le1.
 \label{eq:noncritical-initial-far-tail}
\]

We verify the tail initialization.  On the physical annuli between
\(r_0\) and \(r_*\), differentiability gives
\[
 |U(x)-U(x_0)-q\cdot(x-x_0)|\le H_*|x-x_0|.
\]
After division by \(r_0H_*\), these annuli have a common linear
majorant.  Lemma~\ref{lem:large-slope-scalar-angular} therefore gives a
uniform relative-flux contribution.  Outside \(B_{r_*}(x_0)\), the
global Lipschitz estimate, affine cancellation, and radial integration
give
\[
 C\frac{r_0^{sp-p+1}}{H_*}
 \int_{r_*}^{\infty}t^{p-2-sp}\dd t
 \le C
\]
by \eqref{eq:noncritical-initial-far-tail}.  Thus the pointwise tail,
the local \(L^\infty\) norm, and the normalized energy all satisfy the
hypotheses of Proposition
\ref{prop:large-slope-affine-improvement}.

Iterate that proposition with
\[
 r_k=r_0\rho_\infty^k,\qquad
 H_k=H_*\rho_\infty^{k\alpha}.
\]
The normalized affine increments have bounded slopes.  Their physical
sum is at most \(CH_*\le |q|/4\); hence every updated physical slope has
length at least \(|q|/2\), and the normalized tilt remains above \(Q_0\)
at all levels.  We record the tail induction.  If \(i<k\), summation of
the affine increments from level \(i\) to level \(k\) and the estimate at
level \(i\) give
\begin{equation}
 \|U-\ell_k\|_{L^\infty(B_{2r_i}(x_0))}
 \le Cr_iH_i.
 \label{eq:noncritical-old-annulus-control}
\end{equation}
On the physical annulus \(r_{i+1}<|y-x_0|\le r_i\), every updated
physical slope is bounded below by \(|q|/2\).  We may therefore use the
angular linearized estimate \eqref{eq:large-slope-angular-flux} at every
stage.  After normalization by \(H_k\), the contribution of this
annulus is bounded by
\[
 C\left(\frac{r_k}{r_i}\right)^\gamma\frac{H_i}{H_k}
 =C\rho_\infty^{(k-i)(\gamma-\alpha)},
 \qquad \gamma=sp-p+1.
\]
The angular singularity is integrable because \(p-2>-1\).  Since
\(\alpha<\gamma\), summation in \(i\) gives a uniform geometric bound,
independently of \(k\).  Notice that the global
\((p-1)\)-H\"older estimate is not used here: after multiplication by
the intrinsic factor \(\lambda_{Q_k}^{2-p}\), it would produce
\((r_k/r_i)^\gamma H_i^{p-1}/H_k\), which is not the state propagated
by the large-slope iteration.  The region exterior to the initial ball
carries the decaying factor
\[
 \frac{r_k^{sp-p+1}}{H_k}
 =\frac{r_0^{sp-p+1}}{H_*}
   \rho_\infty^{k(sp-p+1-\alpha)},
\]
and therefore also decreases with \(k\).  Thus the relative tail and the normalized Caccioppoli
energy are reproduced at every step.  Consequently
\[
 \left(\fint_{B_{r_k}(x_0)}
 |U-\ell_k|^p\dd x\right)^{1/p}
 \le Cr_kH_k.
\]
Summing the affine increments and using the differentiability of \(U\)
at \(x_0\) identifies the limiting affine map with
\(U(x_0)+q\cdot(x-x_0)\).  Interpolation between consecutive radii gives
\eqref{eq:large-slope-noncritical-expansion}, with
\(C_{x_0}=CH_*r_0^{-\alpha}\).
\end{proof}

\begin{proposition}[One affine blow-down is rigid]
\label{prop:affine-blowdown-rigidity}
Let \(U\) be a globally Lipschitz entire weak solution, normalized by
\(U(0)=0\), and set \(U_R(X)=U(RX)/R\).  If, for some \(R_j\to\infty\)
and \(b\in\R^n\),
\begin{equation}
 U_{R_j}\longrightarrow b\cdot X
 \quad\text{locally uniformly in }\R^n,
 \label{eq:affine-blowdown-convergence}
\end{equation}
then \(U(X)=b\cdot X\) in \(\R^n\).
\end{proposition}

\begin{proof}
Assume first that \(b=0\).  If
\(\Lambda=\Lip(U;\R^n)\), then \(|U_{R_j}(Y)|\le\Lambda|Y|\).  For
\(M>2\), local uniform convergence and this linear majorant give
\begin{align*}
 \int_{\R^n\setminus B_2}
 \frac{|U_{R_j}(Y)|^{p-1}}{|Y|^{n+sp}}\dd Y
 &\le o_j(1)
 +C\Lambda^{p-1}\int_M^\infty r^{p-2-sp}\dd r\\
 &\le o_j(1)+C\Lambda^{p-1}M^{p-1-sp}.
\end{align*}
Letting first \(j\to\infty\) and then \(M\to\infty\) shows convergence
to zero in the full tail norm.  Since $sp>p-1$, the local Lipschitz
estimate \cite[Theorem~2.1]{BiswasTopp}, applied to \(U_{R_j}\) in
\(B_2\), yields
\[
 [U_{R_j}]_{C^{0,1}(B_1)}\longrightarrow0.
\]
But the left-hand side equals
\([U]_{C^{0,1}(B_{R_j})}\).  Hence \(U\equiv0\).

Suppose now that \(b\ne0\), and put \(V_j=U_{R_j}-b\cdot X\).  Define
\begin{align*}
 d_j&:=\|V_j\|_{L^\infty(B_4)},\\
 E_j&:=(1-s)\iint_{B_2\times B_2}
 \frac{D_\Phi(b\cdot(X-Y);\delta V_j(X,Y))}
 {|X-Y|^{n+sp}}\dd X\!\dd Y,\\
 T_j&:=\max\left\{
 {\operatorname*{ess\,sup}}_{X\in B_{5/4}}
 \int_{\R^n\setminus B_{4/3}}
 \frac{|\mathcal D_bV_j(X,Y)|}{|X-Y|^{n+sp}}\dd Y,\right.\\
 &\hspace{26mm}\left.
 {\operatorname*{ess\,sup}}_{X\in B_{3/2}}
 \int_{\R^n\setminus B_2}
 \frac{|\mathcal D_bV_j(X,Y)|}{|X-Y|^{n+sp}}\dd Y\right\}.
\end{align*}
Then
\begin{equation*}
 d_j+E_j+T_j\longrightarrow0.
\end{equation*}
Indeed, \(d_j\to0\) by hypothesis.  The common Lipschitz bound, the
estimate \(D_\Phi(a;t)\le C|t|^p\), and a radial split at \(d_j\) give
\(E_j\le Cd_j^{p-sp}\).  For both components of \(T_j\), local uniform
convergence handles every fixed exterior ball, whereas
\(|\mathcal D_bV_j(X,Y)|\le C|X-Y|^{p-1}\) gives the uniformly integrable
far-field majorant \(Cr^{p-2-sp}\).

Choose a large structural constant \(C_0\) and put
\[
 a_j=C_0\bigl(d_j+E_j^{1/2}+T_j+R_j^{-1}\bigr),\qquad
 v_j=\frac{V_j}{a_j},\qquad Q_j=\frac b{a_j}.
\]
Then \(a_j\to0\), \(|Q_j|\to\infty\), and the homogeneity identities
\[
 D_\Phi(A/a_j;B/a_j)=a_j^{-p}D_\Phi(A;B),\qquad
 \Jp(A/a_j)=a_j^{1-p}\Jp(A)
\]
show, after increasing \(C_0\), that \(v_j,Q_j\) satisfy all hypotheses
of Proposition~\ref{prop:large-slope-affine-improvement} with constants
independent of \(j\).

We record the far-field calculation needed to iterate that proposition
at every level.  Put \(r_k=\rho_\infty^k\),
\(H_k=\rho_\infty^{k\alpha}\), and let
\(m_{j,k}=c_{j,k}+q_{j,k}\cdot X\) be the accumulated affine
approximation.  The coefficient estimate gives \(|q_{j,k}|\le C\), and
the new tilt is
\[
 P_{j,k}=\frac{Q_j+q_{j,k}}{H_k}.
\]
Thus the large-slope regime persists.  In the part of the pointwise tail
coming from beyond the initial ball, the changes of variables
\(Z=r_kX\), \(W=r_kY\) give the exact factor
\begin{align}
 &\int_{\{|Y|>4/(3r_k)\}}
 \frac{\lambda_{P_{j,k}}^{2-p}
 |\mathcal D_{P_{j,k}}w_{j,k}(X,Y)|}
 {|X-Y|^{n+sp}}\dd Y\notag\\
 &\quad\le
 C\frac{r_k^{sp-p+1}}{a_jH_k}
 \int_{|W|>4/3}
 \frac{|\mathcal D_bV_j(Z,W)|
 +|\Jp((b+a_jq_{j,k})\cdot(Z-W))
       -\Jp(b\cdot(Z-W))|}
 {|Z-W|^{n+sp}}\dd W.
 \label{eq:affine-blowdown-far-field}
\end{align}
The first numerator is controlled by \(T_j\).  For the affine correction,
angular integration in
\eqref{eq:large-slope-angular-integrability} gives
\[
 \int_{\mathbb S^{n-1}}
 |\Jp((b+a_jq_{j,k})\cdot r\omega)
 -\Jp(b\cdot r\omega)|\dd\omega
 \le Ca_j|q_{j,k}|r^{p-1}.
\]
Consequently the right-hand side of
\eqref{eq:affine-blowdown-far-field} is bounded by
\[
 Cr_k^{sp-p+1-\alpha}
 \left(\frac{T_j}{a_j}+|q_{j,k}|\right)\le C.
\]
The intervening annuli are controlled by the affine increments already
obtained at the preceding levels.  Caccioppoli and local boundedness then
reproduce the remaining hypotheses.  Hence the large-slope improvement
iterates for every \(k\), uniformly in \(j\), and supplies affine maps
with
\[
 \left(\fint_{B_{2r_k}}|v_j-m_{j,k}|^p\dd X\right)^{1/p}
 \le Cr_k^{1+\alpha}.
\]

Fix \(L\ge1\), and choose \(k=k(j,L)\) so that
\(L\le R_jr_k<\rho_\infty^{-1}L\).  Returning to the original variables,
one obtains an affine map \(\ell_{j,L}\) satisfying
\[
 \left(\fint_{B_L}
 |U-b\cdot x-\ell_{j,L}|^p\dd x\right)^{1/p}
 \le Ca_jL^{1+\alpha}R_j^{-\alpha}\longrightarrow0.
\]
Affine functions form a closed finite-dimensional subspace of
\(L^p(B_L)\).  Thus \(U-b\cdot x\) is affine on every ball and hence
globally affine.  Finally \eqref{eq:affine-blowdown-convergence} forces the
additional affine slope to vanish, while the normalization \(U(0)=0\)
eliminates the remaining constant.
\end{proof}

\section{Dilation-recurrent blow-downs and the Liouville theorem}
\label{sec:dilation-liouville}

The analytic ingredients for this section were established in
Section~\ref{sec:large-slope-and-blowdown}.  We use the fixed-order
infinite-slope compactness of Proposition
\ref{prop:fixed-order-infinite-slope-compactness}, the large-slope affine
improvement of Proposition~\ref{prop:large-slope-affine-improvement}, the
noncritical pointwise expansion of Proposition
\ref{prop:noncritical-pointwise-expansion}, and the affine blow-down
rigidity of Proposition~\ref{prop:affine-blowdown-rigidity}.  None of
these results uses the nonlinear Liouville theorem proved below.

We first construct a recurrent centered blow-down and isolate a
record-scale profile on which the global Lipschitz seminorm is attained
at the origin.  We then analyze the singular linearized kernel in a
noncritical ball, using directional coarea to control its crossed part
in a Morrey--Kato class.  The resulting contact propagation closes the
Liouville argument in Subsection~\ref{subsec:liouville-completion}.

\subsection{Centered dilation recurrence}
\label{subsec:centered-dilation-recurrence}

Translation averages cannot detect a sparse sequence of bad centers.  We
therefore turn to the centered dilation action
\[
 (D_rV)(x):=\frac{V(rx)}r,
 \qquad r>0,
\]
on normalized profiles satisfying $V(0)=0$.  Proposition
\ref{prop:affine-blowdown-rigidity} shows that if one centered blow-down
of \(U\) is affine, then \(U\) itself is affine.

For a Lipschitz function $V$, write
\[
 \Lambda_V(r):=[V]_{C^{0,1}(B_r)},
 \qquad
 \Lambda_V(0+):=\lim_{r\downarrow0}\Lambda_V(r).
\]

\begin{proposition}[Dilation-recurrent profile]
\label{prop:dilation-recurrent-cascade}
Let $U$ be a nonaffine,
globally Lipschitz entire solution of \eqref{eq:equation}, normalized by
$U(0)=0$.  Then its centered blow-down hull contains a nonaffine profile
$V$ and a sequence $r_j\downarrow0$ such that
\begin{equation}
 D_{r_j}V\longrightarrow V
 \qquad\text{locally uniformly in }\R^n.
 \label{eq:dilation-recurrence}
\end{equation}
Moreover,
\begin{equation}
 [V]_{C^{0,1}(\R^n)}
 =\Lambda_V(R)=\Lambda_V(0+)=:L_V>0
 \qquad\text{for every }R>0.
 \label{eq:all-centered-scales-full-slope}
\end{equation}
If $a_k,b_k\to0$, $a_k\ne b_k$, and
\begin{equation}
 \frac{|V(a_k)-V(b_k)|}{|a_k-b_k|}\longrightarrow L_V,
 \label{eq:centered-extremizing-chords}
\end{equation}
then necessarily
\begin{equation}
 \frac{|a_k-b_k|}{|a_k|+|b_k|}\longrightarrow0.
 \label{eq:centered-microscopic-ratio}
\end{equation}
In particular, such an extremizing sequence exists, and every one of its
chords is asymptotically microscopic relative to its distance from the
center.
\end{proposition}

\begin{proof}
Set
\[
 U_R(x):=\frac{U(Rx)-U(0)}R.
\]
The family $(U_R)_{R\ge1}$ is equi-Lipschitz and normalized at the origin.
Its set $\Omega_\infty(U)$ of local uniform limits as $R\to\infty$ is
therefore nonempty and compact.  The affine-tail condition
$sp>p-1$ permits passage to the weak equation, so every member of
$\Omega_\infty(U)$ is again an entire solution.  Furthermore,
$\Omega_\infty(U)$ is invariant under every $D_r$, $r>0$: if
$U_{R_i}\to W$, then
\[
 U_{rR_i}=D_rU_{R_i}\longrightarrow D_rW,
 \qquad rR_i\longrightarrow\infty.
\]

Choose a nonempty compact subset $\mathscr M\subset\Omega_\infty(U)$
that is minimal among the closed sets invariant under all dilations.  It
exists by Zorn's lemma, since the intersection of a decreasing chain of
nonempty compact invariant sets is nonempty and invariant.  To make the
recurrence under contractions explicit, write \(T_t=D_{e^t}\).  For
\(V\in\mathscr M\), minimality implies that the orbit
\(\{T_tV:t\in\R\}\) is dense in \(\mathscr M\).  Hence, for every
neighborhood \(\mathcal O\) of \(V\), the family
\(\{T_{-t}\mathcal O:t\in\R\}\) covers \(\mathscr M\); compactness gives a
finite subcover corresponding to \(t_1,\ldots,t_N\).  Given \(T>0\),
apply this cover to \(T_{-T}V\).  For some \(i\),
\(T_{t_i-T}V\in\mathcal O\), and \(t_i-T\to-\infty\) as
\(T\to\infty\).  A diagonal choice over a decreasing neighborhood basis
therefore gives \(s_j\to-\infty\) with \(T_{s_j}V\to V\).  Setting
\(r_j=e^{s_j}\) proves \eqref{eq:dilation-recurrence}.

No element of $\mathscr M$ is affine.  Otherwise it would be an affine
centered blow-down of $U$, and Proposition
\ref{prop:affine-blowdown-rigidity} would make $U$ affine.  In
particular, $V$ is nonconstant.  Put
$L_0=\Lambda_V(0+)$.  For fixed $R>0$ and distinct $x,y\in B_R$,
\eqref{eq:dilation-recurrence} gives
\begin{align*}
 \frac{|V(x)-V(y)|}{|x-y|}
 &=\lim_{j\to\infty}
 \frac{|D_{r_j}V(x)-D_{r_j}V(y)|}{|x-y|}\\
 &\le\lim_{j\to\infty}\Lambda_V(r_jR)=L_0.
\end{align*}
Taking the supremum over $x,y\in B_R$ yields
$\Lambda_V(R)\le L_0$.  The reverse inequality follows from the monotonicity
of $R\mapsto\Lambda_V(R)$.  Letting $R\to\infty$ proves
\eqref{eq:all-centered-scales-full-slope}; its common value is positive
because $V$ is nonconstant.

By \eqref{eq:all-centered-scales-full-slope}, one can choose
$a_k,b_k\in B_{1/k}$ satisfying
\eqref{eq:centered-extremizing-chords}.  It remains to prove that every
such sequence obeys \eqref{eq:centered-microscopic-ratio}.  If not, pass to
a subsequence for which
\[
 |a_k-b_k|\ge c\bigl(|a_k|+|b_k|\bigr)
\]
with $c>0$.  Put $\rho_k=|a_k|+|b_k|$ and
$W_k=D_{\rho_k}V$.  Since $\rho_k\downarrow0$ and $\mathscr M$ is compact
and dilation invariant, a subsequence converges locally uniformly to some
$W\in\mathscr M$.  The rescaled endpoints
$\alpha_k=a_k/\rho_k$ and $\beta_k=b_k/\rho_k$ lie in $\overline B_1$ and
satisfy $|\alpha_k-\beta_k|\ge c$.  After one more subsequence they converge
to distinct points $\alpha,\beta$, and
\eqref{eq:centered-extremizing-chords} becomes
\[
 |W(\alpha)-W(\beta)|=L_V|\alpha-\beta|.
\]
The global Lipschitz seminorm of $W$ is at most $L_V$.  The displayed
equality therefore shows both that its seminorm equals $L_V$ and that it is
attained on a nontrivial secant.  Theorem
\ref{thm:extremal-secant-rigidity} makes $W$ affine, contradicting
$W\in\mathscr M$.  This proves \eqref{eq:centered-microscopic-ratio}.
\end{proof}

\subsection{Collapse of affine radii}
\label{subsec:collapse-affine-radii}

For the rest of the paper, retain the minimal compact dilation hull
$\mathscr M$ selected in the proof of
Proposition~\ref{prop:dilation-recurrent-cascade}, together with the recurrent
profile $V\in\mathscr M$.

For a differentiability point $x\ne0$ of a Lipschitz function $V$, put
$q_x=\nabla V(x)$ and, for $0<\eta<1$, define its relative affine radius by
\begin{equation}
\begin{split}
 \mathfrak r_\eta(V;x):=\sup\biggl\{r\in(0,|x|/2]:\ &
 \sup_{0<\rho\le r}\frac1\rho
 \|V-V(x)-q_x\cdot(\,\cdot-x)\|_{L^\infty(B_\rho(x))}
 \le\eta\biggr\}.
\end{split}
 \label{eq:relative-affine-radius}
\end{equation}
Fr\'echet differentiability makes this number positive.  The following
consequence provides an additional pointwise scale separation for the
profile in Proposition~\ref{prop:dilation-recurrent-cascade}.

\begin{corollary}[Collapse of affine radii at near-extremal points]
\label{cor:collapse-affine-radii}
Let $V$ be the recurrent profile in Proposition
\ref{prop:dilation-recurrent-cascade}.  There are differentiability points
$x_k\to0$ such that, after passing to a subsequence,
\begin{equation}
 q_k:=\nabla V(x_k)\longrightarrow q,
 \qquad |q|=L_V.
 \label{eq:near-extremal-gradients}
\end{equation}
For every such sequence one necessarily has
\begin{equation}
 \frac{\mathfrak r_{1/k}(V;x_k)}{|x_k|}\longrightarrow0.
 \label{eq:affine-radius-collapse}
\end{equation}
More generally, \eqref{eq:affine-radius-collapse} holds with any sequence
$\eta_k\downarrow0$ in place of $1/k$.
\end{corollary}

\begin{proof}
On every ball, the Lipschitz seminorm of a Lipschitz function equals the
essential supremum of the length of its gradient.  Indeed, one inequality
is Rademacher's theorem, and the other follows by restricting the function
to line segments in the convex ball.  Equation
\eqref{eq:all-centered-scales-full-slope} therefore gives
\[
 \operatorname*{ess\,sup}_{B_r}|\nabla V|=L_V
 \qquad(r>0).
\]
We may consequently choose differentiability points $x_k\in B_{1/k}$ such
that $|\nabla V(x_k)|>L_V-1/k$.  Compactness of the closed ball of radius
$L_V$ gives a subsequence satisfying
\eqref{eq:near-extremal-gradients}.

Write $d_k=|x_k|$ and suppose that
\eqref{eq:affine-radius-collapse} fails.  For some $c>0$ and a subsequence,
\[
 \mathfrak r_{1/k}(V;x_k)\ge cd_k.
\]
After decreasing $c$, assume $c<1/2$.  Set
$W_k=D_{d_k}V$ and $\xi_k=x_k/d_k\in\mathbb S^{n-1}$.  The definition
\eqref{eq:relative-affine-radius} yields
\begin{equation}
 \left\|W_k-\frac{V(x_k)}{d_k}
 -q_k\cdot(\,\cdot-\xi_k)\right\|_{L^\infty(B_c(\xi_k))}
 \le\frac ck.
 \label{eq:rescaled-affine-ball}
\end{equation}
Since the minimal hull $\mathscr M$ is compact and dilation invariant,
we may assume
\[
 W_k\longrightarrow W\in\mathscr M
 \quad\text{locally uniformly},
 \qquad \xi_k\to\xi\in\mathbb S^{n-1}.
\]
The constants $V(x_k)/d_k$ are bounded by $L_V$.  Passing to one more
subsequence in \eqref{eq:rescaled-affine-ball} gives
\[
 W(X)=a+q\cdot(X-\xi)
 \qquad(X\in B_c(\xi))
\]
for some $a\in\R$.  Since $|q|=L_V$ and
$[W]_{C^{0,1}(\R^n)}\le L_V$, any sufficiently short segment in
$B_c(\xi)$ parallel to $q$ is an extremal secant for $W$.  Theorem
\ref{thm:extremal-secant-rigidity} makes $W$ affine in $\R^n$, which
contradicts $W\in\mathscr M$.  This proves
\eqref{eq:affine-radius-collapse}.  The same proof applies verbatim to
any $\eta_k\downarrow0$.
\end{proof}

\subsection{Nondegeneracy of record-scale profiles}
\label{subsec:record-scale-nondegeneracy}

We next show that the collapse in
\eqref{eq:affine-radius-collapse} cannot occur with vanishing rescaled
amplitude.  Proposition~\ref{prop:noncritical-pointwise-expansion}
implies the following estimate.  If
$x$ is a differentiability point with $q_x\ne0$, then, for every
$\alpha<sp-p+1$,
\[
 \left(\fint_{B_r(x)}
 |V-V(x)-q_x\cdot(\,\cdot-x)|^p\dd y\right)^{1/p}
 \le C_xr^{1+\alpha}
\]
for all sufficiently small $r$.  Since $V$ is globally Lipschitz, the
standard Lipschitz--$L^p$ interpolation gives
\begin{equation}
 \frac1r\|V-V(x)-q_x\cdot(\,\cdot-x)\|_{L^\infty(B_r(x))}
 \le C'_xr^{\alpha p/(n+p)}.
 \label{eq:noncritical-sup-expansion}
\end{equation}
Indeed, a value of size $A$ in $B_r(x)$ persists on a ball of radius
comparable to $A/L_V$; integration on that ball and the displayed
$L^p$ estimate give $A^{n+p}\le C L_V^n r^{n+p(1+\alpha)}$.

Fix from now on
\begin{equation*}
 0<\beta<\frac{p}{n+p}(sp-p+1).
\end{equation*}
For a differentiability point $x$ with tangent $q_x$, define
\begin{equation*}
 \mathfrak e_{V,x}(r):=
 \frac1r\|V-V(x)-q_x\cdot(\,\cdot-x)\|_{L^\infty(B_r(x))},
 \qquad
 \mathfrak H_{V,x}(r):=r^{-\beta}\mathfrak e_{V,x}(r).
\end{equation*}
By \eqref{eq:noncritical-sup-expansion}, after choosing
$\alpha$ so that $\beta<\alpha p/(n+p)$,
$\mathfrak H_{V,x}(r)\to0$ as $r\downarrow0$.  The global Lipschitz bound
also gives $\mathfrak H_{V,x}(r)\to0$ as $r\to\infty$.  Consequently
$\mathfrak H_{V,x}$ attains its maximum at a positive finite radius.
Here \(r\mapsto\mathfrak e_{V,x}(r)\) is continuous, because it is the
normalized supremum of a continuous function over the nested compact
balls \(\overline B_r(x)\).

\begin{proposition}[Nondegeneracy of the rescaled amplitude]
\label{prop:no-linear-record-neck}
Let $V$ and $x_k$ be as in Corollary~\ref{cor:collapse-affine-radii}, and
let $d_k=|x_k|$ and $q_k=\nabla V(x_k)$.  Choose $t_k>0$ so that
\begin{equation}
 \mathfrak H_{V,x_k}(t_k)
 =\max_{r>0}\mathfrak H_{V,x_k}(r),
 \qquad
 a_k:=\mathfrak e_{V,x_k}(t_k).
 \label{eq:record-scale-choice}
\end{equation}
Then $t_k\to0$.  After passage to a subsequence,
\begin{equation}
 a_k\longrightarrow a_*>0.
 \label{eq:positive-record-amplitude}
\end{equation}
Moreover, the rescaled solutions
\[
 U_k(X):=\frac{V(x_k+t_kX)-V(x_k)}{t_k}
       =q_k\cdot X+a_kv_k(X)
\]
converge locally uniformly to a nonaffine, globally $L_V$-Lipschitz entire
solution
\begin{equation}
 U_*(X)=q\cdot X+a_*v_*(X),
 \qquad |q|=L_V,
 \label{eq:gradient-contact-profile}
\end{equation}
such that
\begin{equation}
 v_*(0)=0,\qquad
 \|v_*\|_{L^\infty(B_1)}=1,\qquad
 \|v_*\|_{L^\infty(B_R)}\le R^{1+\beta}\quad(R>0).
 \label{eq:record-profile-growth}
\end{equation}
In particular, $U_*$ is differentiable at the origin,
\begin{equation}
 \nabla U_*(0)=q,
 \qquad |\nabla U_*(0)|=[U_*]_{C^{0,1}(\R^n)}=L_V.
 \label{eq:attained-maximal-gradient}
\end{equation}
Thus the simultaneous collapse in
\eqref{eq:centered-microscopic-ratio} and
\eqref{eq:affine-radius-collapse} cannot be realized by a linearized
profile with vanishing amplitude.
\end{proposition}

\begin{proof}
We first establish uniform nonaffinity at a scale comparable to $d_k$.
Fix $c\in(0,1/4)$.  We claim that there is
$\varepsilon_0>0$ such that
\begin{equation}
 \mathfrak e_{V,x_k}(cd_k)\ge\varepsilon_0
 \label{eq:outer-record-nonaffinity}
\end{equation}
for all large $k$.  Otherwise, put $W_k=D_{d_k}V$ and
$\xi_k=x_k/d_k\in\mathbb S^{n-1}$.  Compactness and dilation invariance of
$\mathscr M$ give, along a subsequence,
$W_k\to W\in\mathscr M$ locally uniformly and $\xi_k\to\xi$.  The failure
of \eqref{eq:outer-record-nonaffinity}, together with $q_k\to q$, gives
\[
 W(X)=W(\xi)+q\cdot(X-\xi)
 \qquad (X\in B_c(\xi)).
\]
Since $|q|=L_V$ and $[W]_{C^{0,1}(\R^n)}\le L_V$, the ball contains an
extremal secant of $W$.  Theorem~\ref{thm:extremal-secant-rigidity} makes
$W$ affine, contradicting $W\in\mathscr M$.

Equations \eqref{eq:record-scale-choice} and
\eqref{eq:outer-record-nonaffinity} imply
\[
 \mathfrak H_{V,x_k}(t_k)
 \ge\varepsilon_0(cd_k)^{-\beta}\longrightarrow\infty.
\]
In particular, \(a_k>0\) for all sufficiently large \(k\), so all the
normalizations below are well defined.
Since $a_k\le2L_V$, the identity
$t_k^\beta=a_k/\mathfrak H_{V,x_k}(t_k)$ proves $t_k\to0$.
The definition of the maximizing scale gives, for every $R>0$,
\begin{equation}
 \|v_k\|_{L^\infty(B_R)}
 =\frac{R\mathfrak e_{V,x_k}(Rt_k)}{a_k}
 \le R^{1+\beta},
 \qquad
 \|v_k\|_{L^\infty(B_1)}=1.
 \label{eq:record-growth-before-limit}
\end{equation}

Since \(0\le a_k\le2L_V\), pass to a subsequence such that
\(a_k\to a_*\ge0\).  It remains to exclude \(a_*=0\).  Suppose this
happens.  Then
$Q_k=q_k/a_k$ is an infinite tilt and its direction tends to
$e=q/L_V$.  The polynomial bound
\eqref{eq:record-growth-before-limit} has exponent
$1+\beta<sp-p+2$.  We verify that this estimate controls the full shifted
flux tail, including the region in which the linearization is not
uniform.  Put
\[
 \gamma:=sp-p+1,
 \qquad R_k:=a_k^{-1/\beta}.
\]
For \(X\in B_{5/4}\), \(R\ge2\), and
\(Y=X+R\omega\), set
\begin{equation*}
 \mathcal F_k(X,R,\omega):=(1+|Q_k|)^{2-p}
 \left|\Jp\bigl(Q_k\cdot(X-Y)+\delta v_k(X,Y)\bigr)
       -\Jp\bigl(Q_k\cdot(X-Y)\bigr)\right|.
\end{equation*}
Since \(|q_k|\to L_V>0\), one has
\((1+|Q_k|)^{2-p}\asymp a_k^{p-2}\).  On \(2\le R\le R_k\), split the
sphere into
\begin{equation*}
 \bigl\{|e_k\cdot\omega|>C a_kR^\beta\bigr\}
 \quad\hbox{and}\quad
 \bigl\{|e_k\cdot\omega|\le C a_kR^\beta\bigr\},
 \qquad e_k:=q_k/|q_k|.
\end{equation*}
On the first set, the linearized flux estimate and
\eqref{eq:record-growth-before-limit} give
\begin{equation*}
 \mathcal F_k(X,R,\omega)
 \le C R^{p-2}|e_k\cdot\omega|^{p-2}R^{1+\beta}.
\end{equation*}
The angular integral is bounded by \(CR^{p-1+\beta}\), because
\(p-2>-1\).  On the second set, whose measure is at most
\(Ca_kR^\beta\), the global \((p-1)\)-H\"older estimate for \(\Jp\)
gives the same bound:
\begin{align*}
 &a_k^{p-2}R^{(1+\beta)(p-1)}a_kR^\beta\\
 &\qquad=R^{p-1+\beta}(a_kR^\beta)^{p-1}
 \le R^{p-1+\beta}.
\end{align*}
Consequently,
\begin{equation*}
 \sup_{X\in B_{5/4}}
 \int_2^{R_k}\int_{\mathbb S^{n-1}}
 \mathcal F_k(X,R,\omega)\dd\omega\,
 \frac{\dd R}{R^{sp+1}}
 \le C\int_2^\infty R^{\beta-\gamma-1}\dd R<\infty.
\end{equation*}
On \(R>R_k\), use the global Lipschitz bound for
\(U_k=q_k\cdot X+a_kv_k\).  Homogeneity and the global H\"older estimate
for the \(p\)-flux yield
\[
 \int_{R_k}^\infty
 a_k^{-1}R^{p-1}R^{-sp-1}\dd R
 \le C a_k^{\gamma/\beta-1}\longrightarrow0.
\]
Thus the shifted tails are uniformly bounded.  Moreover, for every fixed
\(R\ge2\), the portion exterior to \(B_R\) is bounded, uniformly for all
large \(k\), by
\begin{equation}
 C R^{\beta-\gamma}+o_k(1)
 =C R^{1+\beta-\tau_s}+o_k(1).
 \label{eq:vanishing-record-linearized-tail}
\end{equation}

We make the all-scale content of this estimate explicit.  For \(A\ge2\)
put
\begin{align*}
 \mathfrak T_k(A):=\max\biggl\{&
 {\operatorname*{ess\,sup}}_{X\in B_{5A/4}}
 \int_{\R^n\setminus B_{4A/3}}
 \frac{\lambda_{Q_k}^{2-p}|\mathcal D_{Q_k}v_k(X,Y)|}
 {|X-Y|^{n+sp}}\dd Y,\\
 &{\operatorname*{ess\,sup}}_{X\in B_{3A/2}}
 \int_{\R^n\setminus B_{2A}}
 \frac{\lambda_{Q_k}^{2-p}|\mathcal D_{Q_k}v_k(X,Y)|}
 {|X-Y|^{n+sp}}\dd Y\biggr\}.
\end{align*}
The preceding angular splitting is unchanged when its radial
integration begins at \(cA\).  Indeed, if \(|X|\le2A\) and
\(|X-Y|\ge cA\), then \eqref{eq:record-growth-before-limit} gives
\[
 |\delta v_k(X,Y)|
 \le C\bigl(A+|X-Y|\bigr)^{1+\beta}
 \le C|X-Y|^{1+\beta}.
\]
The part between \(cA\) and \(R_k\) is therefore bounded by
\[
 C\int_{cA}^{R_k}R^{\beta-\gamma-1}\dd R
 \le CA^{\beta-\gamma},
\]
whereas the part beyond \(R_k\) is
\(O(a_k^{\gamma/\beta-1})\).  Consequently,
\begin{equation}
 \limsup_{k\to\infty}\mathfrak T_k(A)
 \le CA^{\beta-\gamma}
 \qquad(A\ge2).
 \label{eq:record-tail-all-scales}
\end{equation}
This contains both exterior regions in
\eqref{eq:large-slope-pointwise-tail}; in particular, no enlargement of
the \(X\)-region is implicit.

For each fixed \(A\ge2\), define
\[
 \widehat v_{k,A}(X):=A^{-1-\beta}v_k(AX),
 \qquad
 \widehat Q_{k,A}:=A^{-\beta}Q_k.
\]
Then \(\widehat v_{k,A}\) solves the equation shifted by
\(\widehat Q_{k,A}\), and
\[
 \|\widehat v_{k,A}\|_{L^\infty(B_R)}\le R^{1+\beta},
 \qquad
 |\widehat Q_{k,A}|\longrightarrow\infty,
 \qquad
 \frac{\widehat Q_{k,A}}{|\widehat Q_{k,A}|}\longrightarrow e.
\]
A change of variables in the two tail components gives
\begin{align}
 \mathcal T^{\rm pt}_{s,\widehat Q_{k,A}}
       (\widehat v_{k,A})
 &=
 \left(\frac{\lambda_{\widehat Q_{k,A}}}{\lambda_{Q_k}}\right)^{2-p}
 A^{sp-(1+\beta)(p-1)}\mathfrak T_k(A)\notag\\
 &\le C+o_k(1).
 \label{eq:scaled-record-tail}
\end{align}
Indeed, the prefactor before \(\mathfrak T_k(A)\) converges to
\(A^{\gamma-\beta}\), which cancels the power in
\eqref{eq:record-tail-all-scales}.  We also need the corresponding
bound with a larger \(X\)-region.  The second component of
\(\mathfrak T_k(2A)\), after the same change of variables, gives
\begin{equation}
 \begin{split}
 &{\operatorname*{ess\,sup}}_{X\in B_3}
 \int_{\R^n\setminus B_4}
 \frac{\lambda_{\widehat Q_{k,A}}^{2-p}
 |\mathcal D_{\widehat Q_{k,A}}\widehat v_{k,A}(X,Y)|}
 {|X-Y|^{n+sp}}\dd Y\\
 &\hspace{35mm}\le C+o_k(1).
 \end{split}
 \label{eq:scaled-record-expanded-tail}
\end{equation}
Indeed, its prefactor is again \(A^{\gamma-\beta}+o_k(1)\), whereas
\eqref{eq:record-tail-all-scales} at \(2A\) is
\(O((2A)^{\beta-\gamma})\).

Choose \(\zeta\in C_c^\infty(B_3)\), \(\zeta=1\) on \(B_2\), and test
the shifted equation with
\(\zeta^2\widehat v_{k,A}\), using a standard bounded truncation and
then letting the truncation level tend to infinity.  Shifted
monotonicity and the product estimate in
Lemma~\ref{lem:angularly-averaged-shifted-modular} give
\begin{align*}
 &\mathcal E_{s,\widehat Q_{k,A}}
       (\widehat v_{k,A};B_2)\\
 &\quad\le C(1-s)
 \iint_{B_4^2}
 \frac{\mathcal H_{\widehat Q_{k,A},X-Y}
 \bigl((|\widehat v_{k,A}(X)|
       +|\widehat v_{k,A}(Y)|)\delta\zeta(X,Y)\bigr)}
 {|X-Y|^{n+sp}}\dd X\!\dd Y\\
 &\qquad+C\|\widehat v_{k,A}\|_{L^\infty(B_3)}
 {\operatorname*{ess\,sup}}_{X\in B_3}
 \int_{\R^n\setminus B_4}
 \frac{\lambda_{\widehat Q_{k,A}}^{2-p}
 |\mathcal D_{\widehat Q_{k,A}}\widehat v_{k,A}(X,Y)|}
 {|X-Y|^{n+sp}}\dd Y.
\end{align*}
The growth bound controls \(\widehat v_{k,A}\) on \(B_4\).  To make
the cutoff estimate uniform in \(k\), write \(X-Y=r\omega\) and
\[
 (|\widehat v_{k,A}(X)|+|\widehat v_{k,A}(Y)|)
 \delta\zeta(X,Y)=r\,t_{k,A,X,r}(\omega).
\]
For fixed \(A\), the preceding local bound and the Lipschitz norm of
\(\zeta\) give \(|t_{k,A,X,r}(\omega)|\le C_A\).  If
\(|\widehat Q_{k,A}|<1\), use the global \(p\)-growth estimate.  If
\(\widehat Q_{k,A}=q_{k,A}e_{k,A}\), \(q_{k,A}\ge1\), use instead
\[
 r^{-p}\mathcal H_{\widehat Q_{k,A},r\omega}
       (r t_{k,A,X,r}(\omega))
 \le C|e_{k,A}\cdot\omega|^{p-2}
       |t_{k,A,X,r}(\omega)|^2.
\]
The angular integral is bounded because \(p-2>-1\).  Consequently,
the spherical integral of the cutoff density is at most \(C_A r^p\),
and
\[
 (1-s)\int_0^8r^{p-sp-1}\dd r\le C(p).
\]
This controls the cutoff term without treating the variable amplitude
as a fixed scalar.  Estimate \eqref{eq:scaled-record-expanded-tail}
controls the crossed term.  Both bounds are independent of \(k\).
Hence
\begin{equation}
 \mathcal E_{s,\widehat Q_{k,A}}
       (\widehat v_{k,A};B_2)\le C_A.
 \label{eq:scaled-record-energy}
\end{equation}
For fixed \(A\), Proposition
\ref{prop:fixed-order-infinite-slope-compactness} now applies to
\(\widehat v_{k,A}\).  Applying it successively for \(A=2^m\),
\(m\in\mathbb N\), and taking a diagonal subsequence yields
\begin{equation}
 v_k\longrightarrow v
 \quad\hbox{locally uniformly in }\R^n.
 \label{eq:record-global-local-uniform-convergence}
\end{equation}
The limits obtained at different values of \(A\) agree on overlaps,
because they arise from the same sequence.

We finally identify the equation without suppressing the exterior
limit.  For each dyadic \(A=2^m\), Proposition
\ref{prop:fixed-order-infinite-slope-compactness} and
\eqref{eq:record-tail-all-scales} give, in \(B_{5A/4}\),
\[
 L^{B_{4A/3}}_{s,e}v=f_A,
 \qquad
 \|f_A\|_{L^\infty(B_{5A/4})}
 \le CA^{\beta-\gamma}.
\]
More explicitly, if \(\widehat f_A\) is the regional source obtained
for \(\widehat v_{k,A}\) on \(B_{5/4}\), then stable scaling gives
\[
 f_A(X)=A^{\beta-\gamma}\widehat f_A(X/A),
 \qquad \|\widehat f_A\|_{L^\infty(B_{5/4})}\le C.
\]
Since \(\beta<\gamma\), \(f_A\to0\) locally uniformly in norm as
\(A\to\infty\).  Moreover, the growth inherited from
\eqref{eq:record-growth-before-limit} gives
\begin{equation}
 \int_{\R^n\setminus B_A}
 \frac{|v(Y)|}{|Y|^{n+\tau_s}}\dd Y
 \le CA^{1+\beta-\tau_s}
 =CA^{\beta-\gamma}.
 \label{eq:record-limit-stable-tail}
\end{equation}
The same decay holds with the spectral density of the limiting
operator.  Indeed, polar coordinates, the growth estimate, and
\(|e\cdot\omega|^{p-2}\in L^1(\mathbb S^{n-1})\) give
\begin{equation*}
 \int_{\R^n\setminus B_A}
 \frac{|e\cdot(Y/|Y|)|^{p-2}|v(Y)|}
 {|Y|^{n+\tau_s}}\dd Y
 \le CA^{\beta-\gamma}.
\end{equation*}
Thus the regional forms converge to the full form on every compactly
supported test function.  Letting \(A\to\infty\) shows that \(v\)
solves the homogeneous translation-invariant stable equation associated
with the kernel
\begin{equation*}
 K_e(z)=\frac{|e\cdot z|^{p-2}}
 {|z|^{n+sp}}
 =\frac{|e\cdot(z/|z|)|^{p-2}}
 {|z|^{n+sp-p+2}}.
\end{equation*}
The angular density is integrable and spectrally elliptic.  Local
uniform convergence in
\eqref{eq:record-global-local-uniform-convergence} preserves both
statements in \eqref{eq:record-growth-before-limit}; thus
\begin{equation}
 v(0)=0,\qquad
 \|v\|_{L^\infty(B_1)}=1,\qquad
 \|v\|_{L^\infty(B_R)}\le R^{1+\beta}.
 \label{eq:nontrivial-linear-record-limit}
\end{equation}

Since \(1+\beta<\tau_s\), the whole-space Liouville theorem for
symmetric stable operators
\cite[Theorem~2.1]{RosOtonSerraStable}, applied with the growth in
\eqref{eq:nontrivial-linear-record-limit}, shows directly that \(v\)
is a polynomial of degree at most one.  Hence \(v\) is affine.  The
bound $|v(X)|\le|X|^{1+\beta}$ as $X\to0$ forces its affine part to
vanish, contradicting
$\|v\|_{L^\infty(B_1)}=1$.  This proves
\eqref{eq:positive-record-amplitude}.

We may now assume $a_k\to a_*>0$.  The functions $U_k$ are globally
$L_V$-Lipschitz, and the equation is stable under local uniform convergence
because $sp>p-1$.  Thus, after a subsequence, $U_k\to U_*$ locally uniformly
and $U_*$ is an entire solution.  Since $a_k$ stays away from zero,
\eqref{eq:record-growth-before-limit} also gives local uniform convergence
$v_k\to v_*$.  Passing to the limit yields
\eqref{eq:gradient-contact-profile} and
\eqref{eq:record-profile-growth}.  The small-radius part of
\eqref{eq:record-profile-growth} gives $\nabla v_*(0)=0$, and hence
\eqref{eq:attained-maximal-gradient}.  Finally, $v_*$ is nonconstant by its
unit norm on $B_1$, so $U_*$ is nonaffine.  The proposition follows.
\end{proof}

\subsection{Noncritical regularity at the maximal-gradient profile}
\label{subsec:noncritical-maximal-gradient}

\begin{corollary}[Dissipation along the contact direction]
\label{cor:tangent-ball-dissipation}
Let $U_*=q\cdot X+a_*v_*$ be the profile in
Proposition~\ref{prop:no-linear-record-neck}, and put $e=q/L_V$.  Then
$v_*\in W^{1,\infty}_{\rm loc}(\R^n)$ and
\begin{equation}
 \partial_ev_*
 \le-\frac{a_*}{2L_V}|\nabla v_*|^2\le0
 \qquad\text{a.e. in }\R^n.
 \label{eq:tangent-ball-dissipation}
\end{equation}
Consequently, for almost every $X\in e^\perp$ and $t_1<t_2$,
\begin{equation}
 \frac{a_*}{2L_V}
 \int_{t_1}^{t_2}|\nabla v_*(X+te)|^2\dd t
 \le v_*(X+t_1e)-v_*(X+t_2e).
 \label{eq:line-dissipation}
\end{equation}
Moreover,
\begin{equation}
 \limsup_{R\to\infty}\frac1R
 \|v_*\|_{L^\infty(B_R)}>0,
 \label{eq:mandatory-linear-far-field}
\end{equation}
and no centered blow-down of $U_*$ is affine.
\end{corollary}

\begin{proof}
Rademacher's theorem and the global Lipschitz bound give
$|q+a_*\nabla v_*|\le L_V=|q|$ almost everywhere.  Squaring this inequality
and dividing by $2a_*L_V$ proves
\eqref{eq:tangent-ball-dissipation}; integration on lines gives
\eqref{eq:line-dissipation}.

If \eqref{eq:mandatory-linear-far-field} failed, then
$U_*(X)-q\cdot X=o(|X|)$ uniformly on balls.  Every centered blow-down would
therefore be the affine map $q\cdot X$.  The affine blow-down rigidity used
in Proposition~\ref{prop:dilation-recurrent-cascade} would make $U_*$
affine, contradicting Proposition~\ref{prop:no-linear-record-neck}.  The
same rigidity theorem shows directly that no individual centered blow-down
can be affine.
\end{proof}

\begin{lemma}[Regional angular compactness at the gradient contact]
\label{lem:regional-angular-compactness}
Let $1<p<2$, $0<\alpha<1$, and let $U\in C^{1,\alpha}(B_{2r_0})$ satisfy
\begin{equation}
 \|\nabla U-q\|_{L^\infty(B_{2r_0})}\le\frac14|q|,
 \qquad q\ne0.
 \label{eq:noncritical-gradient-tube}
\end{equation}
Fix a unit vector $\xi$.  If $x$, $x-\rho\omega$ and their translates by
$h\xi$ belong to $B_{2r_0}$, put
\begin{align*}
 A_0(x,\rho,\omega)&:=U(x)-U(x-\rho\omega),\\
 A_h(x,\rho,\omega)&:=U(x+h\xi)-U(x-\rho\omega+h\xi),
\end{align*}
and
\begin{equation*}
 b_h(x,\rho,\omega)
 :=(p-1)\rho^{2-p}\int_0^1
 |(1-t)A_0+tA_h|^{p-2}\dd t.
\end{equation*}
Then, uniformly for the admissible $x$, $\rho$ and $h$,
\begin{equation}
 c\le b_h(x,\rho,\omega)\quad\text{for a.e. }\omega,
 \qquad
 \int_{\mathbb S^{n-1}}b_h(x,\rho,\omega)\dd\omega\le C.
 \label{eq:regional-angular-ellipticity}
\end{equation}
Moreover, if $|h|\le\rho$, then
\begin{equation}
 \int_{\mathbb S^{n-1}}
 \left|b_h(x,\rho,\omega)
 -(p-1)\left|\frac{A_0(x,\rho,\omega)}{\rho}\right|^{p-2}
 \right|\dd\omega
 \le C|h|^{\alpha(p-1)}.
 \label{eq:regional-angular-convergence}
\end{equation}
Here the constants depend only on $n,p,|q|$ and
$[\nabla U]_{C^\alpha(B_{2r_0})}$, after the geometry has been normalized
so that the displayed segments stay a fixed positive distance from the
boundary.
\end{lemma}

\begin{proof}
Let $\nu=p-2\in(-1,0)$.  We first record the angular estimate used
below.  Suppose that \(F,G\in C^1(\mathbb S^{n-1})\) and
\begin{equation*}
 \max\left\{
 \|F-q\cdot\omega\|_{C^1(\mathbb S^{n-1})},
 \|G-q\cdot\omega\|_{C^1(\mathbb S^{n-1})}
 \right\}\le\frac14|q|.
\end{equation*}
At a zero of either function, one has
\(|q\cdot\omega|\le|q|/4\), and hence the tangential projection of \(q\)
has length at least \(\sqrt{15}|q|/4\).  The tangential gradient of the
function is therefore bounded from below by
\((\sqrt{15}-1)|q|/4\) on its zero set.  The
implicit function theorem, a finite covering of the equator, and
integration in the normal coordinate give
\begin{equation}
 \bigl|\{|F|<t\}\bigr|+\bigl|\{|G|<t\}\bigr|\le Ct
 \qquad(0<t<1).
 \label{eq:angular-sublevel-bound}
\end{equation}
Consequently, if $\|F-G\|_{L^\infty}\le\delta\le1$, splitting into
$\{|F|\le2\delta\}$ and its complement and using the mean value theorem on
the latter yields
\begin{equation}
 \int_{\mathbb S^{n-1}}
 \bigl||F|^\nu-|G|^\nu\bigr|\dd\omega
 \le C\delta^{\nu+1}=C\delta^{p-1}.
 \label{eq:singular-angular-modulus}
\end{equation}
Indeed, the near-zero contribution follows from
\eqref{eq:angular-sublevel-bound} by layer-cake integration, while the other
part is bounded by
$C\delta\int_\delta^1 t^{\nu-1}\dd t\le C\delta^{\nu+1}$.

Set $F_j=A_j/\rho$, $j\in\{0,h\}$.  The fundamental theorem of calculus
gives
\[
 F_0(x,\rho,\omega)
 =\int_0^1\nabla U(x-t\rho\omega)\cdot\omega\dd t,
\]
and the analogous formula for $F_h$.  If
\(\zeta\in T_\omega\mathbb S^{n-1}\), direct differentiation gives
\begin{align*}
 \partial_\zeta F_0(x,\rho,\omega)
 &=\nabla U(x-\rho\omega)\cdot\zeta,\\
 \partial_\zeta F_h(x,\rho,\omega)
 &=\nabla U(x-\rho\omega+h\xi)\cdot\zeta.
\end{align*}
Consequently, \eqref{eq:noncritical-gradient-tube} controls both the
values and the tangential gradients of \(F_0\) and \(F_h\) relative to
\(q\cdot\omega\).  Every convex combination
$(1-t)F_0+tF_h$ therefore has the uniform transversality used in
\eqref{eq:angular-sublevel-bound}.  It follows at once that its
$\nu$-power has a uniformly bounded angular integral.  On the other hand,
the Lipschitz bound on $U$ and $\nu<0$ give the pointwise lower bound in
\eqref{eq:regional-angular-ellipticity}.

Finally, H\"older continuity of the gradient gives
\[
 \|F_h-F_0\|_{L^\infty(\mathbb S^{n-1})}
 \le C|h|^\alpha.
\]
Apply \eqref{eq:singular-angular-modulus} to $F_0$ and each convex
combination $(1-t)F_0+tF_h$, and then integrate in $t$.  This proves
\eqref{eq:regional-angular-convergence}.
\end{proof}

\begin{lemma}[A noncritical neighborhood of the maximal-gradient profile]
\label{lem:record-noncritical-neighborhood}
Let \(U_*=q\cdot X+a_*v_*\) be the profile in Proposition
\ref{prop:no-linear-record-neck}.  If
\begin{equation*}
 \beta<\alpha<sp-p+1,
\end{equation*}
then \(U_*\in C^{1,\alpha}(B_{2r_0})\) for some \(r_0>0\).  After
decreasing \(r_0\),
\begin{equation}
 \|\nabla U_*-q\|_{L^\infty(B_{2r_0})}\le\frac14|q|.
 \label{eq:record-noncritical-gradient-tube}
\end{equation}
\end{lemma}

\begin{proof}
Fix a small \(r>0\) and a center \(x\in B_r\).  Subtract
\[
 \ell_x(y):=U_*(x)+q\cdot(y-x)
\]
and use the initial amplitude \(H_0=C_2a_*r^\beta\), where \(C_2\) will
be fixed independently of \(x\) and \(r\).  Define
\begin{equation*}
 z_{x,r}(X):=
 \frac{U_*(x+rX)-\ell_x(x+rX)}{rH_0},
 \qquad Q_{x,r}:=\frac q{H_0}.
\end{equation*}
Then \(z_{x,r}\) solves the equation shifted by \(Q_{x,r}\), in the
sense defined at the beginning of
Section~\ref{sec:large-slope-and-blowdown}.  The growth estimate
\eqref{eq:record-profile-growth} gives
\begin{equation}
 \|U_*-\ell_x\|_{L^\infty(B_{4r}(x))}
 \le C a_*r^{1+\beta}.
 \label{eq:record-translated-affine-tube}
\end{equation}
In particular, \(\|z_{x,r}\|_{L^\infty(B_4)}\le C/C_2\).

We next verify the exterior estimate uniformly in the center.  For
\(t\ge r\), \(x\in B_r\), and \(y\in B_{2t}\),
\eqref{eq:record-profile-growth} yields
\[
 |U_*(y)-q\cdot y-U_*(x)+q\cdot x|
 \le C a_*t^{1+\beta}.
\]
Let \(t_c=(|q|/a_*)^{1/\beta}\), decreasing \(r\) so that
\(r<t_c/4\).  On the annuli \(t_j=2^jr\le t_c\), the angular decomposition used in
the proof of Proposition~\ref{prop:no-linear-record-neck} gives the
relative-flux contribution
\begin{equation*}
 C\left(\frac r{t_j}\right)^{sp-p+1-\beta}.
\end{equation*}
For \(t>t_c\), the global Lipschitz bound for \(U_*\), homogeneity of
\(\Jp\), and radial integration give
\begin{equation*}
 C\frac{r^{sp-p+1}}{H_0}
 \int_{t_c}^\infty t^{p-2-sp}\dd t
 \le C\left(\frac r{t_c}\right)^{sp-p+1-\beta}.
\end{equation*}
Hence, with \(\gamma=sp-p+1\),
\begin{equation}
 \begin{split}
 &{\operatorname*{ess\,sup}}_{X\in B_{5/4}}
 \int_{\R^n\setminus B_{4/3}}
 \frac{(1+|Q_{x,r}|)^{2-p}}
 {|X-Y|^{n+sp}}\\
 &\quad\times
 \left|\Jp\bigl(Q_{x,r}\cdot(X-Y)
                   +\delta z_{x,r}(X,Y)\bigr)
       -\Jp\bigl(Q_{x,r}\cdot(X-Y)\bigr)\right|\dd Y
 \le C\sum_{j\ge0}2^{-j(\gamma-\beta)}+C\le C.
 \end{split}
 \label{eq:record-uniform-relative-tail}
\end{equation}
The constant is independent of \(x\) and \(r\).  The shifted
flux over \(X\in B_{3/2}\), \(Y\notin B_2\) satisfies the same bound:
the two sets are separated and the identical annular decomposition
starts at a different fixed radius.  Thus both components of
\(\mathcal T^{\rm pt}_{s,Q_{x,r}}(z_{x,r})\) are controlled.
The shifted
Caccioppoli inequality, applied after
\eqref{eq:record-uniform-relative-tail}, now bounds the normalized local
Bregman energy.  Increasing \(C_2\) if necessary, the local
\(L^p\) and energy quantities satisfy the unit bounds in Proposition
\ref{prop:large-slope-affine-improvement}; the tail bound is the fixed
constant in \eqref{eq:record-uniform-relative-tail}.

We now specify the iteration and its tail budget.  Choose
\(\bar\alpha\) with
\[
 \alpha<\bar\alpha<\gamma.
\]
Let \(C_T\) be the uniform constant in
\eqref{eq:record-uniform-relative-tail} for the two tail geometries and
set \(L_{\rm rec}=2\max\{1,C_T\}\).  Corollary
\ref{cor:prescribed-large-slope-block-main}, with
\(L=L_{\rm rec}\), first fixes \(C_{\rm L}\); enlarge it once to
dominate the accompanying affine coefficient bound.  After that
choice, let \(C_0\) dominate the constants in the angular flux and
affine-base estimates for both tail geometries, and fix
a dyadic \(\rho_*\in(0,1/16)\) so small that
\begin{align}
 C_{\rm L}\rho_*^{\bar\alpha-\alpha}&\le\frac12,
 \label{eq:record-iteration-endpoint}\\
 C_0\bigl[\rho_*^{\bar\alpha-\alpha}
          +\rho_*^{\gamma-\alpha}\bigr]
 +C_0L_{\rm rec}\rho_*^{\gamma-\alpha}
 &\le L_{\rm rec}.
 \label{eq:record-iteration-tail-budget}
\end{align}
The prescribed-scale corollary then supplies a threshold \(Q_{\rm rec}\).
Finally decrease \(r\), uniformly for \(x\in B_r\), so that
\begin{equation}
 |Q_{x,r}|=\frac{|q|}{H_0}\ge2Q_{\rm rec},
 \qquad
 \frac{C_{\rm L}H_0}{1-\rho_*^\alpha}\le\frac14|q|.
 \label{eq:record-slope-stability}
\end{equation}

For completeness, define every normalized state.  Put
\[
 r_k=r\rho_*^k,\qquad H_k=H_0\rho_*^{k\alpha},
\]
let \(\ell_0=\ell_x\), and, writing \(q_k=\nabla\ell_k\), set
\begin{equation}
 v_k(X):=
 \frac{U_*(x+r_kX)-\ell_k(x+r_kX)}{r_kH_k},
 \qquad Q_k:=\frac{q_k}{H_k}.
 \label{eq:record-explicit-normalized-state}
\end{equation}
The preceding \(L^p\), energy, and tail estimates give at \(k=0\)
\[
 \fint_{B_2}|v_0|^p\dd X\le1,\qquad
 \mathcal E_{s,Q_0}(v_0;B_2)\le1,\qquad
 \mathcal T^{\rm pt}_{s,Q_0}(v_0)\le L_{\rm rec}.
\]
Suppose these bounds hold at level \(k\).  The prescribed-scale
large-slope estimate gives \(m_k=c_k+b_k\cdot X\), with
\(|c_k|+|b_k|\le C_{\rm L}\), and an admissible intrinsic amplitude
\(a_k<\eta_*:=\rho_*^\alpha\).  Define
\begin{align*}
 \ell_{k+1}(y)
 &:=\ell_k(y)+r_kH_k
 m_k\left(\frac{y-x}{r_k}\right),\\
 q_{k+1}&=q_k+H_kb_k,\qquad
 v_{k+1}(X)
 :=\frac{v_k(\rho_*X)-m_k(\rho_*X)}
 {\rho_*\eta_*},\\
 Q_{k+1}&=\frac{Q_k+b_k}{\eta_*}.
\end{align*}
After rescaling, the zero-order modular is bounded by
\((a_k/\eta_*)^p\).  If \(q_k'=|Q_k+b_k|\), exact homogeneity gives
\[
 \mathcal E_{s,Q_{k+1}}(v_{k+1};B_2)
 \le
 \left(\frac{q_k'+\eta_*}{q_k'+a_k}\right)^{2-p}
 \left(\frac{a_k}{\eta_*}\right)^2
 \le\left(\frac{a_k}{\eta_*}\right)^p\le1.
\]

It remains to record the tail step.  By
\eqref{eq:record-slope-stability},
\[
 |q_k-q|\le C_{\rm L}\sum_{i<k}H_i\le\frac14|q|,
\]
so every affine base lies in the same large-slope cone.  Keeping the
factor \(\lambda_{Q_{k+1}}^{2-p}\), the angular estimate on the
annuli processed at this step and the exact old-tail rescaling give
\[
 \lambda_{Q_{k+1}}^{2-p}
 \frac{|\mathcal D_{Q_{k+1}}v_{k+1}(X,Y)|}
 {|X-Y|^{n+sp}}\dd Y
 =
 \rho_*^\gamma\eta_*^{1-p}\lambda_{Q_{k+1}}^{2-p}
 \frac{|\mathcal D_{Q_k+b_k}h_k(Z,W)|}
 {|Z-W|^{n+sp}}\dd W,
\]
where \(h_k=v_k-m_k\), \(Z=\rho_*X\), and \(W=\rho_*Y\).  In
particular,
\[
 \rho_*^\gamma\eta_*^{1-p}
 \left(\frac{\lambda_{Q_{k+1}}}{\lambda_{Q_k}}\right)^{2-p}
 \le C\rho_*^\gamma\eta_*^{-1}
 =C\rho_*^{\gamma-\alpha}
\]
and hence
\begin{equation}
 \mathcal T^{\rm pt}_{s,Q_{k+1}}(v_{k+1})
 \le C_0\bigl[\rho_*^{\bar\alpha-\alpha}
              +\rho_*^{\gamma-\alpha}\bigr]
 +C_0\rho_*^{\gamma-\alpha}
   \mathcal T^{\rm pt}_{s,Q_k}(v_k).
 \label{eq:record-explicit-tail-recurrence}
\end{equation}
Indeed, the angular contribution inherited from a level \(i<k\) is
\[
 C\left(\frac{r_k}{r_i}\right)^\gamma\frac{H_i}{H_k}
 =C\rho_*^{(k-i)(\gamma-\alpha)};
\]
the spectral factor \(|e\cdot\omega|^{p-2}\) is integrable, and the
unprocessed far field has the same endpoint factor
\(\rho_*^{\gamma-\alpha}\).  The two endpoint terms at the new annuli
are those displayed in
\eqref{eq:record-explicit-tail-recurrence}.  Thus
\eqref{eq:record-iteration-tail-budget} preserves the budget
\(L_{\rm rec}\).  At the same time,
\(|q_k|\ge3|q|/4\) and \(H_k\le H_0\) imply
\(|Q_k|\ge Q_{\rm rec}\), so the induction never leaves the
large-slope regime.

The endpoint \(L^\infty\) estimate and the affine increments now give,
uniformly for \(x\in B_r\),
\begin{equation*}
 \inf_{m\in\mathcal A}
 \|U_*-m\|_{L^\infty(B_t(x))}
 \le C H_0r^{-\alpha}t^{1+\alpha}
 \qquad(0<t\le r/2).
\end{equation*}
The Campanato characterization gives
\(U_*\in C^{1,\alpha}(B_{r/2})\); here one uses the usual
overlapping-ball argument.  Set \(r_0=r/4\).  Since
\(\nabla U_*(0)=q\), continuity of the gradient gives
\eqref{eq:record-noncritical-gradient-tube} after a final decrease of
the radius.
\end{proof}

\subsection{Morrey--Kato propagation and weak linearization}
\label{subsec:morrey-kato-linearization}

Lemma~\ref{lem:record-noncritical-neighborhood} verifies the hypotheses
of Lemma~\ref{lem:regional-angular-compactness} for \(U_*\), with
\eqref{eq:noncritical-gradient-tube} replaced by the identical estimate
\eqref{eq:record-noncritical-gradient-tube}.

\begin{proposition}[Estimate of the exterior interaction]
\label{prop:crossed-sublinear-absorption}
Let $U_*$, $e=q/L_V$ and $r_0$ be as above.  For $0<h<r_0/4$, set
\begin{equation*}
 V_h(X):=U_*(X+he)-L_Vh,
 \qquad Z_h:=U_*-V_h.
\end{equation*}
Then $V_h$ is an entire solution, $0\le Z_h\le2L_Vh$ in $\R^n$, and
the regional secant kernel in $B_{r_0}$ satisfies the uniform stable bounds
of order $sp-p+2$ in Lemma
\ref{lem:regional-angular-compactness}.  Moreover, for every
$\varphi\in C_c^\infty(B_{r_0/2})$, $\varphi\ge0$,
\begin{equation}
 \begin{split}
 &(1-s)\iint_{B_{r_0}\times B_{r_0}}
 \frac{c_h(X,Y)(Z_h(X)-Z_h(Y))
       (\varphi(X)-\varphi(Y))}
 {|X-Y|^{n+sp}}\dd X\!\dd Y\\
 &\hspace{35mm}\ge
 -C\int_{B_{r_0/2}}Z_h^{p-1}\varphi\dd X,
 \end{split}
 \label{eq:regional-sublinear-absorption}
\end{equation}
where
\begin{equation*}
 c_h(X,Y):=(p-1)\int_0^1
 |(1-t)\delta V_h(X,Y)+t\delta U_*(X,Y)|^{p-2}\dd t.
\end{equation*}
On the set where both increments vanish, the integrand and \(c_h\)
are defined to be zero.  This convention does not alter the secant
identity, whose two sides vanish there.
The power $p-1$ in \eqref{eq:regional-sublinear-absorption} is sharp under
the sole global Lipschitz assumption.
\end{proposition}

\begin{proof}
Translation invariance and invariance under addition of constants show that
$V_h$ is a solution.  The global Lipschitz bound gives
\[
 U_*(X+he)-L_Vh\le U_*(X)
\]
and also $Z_h\le2L_Vh$.  Subtract the weak equations of $U_*$ and $V_h$.
For pairs in $B_{r_0}\times B_{r_0}$, the fundamental theorem of calculus
for $\Jp$ gives
\[
 \Jp(\delta U_*)-\Jp(\delta V_h)
 =c_h\,\delta Z_h.
\]
The angular bounds for $c_h$ are precisely
\eqref{eq:regional-angular-ellipticity}, after multiplication by
$|X-Y|^{2-p}$.

It remains to estimate crossed pairs.  If
$X\in B_{r_0/2}$ and $Y\notin B_{r_0}$, write
$a=\delta V_h(X,Y)$ and
$\delta U_*(X,Y)=a+Z_h(X)-Z_h(Y)$.  Since $Z_h(Y)\ge0$ and $\Jp$ is
increasing,
\[
 \Jp(\delta U_*)-\Jp(\delta V_h)
 \le \Jp(a+Z_h(X))-\Jp(a)
 \le C Z_h(X)^{p-1}.
\]
The last inequality is the global $(p-1)$-H\"older continuity of $\Jp$.
The two sets are separated, and hence integration of
$|X-Y|^{-n-sp}$ over $Y\notin B_{r_0}$ gives a finite constant.  Moving the
crossed term to the other side of the regional identity proves
\eqref{eq:regional-sublinear-absorption}.  Finally, the scalar estimate
$|\Jp(a+t)-\Jp(a)|\le C|t|^{p-1}$ is optimal uniformly in $a$, as is seen
at $a=0$; thus the exponent cannot be improved without information on the
exterior level sets of $U_*$.  
\end{proof}

The exterior coefficient is estimated by integrating in the contact
direction before passing to the first difference quotient.  We use the
following local strong minimum principle, which does not require a
pointwise L\'evy bound for the crossed kernel.

For a symmetric kernel \(K\) on \(D\times D\), let
\(\mathcal F_K(D)\) denote the closure of \(C_c^\infty(D)\) under the
norm
\[
 \|\phi\|_{\mathcal F_K(D)}^2
 :=\|\phi\|_{L^2(D)}^2
 +\iint_{D\times D}K(x,y)|\delta\phi(x,y)|^2\dd x\!\dd y.
\]
We write \(z\in\mathcal F_{K,\mathrm{loc}}(D)\) if
\(\eta z\in\mathcal F_K(D)\) for every
\(\eta\in C_c^\infty(D)\), and use
\(\mathcal F_{K,c}(D)\) for elements of \(\mathcal F_K(D)\) with
compact support in \(D\).
Whenever an element of \(\mathcal F_{K,c}(D)\) is inserted into a
full-space integral, it is understood to be extended by zero on
\(\R^n\setminus D\).

\begin{lemma}[Contact propagation with a Morrey--Kato absorption]
\label{lem:kato-contact-propagation}
Let \(n\ge2\), let \(D\subset\R^n\) be a ball, let
\(1<\tau<2\), and let \(K\) be a symmetric measurable kernel in
\(D\times D\).  The same statement holds when \(D\) is a connected
convex open set.  Suppose that, in polar coordinates about every
$x\in D$,
\begin{equation}
 K(x,x+\rho\omega)
 =\frac{a(x,\rho,\omega)}{\rho^{n+\tau}},
 \qquad
 a\ge\lambda>0,\qquad
 \int_{\mathbb S^{n-1}}a(x,\rho,\omega)\dd\omega\le\Lambda
 \label{eq:spectral-regional-bounds}
\end{equation}
for almost every \((\rho,\omega)\) such that
\(x+\rho\omega\in D\).
It is enough that these bounds hold locally, with uniform constants on
the connected subdomain in which contact is to be propagated.
Let $A\ge0$ satisfy, locally uniformly in $D$,
\begin{equation}
 \int_{B_r(x)}A\dd y\le M r^{n-\kappa},
 \qquad 0<\kappa<\tau.
 \label{eq:morrey-absorption}
\end{equation}
Assume that \(z\ge0\), that
\(z\in C(D)\cap L^\infty(D)\cap\mathcal F_{K,\mathrm{loc}}(D)\), and
assume that
\begin{equation}
 \iint_{D\times D}K(x,y)\delta z(x,y)\delta\varphi(x,y)
 \dd x\!\dd y
 \ge-\int_DAz\varphi\dd x
 \label{eq:regional-supersolution-with-potential}
\end{equation}
for every nonnegative
\(\varphi\in\mathcal F_{K,c}(D)\).  If $z$ vanishes at one
point of $D$, then $z\equiv0$ in $D$.
The conclusion is local in the following sense: if
\eqref{eq:regional-supersolution-with-potential} is known only for tests
supported in an open set \(D_0\Subset D\), then every zero of \(z\) in
\(D_0\) has a neighborhood on which \(z\) vanishes.
\end{lemma}

\begin{proof}
It suffices to prove the local assertion.  Fix a ball
\(B_{4R_0}\Subset D\), with \(R_0\) small enough that
\eqref{eq:morrey-absorption} holds with the same constant \(M\) on all
balls under consideration.  All constants below may depend on this
fixed localization, but not on the regularizing parameter introduced
below.

We first record the form estimate for the potential.  Restrict
\(\dd\mu=A\dd x\) to \(B_{3R_0}\).  The Morrey bound and a dyadic
decomposition give
\begin{equation}
 \sup_{x\in B_{3R_0}}
 \sum_{j=0}^{\infty}(2^{-j}R)^{\tau-n}
 \mu(B_{2^{-j}R}(x))
 \le CMR^{\tau-\kappa},
 \qquad 0<R\le R_0.
 \label{eq:morrey-to-kato}
\end{equation}
For a ball \(B_R\) used below, put
\(\mu_R:=\mu\lfloor B_R\).  Since \(0<\kappa<\tau<n\), Adams' trace
inequality \cite{AdamsTrace}, applied with Riesz order \(\tau/2\),
domain exponent \(2\), and target exponent
\[
 q_\mu:=\frac{2(n-\kappa)}{n-\tau}>2,
\]
gives
\[
 \|\psi\|_{L^{q_\mu}(\dd\mu_R)}^2
 \le CM^{2/q_\mu}
 [\widetilde\psi]_{\dot H^{\tau/2}(\R^n)}^2.
\]
Indeed, the Morrey norm of \(\mu_R\), viewed as a measure of dimension
\(n-\kappa\), is bounded by \(CM\).  H\"older's inequality and
\(\mu(B_R)\le MR^{n-\kappa}\), together with
\((n-\kappa)(1-2/q_\mu)=\tau-\kappa\), therefore yield
\begin{equation}
 \int_{B_R}A\psi^2\dd x
 \le CMR^{\tau-\kappa}
 [\widetilde\psi]_{\dot H^{\tau/2}(\R^n)}^2
 \label{eq:kato-form-bound}
\end{equation}
for \(\psi\) supported in a ball \(B_R\) whose double is contained in
\(B_{3R_0}\), where \(\widetilde\psi\) denotes its zero extension.

We spell out the passage to an infinitesimal form bound.  Fix
\(\psi\in C_c^\infty(B_{2R_0})\), and let
\(\{\zeta_i\}\) be a smooth partition of unity subordinate to a
covering of \(B_{2R_0}\) by balls of radius \(\rho\), chosen so that
the doubled balls are contained in \(B_{3R_0}\), with bounded overlap
and \(\sum_i\zeta_i^2=1\).  The fractional IMS estimate gives
\[
 \sum_i[\zeta_i\psi]_{\dot H^{\tau/2}(\R^n)}^2
 \le C[\widetilde\psi]_{\dot H^{\tau/2}(\R^n)}^2
      +C\rho^{-\tau}\|\psi\|_2^2.
\]
For completeness, this follows from
\[
 \sum_i|\zeta_i(x)\psi(x)-\zeta_i(y)\psi(y)|^2
 \le2|\delta\psi(x,y)|^2
 +2|\psi(y)|^2\sum_i|\zeta_i(x)-\zeta_i(y)|^2
\]
and the bounded-overlap estimate
\[
 \sum_i|\zeta_i(x)-\zeta_i(y)|^2
 \le C\min\left\{1,\frac{|x-y|^2}{\rho^2}\right\}.
\]
Integration against \(|x-y|^{-n-\tau}\dd x\!\dd y\) gives the stated
remainder \(C\rho^{-\tau}\|\psi\|_2^2\).
Applying \eqref{eq:kato-form-bound} to each \(\zeta_i\psi\), summing,
and using bounded overlap, we obtain
\[
 \int A\psi^2\dd x
 \le CM\rho^{\tau-\kappa}
       [\widetilde\psi]_{\dot H^{\tau/2}(\R^n)}^2
      +CM\rho^{-\kappa}\|\psi\|_2^2.
\]
For later use we make the zero-extension term explicit.  If
\(d_0=\operatorname{dist}(\operatorname{supp}\psi,\partial D)>0\),
then the lower spectral bound in
\eqref{eq:spectral-regional-bounds} and direct integration outside
\(D\) give
\begin{equation}
 [\widetilde\psi]_{\dot H^{\tau/2}(\R^n)}^2
 \le C\lambda^{-1}\mathcal E_K(\psi,\psi)
 +Cd_0^{-\tau}\|\psi\|_2^2.
 \label{eq:kato-zero-extension-localization}
\end{equation}
Indeed, the part over \(D\times D\) is controlled by the regional
form, while for \(x\in\operatorname{supp}\psi\),
\[
 \int_{\R^n\setminus D}|x-y|^{-n-\tau}\dd y
 \le Cd_0^{-\tau}.
\]
Choose
\(\rho=c(\epsilon/(1+M))^{1/(\tau-\kappa)}\), with \(c\) fixed small
and with \(\rho\) capped by the preceding geometric radius.  Insert
\eqref{eq:kato-zero-extension-localization} in the partition estimate.
The energy coefficient is at most \(\epsilon\), and the remaining
terms are zeroth order.  Hence
\begin{equation}
 \int A\psi^2\dd x
 \le \epsilon\mathcal E_K(\psi,\psi)
 +C_*(1+M)^{\tau/(\tau-\kappa)}
       \epsilon^{-\kappa/(\tau-\kappa)}\|\psi\|_2^2,
 \qquad 0<\epsilon\le1,
 \label{eq:kato-infinitesimal-form-bound}
\end{equation}
where
\[
 \mathcal E_K(\psi,\psi)
 :=\iint_{D\times D}K(x,y)|\psi(x)-\psi(y)|^2\dd x\!\dd y.
\]

We shall also use the following cutoff bound.  If \(0\le\eta\le1\),
\(\eta\in C_c^{0,1}(B_R)\), and
\(\|\nabla\eta\|_\infty\le C/d\), then
\begin{equation}
 \mathcal K_\eta(x):=
 \int_DK(x,y)|\eta(x)-\eta(y)|^2\dd y
 \le C\Lambda d^{-\tau}
 \qquad \bigl(x\in D\bigr).
 \label{eq:kato-cutoff-kernel-bound}
\end{equation}
Indeed, after extending \(\eta\) by zero outside \(D\), one has
\(|\eta(x)-\eta(y)|\le C\min\{1,|x-y|/d\}\).  The angular upper bound
then gives
\[
 \mathcal K_\eta(x)
 \le C\Lambda\int_0^\infty
 \min\left\{1,\frac{\rho^2}{d^2}\right\}
 \frac{\dd\rho}{\rho^{1+\tau}}
 \le C\Lambda d^{-\tau}.
\]

Put \(v=z+\varepsilon\), where \(\varepsilon>0\).  Let
\(B_r\Subset B_R\Subset B_{4R_0}\), with \(R/2\le r<R\), choose
\(\eta\in C_c^{0,1}(B_R)\), \(0\le\eta\le1\), so that
\(\eta=1\) on \(B_r\) and
\(\|\nabla\eta\|_\infty\le C/(R-r)\), and put
\[
 S=\{\eta>0\},\qquad m=1+2q,\qquad \psi=\eta v^{-q}.
\]
For every fixed \(q_0>0\), the elementary negative-power inequality
gives, for \(q\ge q_0\),
\begin{align}
 &(a-b)(\xi^2a^{-1-2q}-\zeta^2b^{-1-2q})\notag\\
 &\quad\le-\gamma_q
 |\xi a^{-q}-\zeta b^{-q}|^2
 +C_q(a^{-2q}+b^{-2q})|\xi-\zeta|^2,
 \label{eq:kato-interior-negative-power-scalar}
\end{align}
whenever \(a,b>0\) and \(0\le\xi,\zeta\le1\), where
\begin{equation*}
 0<\gamma_q\le1,\qquad
 \gamma_q\ge\frac{c_{q_0}}{1+q},
 \qquad C_q\le C_{q_0}(1+q)^2.
\end{equation*}
To see this, one starts from
\[
 (a-b)(a^{-1-2q}-b^{-1-2q})
 \le-\frac{1+2q}{q^2}(a^{-q}-b^{-q})^2
\]
and applies Young's inequality to the two cutoff terms.  Thus the
coercive constant is a fixed multiple of
\((1+2q)/q^2\); the exact constant from the uncut inequality is not
needed.

The pairs crossing the cutoff require a separate argument.  If
\(x\in S\) and \(y\in D\setminus S\), then the test vanishes at \(y\).
Moreover,
\begin{equation*}
 (v(x)-v(y))\eta(x)^2v(x)^{-1-2q}
 \le
 \begin{cases}
  0, & v(x)\le v(y),\\
  \eta(x)^2v(x)^{-2q}, & v(x)>v(y).
 \end{cases}
\end{equation*}
Consequently, after decreasing \(\gamma_q\) by a fixed factor,
\begin{align}
 &(v(x)-v(y))\eta(x)^2v(x)^{-1-2q}\notag\\
 &\quad\le-\gamma_q|\psi(x)-\psi(y)|^2
 +(1+\gamma_q)v(x)^{-2q}
       |\eta(x)-\eta(y)|^2.
 \label{eq:kato-crossed-negative-power-scalar}
\end{align}
In particular, no negative power of \(v(y)\) occurs outside \(S\).

Use bounded energy approximations of
\(\eta^2v^{-1-2q}\) in
\eqref{eq:regional-supersolution-with-potential}.  Since \(z\le v\),
\[
 \int_DAz\eta^2v^{-1-2q}\dd x
 \le\int_DA\psi^2\dd x.
\]
Integrating \eqref{eq:kato-interior-negative-power-scalar} on
\(S\times S\), using
\eqref{eq:kato-crossed-negative-power-scalar} on
\(S\times(D\setminus S)\), and exploiting symmetry, we obtain
\begin{equation}
 \gamma_q\mathcal E_K(\psi,\psi)
 \le C_q\int_Sv^{-2q}\mathcal K_\eta\dd x
      +\int_SA\psi^2\dd x.
 \label{eq:kato-negative-caccioppoli-preabsorption}
\end{equation}
Apply \eqref{eq:kato-infinitesimal-form-bound} with
\(\epsilon=\gamma_q/2\).  The energy term is absorbed on the left.
Using \eqref{eq:kato-cutoff-kernel-bound}, and recalling that
\(\gamma_q\asymp(1+q)^{-1}\), gives
\begin{equation}
 \mathcal E_K(\psi,\psi)
 \le C(1+q)^\nu
 \bigl((R-r)^{-\tau}+C_0\bigr)
 \int_{B_R}v^{-2q}\dd x,
 \qquad
 \nu=\max\left\{3,\frac{\tau}{\tau-\kappa}\right\}.
 \label{eq:kato-negative-energy}
\end{equation}
Here \(C_0\) depends only on the fixed localization and the constants in
the hypotheses.  The lower spectral bound,
\eqref{eq:kato-zero-extension-localization}, and the fractional
Sobolev inequality now yield, with \(\chi=n/(n-\tau)>1\),
\begin{equation}
 \left(\fint_{B_r}v^{-2q\chi}\dd x\right)^{1/\chi}
 \le C(1+q)^\nu
 \left[1+\left(\frac{R}{R-r}\right)^\tau\right]
 \fint_{B_R}v^{-2q}\dd x.
 \label{eq:kato-negative-power-step}
\end{equation}
Notice that \(\tau<n\), since \(n\ge2\) and \(\tau<2\).

Let \(q_{j+1}=\chi q_j\), and choose nested radii decreasing from
\(3R/4\) to \(R/2\).  After taking the power \(1/(2q_j)\) in the
\(j\)-th inequality, the logarithms of both the cutoff factor and
\((1+q_j)^\nu\) are summable.  Therefore
\begin{equation}
 \operatorname*{ess\,sup}_{B_{R/2}}v^{-1}
 \le C\left(\fint_{B_{3R/4}}v^{-2q_0}\dd x\right)^{1/(2q_0)}
 \qquad(q_0>0).
 \label{eq:kato-reverse-moser}
\end{equation}

It remains to obtain an initial negative moment.  Use the logarithmic
test \(\eta^2/v\).  On \(S\times S\), the scalar inequality
\begin{align*}
 &(a-b)(\xi^2a^{-1}-\zeta^2b^{-1})\\
 &\quad\le-c\min\{\xi^2,\zeta^2\}
       |\log a-\log b|^2+C|\xi-\zeta|^2
\end{align*}
follows from
\((a-b)(a^{-1}-b^{-1})=-(a-b)^2/(ab)\),
\(|\log a-\log b|^2\le(a-b)^2/(ab)\), and Young's inequality.
For \(x\in S\), \(y\notin S\), split once more according to whether
\(v(x)\le v(y)\) or \(v(x)>v(y)\).  The crossed contribution is
nonpositive in the first case and at most \(\eta(x)^2\) in the second.
Since \(z/v\le1\), the cutoff bound and the Morrey estimate imply, on
every sufficiently small ball,
\begin{equation}
 \iint_{B_{3R/4}\times B_{3R/4}}
 \frac{|\log v(x)-\log v(y)|^2}{|x-y|^{n+\tau}}
 \dd x\!\dd y
 \le C R^{n-\tau}+CMR^{n-\kappa}
 \le C R^{n-\tau}.
 \label{eq:kato-logarithmic-energy}
\end{equation}
The last inequality uses \(\kappa<\tau\) and \(R\le R_0\).
Repeating the estimate on each subball, fractional Poincar\'e gives,
for every ball \(B_\varrho\) in the smaller concentric ball,
\[
 \fint_{B_\varrho}
 \left|\log v-(\log v)_{B_\varrho}\right|\dd x
 \le C\varrho^{(\tau-n)/2}
 \left(
 \iint_{B_\varrho^2}
 \frac{|\delta\log v(x,y)|^2}{|x-y|^{n+\tau}}\dd x\!\dd y
 \right)^{1/2}
 \le C.
\]
Thus \(\log v\) has a local BMO seminorm bounded independently of
\(\varepsilon\).  The John--Nirenberg inequality then
gives \(\theta>0\), also independent of \(\varepsilon\), such that
\begin{equation}
 \left(\fint_{B_{3R/4}}v^\theta\dd x\right)
 \left(\fint_{B_{3R/4}}v^{-\theta}\dd x\right)\le C.
 \label{eq:kato-positive-negative-moment}
\end{equation}
Taking \(q_0=\theta/2\) in \eqref{eq:kato-reverse-moser} and using
\eqref{eq:kato-positive-negative-moment} yields
\begin{equation}
 \left(\fint_{B_{3R/4}}(z+\varepsilon)^\theta\dd x\right)^{1/\theta}
 \le C\operatorname*{ess\,inf}_{B_{R/2}}(z+\varepsilon).
 \label{eq:kato-weak-harnack}
\end{equation}

We finally justify the nonlinear tests.  For fixed \(\varepsilon>0\)
and \(m>0\), put \(F_{\varepsilon,m}(t)=(t+\varepsilon)^{-m}\).
Choose \(\chi\in C_c^\infty(D)\), \(0\le\chi\le1\), which equals one
on a neighborhood of the cutoff support.  By hypothesis,
\(\chi z\in\mathcal F_K(D)\).
The function
\(G_{\varepsilon,m}:=F_{\varepsilon,m}-F_{\varepsilon,m}(0)\) is
Lipschitz on the bounded range under consideration and vanishes at
zero.  Extend it from \([0,\|z\|_\infty]\) to a globally Lipschitz
function \(\widehat G_{\varepsilon,m}\) on \(\R\), with
\(\widehat G_{\varepsilon,m}(0)=0\).  The contraction property of the
closed form therefore gives
\[
 \mathcal E_K(\widehat G_{\varepsilon,m}(\chi z),
              \widehat G_{\varepsilon,m}(\chi z))
 \le \operatorname{Lip}(\widehat G_{\varepsilon,m})^2
       \mathcal E_K(\chi z,\chi z).
\]
Multiplication by the cutoff preserves \(\mathcal F_K(D)\); its
commutator is controlled by \eqref{eq:kato-cutoff-kernel-bound}.  The
constant part \(F_{\varepsilon,m}(0)\) times the cutoff has finite form
energy for the same reason.  Hence
\(\eta v^{-q}\), \(\eta^2v^{-1-2q}\), and \(\eta^2/v\) belong to
\(\mathcal F_{K,c}(D)\) and are admissible in
\eqref{eq:regional-supersolution-with-potential}.  The form bound
\eqref{eq:kato-infinitesimal-form-bound} makes the potential term
continuous in the form norm.  Thus the scalar inequalities above pass
to the closed-form tests directly, and every resulting estimate is
uniform as \(\varepsilon\downarrow0\).

If \(z(x_0)=0\), take the preceding balls centered at \(x_0\).
For every \(\delta>0\), continuity makes
\(\{z<\delta\}\cap B_{R/2}(x_0)\) a set of positive measure.
Consequently, \(z\ge0\) gives
\[
 \operatorname*{ess\,inf}_{B_{R/2}(x_0)}(z+\varepsilon)=\varepsilon.
\]
Letting \(\varepsilon\downarrow0\) in
\eqref{eq:kato-weak-harnack} shows that \(z=0\) in a neighborhood of
\(x_0\).  Thus the zero set is relatively open.  It is relatively
closed by continuity; connectedness then gives \(z\equiv0\) in \(D\).
\end{proof}

\begin{proposition}[Directional coarea and the full weak linearization]
\label{prop:coarea-weak-linearization}
Let $U$ be a globally $L$-Lipschitz entire solution of
\eqref{eq:equation}.  Suppose that, for some $q\ne0$ and
$e=q/|q|$,
\begin{equation}
 U\in C^{1,\alpha}(B_{2r_0}),
 \qquad
 \|\nabla U-q\|_{L^\infty(B_{2r_0})}\le\frac14|q|.
 \label{eq:strict-directional-monotonicity}
\end{equation}
In particular, $\partial_eU\ge\lambda_0:=3|q|/4$ in $B_{2r_0}$.
Set
\[
 \tau:=sp-p+2,
 \qquad
 J_U(x,y):=(p-1)
 \frac{|U(x)-U(y)|^{p-2}}{|x-y|^{n+sp}}.
\]
On the set \(U(x)=U(y)\) we set \(J_U(x,y)=0\).  For every fixed
\(y\), this set has zero \(x\)-measure locally in \(B_{2r_0}\), by
strict monotonicity on the \(e\)-lines, so the convention does not
change any of the forms below.
Then $J_U$ is locally a well-defined symmetric Dirichlet kernel.  If
$B_r(x_0)\subset B_{r_0/2}$ and
\[
 A_U(x):=\int_{\R^n\setminus B_{r_0}}J_U(x,y)\dd y,
\]
then
\begin{equation}
 \int_{B_r(x_0)}A_U(x)\dd x\le Cr^{n+p-2}.
 \label{eq:crossed-morrey-bound}
\end{equation}
More generally, for every \(D_0\Subset B_{r_0}\) the same estimate
holds for balls \(B_r(x_0)\subset D_0\), with a constant which may also
depend on \(\operatorname{dist}(D_0,\partial B_{r_0})\).
In particular, $A_U$ is a Morrey--Kato absorption for the regional order
$\tau$, since
\begin{equation}
 \tau-(2-p)=sp>0.
 \label{eq:strict-kato-gap}
\end{equation}

Moreover, the a.e. directional derivative \(w=\partial_eU\) belongs to
\(\mathcal F_{J_U,\mathrm{loc}}(B_{r_0})\) and satisfies
\begin{equation}
 \iint_{\R^n\times\R^n}
 J_U(x,y)\delta w(x,y)\delta\varphi(x,y)\dd x\!\dd y=0
\label{eq:full-weak-linearization}
\end{equation}
for every \(\varphi\in C_c^\infty(B_{r_0})\).  The identity extends to
every \(\varphi\in\mathcal F_{J_U,c}(B_{r_0})\).
\end{proposition}

\begin{proof}
We first record the quantitative form of the directional coarea
estimate used below.  For \(t\in[0,1]\) and \(0\le h<r_0/4\), put
\[
 W_{t,h}(x):=(1-t)U(x)+tU(x+he),
\]
with \(W_{t,0}=U\).  If \(K\Subset B_{r_0}\), then, after decreasing
the upper bound for \(h\) by a number depending only on \(K\),
\begin{equation}
 \sup_{\substack{0\le h<h_K\\t\in[0,1]\\a\in\R}}
 \int_E|W_{t,h}(x)-a|^{p-2}\dd x
 \le C_K|E|^{p-1}
 \qquad(E\subset K).
 \label{eq:quantitative-directional-coarea}
\end{equation}
Indeed, write \(x=x'+\ell e\), with \(x'\in e^\perp\).
On every line section, \(\partial_eW_{t,h}\ge\lambda_0\).  The
one-dimensional area formula, followed by decreasing rearrangement
about \(a\), gives
\[
 \int_{E_{x'}}|W_{t,h}(x'+\ell e)-a|^{p-2}\dd\ell
 \le C|E_{x'}|^{p-1}.
\]
Integration in \(x'\), followed by H\"older's inequality on the
projection of \(K\), proves
\eqref{eq:quantitative-directional-coarea}.

For \(E=B_r(x_0)\), the line calculation gives the sharper bound
\(Cr^{n-1}r^{p-1}=Cr^{n+p-2}\).  Apply it with \(h=0\) and
\(a=U(y)\), and then integrate in \(y\notin B_{r_0}\).  The two sets
are separated, while for large \(y\) the remaining kernel is bounded
by \(C(1+|y|)^{-n-sp}\), uniformly in \(x\in B_r(x_0)\).
This proves \eqref{eq:crossed-morrey-bound}.
Replacing \(B_{r_0/2}\) by any fixed \(D_0\Subset B_{r_0}\) in the
same argument proves the stated local version, since the two regions
remain positively separated.

We next justify differentiation of the equation.  For $0<h<r_0/4$ put
\begin{align*}
 U_h(x)&:=U(x+he),\qquad w_h:=\frac{U_h-U}{h},\\
 c_h(x,y)&:=(p-1)\int_0^1
 |(1-t)\delta U(x,y)+t\delta U_h(x,y)|^{p-2}\dd t.
\end{align*}
Here and below \(c_h=0\) when
\(\delta U(x,y)=\delta U_h(x,y)=0\).  The integral is finite otherwise,
even if the affine segment between the two increments crosses zero,
because \(p-2>-1\).
Subtraction of the two entire weak equations gives the exact identity
\begin{equation}
 \iint_{\R^n\times\R^n}
 \frac{c_h(x,y)\delta w_h(x,y)\delta\varphi(x,y)}
 {|x-y|^{n+sp}}\dd x\!\dd y=0.
 \label{eq:exact-secant-linear-equation}
\end{equation}
Estimate \eqref{eq:quantitative-directional-coarea} applies uniformly
to every convex combination \((1-t)U+tU_h\).  Together with Lemma
\ref{lem:regional-angular-compactness}, it gives, for each compact
\(K\Subset B_{r_0}\), a number \(h_K>0\) such that
\begin{equation}
 \sup_{0<h<h_K}
 \int_K\int_{\R^n}(1\wedge|x-y|^2)
 \frac{c_h(x,y)}{|x-y|^{n+sp}}\dd y\!\dd x<\infty.
 \label{eq:averaged-secant-levy-bound}
\end{equation}
after reducing the admissible upper bound for \(h\), if necessary.
For clarity, on the near-diagonal part the angular lemma gives
\[
 c_h(x,x-\rho\omega)=\rho^{p-2}b_h(x,\rho,\omega),
 \qquad
 \int_{\mathbb S^{n-1}}b_h(x,\rho,\omega)\dd\omega\le C,
\]
and the corresponding radial integral is
\[
 \int_0^1\rho^{p-2}\rho^2\rho^{-sp-1}\dd\rho
 =\int_0^1\rho^{1-\tau}\dd\rho<\infty.
\]
On separated sets,
\eqref{eq:quantitative-directional-coarea} gives uniform
integrability.  More precisely, if \(K\Subset B_{r_0}\),
\(0<\varepsilon<R<\infty\), and
\[
 G\subset\{(x,y):x\in K,\ \varepsilon\le|x-y|\le R\},
\]
then application of that estimate to each \(x\)-fiber, followed by
H\"older's inequality in \(y\), gives
\begin{equation}
 \sup_{0<h<h_K}\iint_G
 \frac{c_h(x,y)}{|x-y|^{n+sp}}\dd x\!\dd y
 \le C_{K,\varepsilon,R}|G|^{p-1}.
 \label{eq:separated-secant-uniform-integrability}
\end{equation}

Choose a cutoff \(\eta\in C_c^\infty(B_{r_0})\) and test
\eqref{eq:exact-secant-linear-equation} with $\eta^2w_h$.  Since
\(|w_h|\le L\), the identity
\[
 (a-b)(\eta_x^2a-\eta_y^2b)
 =|\eta_xa-\eta_yb|^2-ab|\eta_x-\eta_y|^2
\]
and \eqref{eq:averaged-secant-levy-bound} give
\begin{equation}
 \sup_{0<h<h_\eta}
 \iint_{\R^n\times\R^n}
 \frac{c_h(x,y)|\delta(\eta w_h)(x,y)|^2}
 {|x-y|^{n+sp}}\dd x\!\dd y<\infty.
 \label{eq:uniform-secant-caccioppoli}
\end{equation}
where \(h_\eta\) is the number attached above to a compact
neighborhood of \(\operatorname{supp}\eta\).  We clarify the
admissibility of this test without invoking density for the measurable
secant form.  For each fixed \(h>0\),
\(U,U_h\in W^{s,p}_{\rm loc}\), and
\[
 c_h|\delta w_h|^2
 =h^{-2}\bigl[\Jp(\delta U_h)-\Jp(\delta U)\bigr]
             (\delta U_h-\delta U)
 \le Ch^{-2}|\delta U_h-\delta U|^p.
\]
Moreover, \(\eta^2w_h=\eta^2(U_h-U)/h\in W^{s,p}_0(B_{r_0})\).
Approximate this function in \(W^{s,p}_0(B_{r_0})\) by smooth compactly
supported functions, use the same approximants separately in the weak
equations for \(U_h\) and \(U\), and then subtract.  The local
\(W^{s,p}\) bounds and the global Lipschitz tail, which is integrable
because \(sp>p-1\), justify the limit in the two original weak forms.
The secant identity then gives
\eqref{eq:exact-secant-linear-equation} with the test
\(\eta^2w_h\).  The restrictions
\(|x-y|>\varepsilon\) and \(c_h\le N\) are used only to estimate the
resulting integrals; they do not replace the kernel in the equation.
Letting first \(N\uparrow\infty\) and then
\(\varepsilon\downarrow0\), the averaged L\'evy bound and Fatou's
lemma give \eqref{eq:uniform-secant-caccioppoli}.

As $h\downarrow0$, $w_h\to w$ locally uniformly in $B_{r_0}$ and almost
everywhere in $\R^n$; the latter follows from Rademacher's theorem, or
equivalently from one-dimensional differentiation on almost every line
parallel to \(e\).  For almost every pair with one point in
$B_{r_0}$,
\[
 \frac{c_h(x,y)}{|x-y|^{n+sp}}\longrightarrow J_U(x,y).
\]
Indeed, the exceptional set \(U(x)=U(y)\) has zero measure in \(x\)
for each fixed \(y\), by strict monotonicity on the \(e\)-lines.
Fatou's lemma in \eqref{eq:uniform-secant-caccioppoli} first gives the
local \(J_U\)-energy of \(w\).

We verify at this point that the limiting energy belongs to the local
closed form domain, rather than only to the maximal finite-energy
class.  We first prove the required core property directly on
\(D_0:=B_{r_0}\).  Define
\[
 \Theta(x):=P_{e^\perp}x+U(x)e,
 \qquad x\in D_0.
\]
The strict bound \(\partial_eU\ge\lambda_0\), the local Lipschitz
bound, and integration on the \(e\)-lines show that \(\Theta\) is
bi-Lipschitz from \(D_0\) onto \(D^\sharp:=\Theta(D_0)\); moreover,
\(\det D\Theta=\partial_eU\) is bounded above and below.  Indeed,
\(\Theta\) preserves the transverse coordinates, while strict
monotonicity on each \(e\)-line controls the remaining coordinate of
\(\Theta^{-1}\).  More explicitly, after rotating so that \(e=e_n\)
and inserting the intermediate point \((x',y_n)\in B_{2r_0}\), one has
\[
 |x_n-y_n|\le\lambda_0^{-1}
 \bigl(|U(x)-U(y)|+L|x'-y'|\bigr),
\]
which gives \(|x-y|\le C|\Theta(x)-\Theta(y)|\).

Let \(f\in L^2(D_0)\) have compact support in \(D_0\) and finite
regional \(J_U\)-energy, and put
\(\widetilde f=f\circ\Theta^{-1}\).  With
\(\xi=\Theta(x)\) and \(\zeta=\Theta(y)\), change of variables gives
\begin{equation}
 \iint_{D_0\times D_0}J_U(x,y)|\delta f(x,y)|^2\dd x\!\dd y
 \asymp
 \iint_{D^\sharp\times D^\sharp}
 \frac{|e\cdot(\xi-\zeta)|^{p-2}
       |\delta\widetilde f(\xi,\zeta)|^2}
 {|\xi-\zeta|^{n+sp}}\dd\xi\!\dd\zeta.
 \label{eq:noncritical-coordinate-form-equivalence}
\end{equation}
The kernel on the right is the translation-invariant kernel
\[
 L(H):=\frac{|e\cdot H|^{p-2}}{|H|^{n+sp}}
 =\frac{|e\cdot\omega|^{p-2}}{|H|^{n+\tau}},
 \qquad \omega=\frac H{|H|}.
\]
We denote its full-space quadratic form by
\[
 \mathcal E_L(g):=\iint_{\R^n\times\R^n}
 L(\xi-\zeta)|g(\xi)-g(\zeta)|^2\dd\xi\!\dd\zeta.
\]
Its angular density is integrable because \(p-2>-1\), and its Fourier
symbol satisfies
\[
 c_{n,\tau}\int_{\mathbb S^{n-1}}
 |\vartheta\cdot\omega|^\tau|e\cdot\omega|^{p-2}\dd\omega
 \asymp |\vartheta|^\tau.
\]
Thus its full-space form domain is \(H^{\tau/2}(\R^n)\).

Extend \(\widetilde f\) by zero outside \(D^\sharp\).  Since its
support has positive distance from \(\partial D^\sharp\),
\eqref{eq:noncritical-coordinate-form-equivalence} and direct radial
integration show that its full-space \(L\)-energy is finite.  If
\(\rho_\varepsilon\) is a standard mollifier, translation invariance
and Jensen's inequality give
\[
 \mathcal E_L(\rho_\varepsilon*\widetilde f-\widetilde f)
 \le\int_{\R^n}\rho_\varepsilon(h)
       \mathcal E_L(\widetilde f(\,\cdot+h)-\widetilde f)\dd h
 \longrightarrow0.
\]
For completeness, the last convergence follows by writing the energy
as the integral in \(H\) of the squared \(L^2\)-norm of the
corresponding increment.  Translation continuity in \(L^2\) gives
pointwise convergence in \(H\), while
\(4L(H)\|\widetilde f(\,\cdot+H)-\widetilde f\|_2^2\) is an integrable
majorant.  For small \(\varepsilon\), the convolution remains compactly
supported in \(D^\sharp\).

Pulling it back by \(\Theta\) gives a compactly supported
\(C^{1,\alpha}\) function in \(D_0\), converging to \(f\) in
\(L^2\) and in the regional \(J_U\)-energy by
\eqref{eq:noncritical-coordinate-form-equivalence}.  Mollifying once
more in the original variables produces functions in
\(C_c^\infty(D_0)\).  For this second mollification the differences
converge in \(C^1\); Lemma~\ref{lem:regional-angular-compactness} and
\(\int_0^1r^{1-\tau}\dd r<\infty\) then give convergence in the
\(J_U\)-energy.  If \(f\ge0\), both mollifications preserve
nonnegativity.  We have
therefore proved that every compactly supported finite-energy function
in \(D_0\) belongs to \(\mathcal F_{J_U}(D_0)\).  The approximants may
be chosen with supports in one fixed compact subset of \(D_0\).

Fatou's lemma applied to \eqref{eq:uniform-secant-caccioppoli} gives
finite \(J_U\)-energy for \(\eta w\), not merely for \(w\) on pairs
inside the cutoff set.  Applying the preceding core property to
\(f=\eta w\) proves
\(\eta w\in\mathcal F_{J_U}(B_{r_0})\).  Thus
\(w\in\mathcal F_{J_U,\mathrm{loc}}(B_{r_0})\).

To pass to the equation, put
\(K=\operatorname{supp}\varphi\) and choose the preceding cutoff
\(\eta\) to be one in a fixed neighborhood of \(K\).  We spell out the
three truncations used in \eqref{eq:exact-secant-linear-equation}.

First, on
\[
 \{(x,y):x\in K,
       \varepsilon\le|x-y|\le R\},
\]
the uniform integrability following
\eqref{eq:quantitative-directional-coarea}, the almost everywhere
convergence of the kernels, and \(|w_h|\le L\) permit the use of
Vitali's theorem.  Hence the integral converges to the corresponding
integral with kernel \(J_U\).

Second, if \(|y|>R\), then \(\delta\varphi(x,y)=\varphi(x)\) and
\(|\delta w_h(x,y)|\le2L\).  Directional coarea in the \(x\)-variable
gives
\begin{equation*}
 \sup_{\substack{0<h<h_K\\t\in[0,1]}}\int_K
 |(1-t)\delta U+t\delta U_h|^{p-2}\dd x\le C
\end{equation*}
for each \(y\).  After multiplication by
\(|x-y|^{-n-sp}\) and integration in \(y\), the contribution is
\(O(R^{-sp})\), uniformly in \(h\).

Third, choose \(\varepsilon\) smaller than the distance from \(K\) to
\(\{\eta\ne1\}\).  Then every pair with \(|x-y|<\varepsilon\) which
meets \(K\) lies in \(\{\eta=1\}\).  Cauchy--Schwarz,
\eqref{eq:uniform-secant-caccioppoli}, and the angular upper bound
therefore give
\[
 C\left(\int_0^\varepsilon r^{1-\tau}\dd r\right)^{1/2}
 =C\varepsilon^{(2-\tau)/2}\longrightarrow0.
\]
The estimate is uniform in \(h\).  Letting successively
\(h\downarrow0\), \(R\uparrow\infty\), and
\(\varepsilon\downarrow0\) permits passage to the limit in
\eqref{eq:exact-secant-linear-equation}.  This proves
\eqref{eq:full-weak-linearization}.

It remains to justify the asserted extension of the identity.  Let
\(\varphi\in\mathcal F_{J_U,c}(B_{r_0})\).  Multiplying a form-norm
approximating sequence by a fixed cutoff which equals one near
\(\operatorname{supp}\varphi\), we may choose
\(\varphi_j\in C_c^\infty(B_{r_0})\) with supports in a common compact
set and
\[
 \|\varphi_j-\varphi\|_2^2
 +\iint_{B_{r_0}\times B_{r_0}}J_U
       |\delta(\varphi_j-\varphi)|^2\dd x\!\dd y\longrightarrow0.
\]
Choose \(K\Subset K_1\Subset B_{r_0}\) so that all the supports lie in
\(K\).  On \(K_1\times K_1\), the regional part of
\eqref{eq:full-weak-linearization} passes to the limit by
Cauchy--Schwarz and the local form membership of \(w\).  Pairs meeting
\(K\) and \(B_{r_0}\setminus K_1\) are separated; directional coarea
gives the same local Morrey bound for their coefficient, and
\eqref{eq:kato-infinitesimal-form-bound} makes their contribution tend
to zero.  For the crossed part, write
\(A_U(x)=\int_{\R^n\setminus B_{r_0}}
J_U(x,y)\dd y\).  Since \(|w|\le L\),
\begin{align*}
 &\left|\int_{B_{r_0}}\int_{\R^n\setminus B_{r_0}}
 J_U(x,y)\delta w(x,y)(\varphi_j-\varphi)(x)\dd y\!\dd x\right|\\
 &\qquad\le2L
 \left(\int_KA_U\dd x\right)^{1/2}
 \left(\int_KA_U|\varphi_j-\varphi|^2\dd x\right)^{1/2}.
\end{align*}
The same
directional-coarea argument used for \eqref{eq:crossed-morrey-bound},
now with constants depending also on
\(\operatorname{dist}(K,\partial B_{r_0})\), gives the required local
Morrey bound for \(A_U\) on a neighborhood of \(K\).  The last factor
therefore tends to zero by the density extension of
\eqref{eq:kato-infinitesimal-form-bound}, applied to the kernel
\(J_U\).  Hence
\eqref{eq:full-weak-linearization} holds for every compactly supported
form-domain test.
\end{proof}

At the contact point, local $C^{1,\alpha}$ regularity gives
\begin{equation}
 0\le Z_h(0)=L_Vh-U_*(he)+U_*(0)\le C h^{1+\alpha}.
 \label{eq:ordered-gap-contact-rate}
\end{equation}
Nevertheless, division of
\eqref{eq:regional-sublinear-absorption} by $h$ produces the coefficient
$h^{p-2}$ in front of $(Z_h/h)^{p-1}$.  It diverges as $h\downarrow0$.
Thus ordinary first difference quotients lose exactly the fixed gain
provided by \eqref{eq:ordered-gap-contact-rate}; this is the precise
singular crossed-tail endpoint for a pointwise tail estimate.  Proposition
\ref{prop:coarea-weak-linearization} avoids that loss by integrating in
the strictly monotone direction before passing to the quotient.

\subsection{Completion of the Liouville theorem}
\label{subsec:liouville-completion}

\begin{corollary}[Rigidity at a maximal gradient point]
\label{cor:gradient-contact-bernstein}
Let $U_*$ be the nonaffine profile in Proposition
\ref{prop:no-linear-record-neck}.  Then such a profile cannot exist.
\end{corollary}

\begin{proof}
Put $e=q/L_V$ and $w=\partial_eU_*$.  The global Lipschitz bound gives
$w\le L_V$ almost everywhere, while
\eqref{eq:attained-maximal-gradient} gives $w(0)=L_V$.  In the ball
selected after Lemma~\ref{lem:regional-angular-compactness}, condition
\eqref{eq:strict-directional-monotonicity} holds with, for instance,
$\lambda_0=L_V/2$.  Proposition
\ref{prop:coarea-weak-linearization} applies.

Set $z=L_V-w\ge0$.  It is continuous in $B_{r_0}$ and $z(0)=0$.
For every \(\eta\in C_c^\infty(B_{r_0})\), the cutoff estimate and
Proposition~\ref{prop:coarea-weak-linearization} give
\[
 \eta z=L_V\eta-\eta w\in\mathcal F_{J_{U_*}}(B_{r_0}).
\]
Thus \(z\in\mathcal F_{J_{U_*},\mathrm{loc}}(B_{r_0})\); in
particular, its local energy is finite.  Split
\eqref{eq:full-weak-linearization} into pairs in
$B_{r_0}\times B_{r_0}$ and crossed pairs.  For every nonnegative test
\(\varphi\in\mathcal F_{J_{U_*},c}(B_{r_0})\) supported in
$B_{r_0/2}$, the form-domain extension in Proposition
\ref{prop:coarea-weak-linearization} is applicable, and one has
\(\delta z(x,y)=z(x)-z(y)\le z(x)\) whenever
\(x\in B_{r_0/2}\) and \(y\notin B_{r_0}\).  Symmetry of the full form
and the inequality \(z(y)\ge0\) therefore give
\[
 \iint_{B_{r_0}\times B_{r_0}}
 J_{U_*}(x,y)\delta z\,\delta\varphi\dd x\!\dd y
 \ge-2\int_{B_{r_0/2}}A_{U_*}z\varphi\dd x.
\]
Lemma~\ref{lem:regional-angular-compactness} gives
\eqref{eq:spectral-regional-bounds} with order $\tau=sp-p+2$, while
\eqref{eq:crossed-morrey-bound} gives
\eqref{eq:morrey-absorption} with $\kappa=2-p<\tau$.  Lemma
\ref{lem:kato-contact-propagation}, in its localized form with tests
supported in \(B_{r_0/2}\), therefore shows that $z=0$ in a ball about
the origin.  Thus $\partial_eU_*=L_V$ there.  Since
$|\nabla U_*|\le L_V$, all transverse derivatives vanish there, and
$U_*$ is affine with slope $L_Ve$ on that ball.  A short segment in the
ball parallel to $e$ is an extremal secant.  Theorem
\ref{thm:extremal-secant-rigidity} makes $U_*$ affine in $\R^n$, contrary
to Proposition~\ref{prop:no-linear-record-neck}.
\end{proof}

\begin{theorem}[Fixed-order Lipschitz Liouville theorem]
\label{thm:fixed-order-centered-closure}
If \(n\ge2\), $1<p<2$, $0<s<1$, and $sp>p-1$, every
globally Lipschitz entire weak solution of \eqref{eq:equation} is affine.
\end{theorem}

\begin{proof}
If a nonaffine solution existed, Proposition
\ref{prop:dilation-recurrent-cascade} and Corollary
\ref{cor:collapse-affine-radii} would produce the maximizing sequence used in
Proposition~\ref{prop:no-linear-record-neck}.  That proposition gives a
nonaffine entire profile attaining its global Lipschitz slope at the
origin.  Corollary~\ref{cor:gradient-contact-bernstein} excludes precisely
such a profile.
\end{proof}

\begin{remark}[Subcriticality of the crossed singularity]
\label{rem:gradient-contact-endpoint}
The pointwise estimate in Proposition
\ref{prop:crossed-sublinear-absorption} is sharp and produces the
divergent factor $h^{p-2}$.  The geometry of the contact direction
provides the additional information required to control this factor.  Since
$\partial_eU_*$ stays positive in the noncritical ball, every level set is
crossed transversally on an $e$-line.  The singular power $p-2>-1$ is
therefore integrable along that line, and
\eqref{eq:crossed-morrey-bound} assigns the exterior coefficient the
codimension $2-p$.  The regional derivative kernel has order
$\tau=sp-p+2$.  Identity \eqref{eq:strict-kato-gap} is exactly the
strict subcriticality required by the form absorption.  Thus the global
pointwise L\'evy bound used in a general strong-minimum principle such as
\cite{JarohsWethSMP} is unnecessary for this special kernel: regional
angular control and the directional Morrey estimate suffice.
\end{remark}

\begin{remark}[From rigidity to finite-scale decay]
Theorem~\ref{thm:extremal-secant-rigidity} treats exact macroscopic
saturation.  The centered dilation-hull argument reduces every
putative nonaffine entire solution to a recurrent profile whose gradient
attains the global Lipschitz amplitude.  Proposition
\ref{prop:coarea-weak-linearization} and Corollary
\ref{cor:gradient-contact-bernstein} exclude this profile directly.
Proposition
\ref{prop:affine-blowdown-rigidity} supplies the blow-down step in the
centered argument.  Hence the fixed-order Bernstein theorem
is unconditional throughout \(sp>p-1\).

The finite-scale conversion from this Liouville theorem is carried out
below without an affine-tube assumption.  Section
\ref{sec:common-recursion} places the bounded- and large-slope
alternatives in a single intrinsic iteration.
\end{remark}

\section{The intrinsic improvement-of-flatness iteration}
\label{sec:common-recursion}

We formulate the two analytic alternatives in a common normalization.  Put
\begin{equation*}
 \gamma:=sp-p+1>0.
\end{equation*}
At a center $x_0$, a radius $r$, an affine map
$\ell(x)=a+q\cdot(x-x_0)$, and an amplitude $H>0$, set
\begin{equation}
 v(X):=\frac{u(x_0+rX)-\ell(x_0+rX)}{rH},
 \qquad Q:=\frac qH.
 \label{eq:common-intrinsic-profile}
\end{equation}
Fix once and for all
\begin{equation*}
 \ell>\max\left\{p,\frac1{p-1}\right\}.
\end{equation*}
Let \(\mathcal A\) denote the finite-dimensional space of affine
functions on \(\R^n\).
For \(\beta\ge0\) and a fixed constant \(G\ge1\), the
common exterior state is
\begin{equation}
 \mathcal T_{\beta,G}(v,Q)
 :=\max\left\{G^{-1}\mathscr C_{\ell,\beta}(v),
   \mathcal T^{\rm pt}_{s,Q}(v)\right\},
 \label{eq:common-tail-state}
\end{equation}
where
\begin{align*}
 \mathbf c_\ell(v;j)
 &:=2^{-j}\inf_{m\in\mathcal A}
 \left(\fint_{B_{2^{j+1}}}|v-m|^\ell\dd X\right)^{1/\ell},\\
 \mathscr C_{\ell,\beta}(v)
 &:=\sup_{j\ge0}2^{-j\beta}\mathbf c_\ell(v;j),
\end{align*}
and
\begin{align*}
 \mathcal T^{\rm pt}_{s,Q}(v)
 &:=\max\left\{
 {\operatorname*{ess\,sup}}_{X\in B_{5/4}}
 \int_{\R^n\setminus B_{4/3}}
 \frac{(1+|Q|)^{2-p}|\mathcal D_Qv(X,Y)|}
 {|X-Y|^{n+sp}}\dd Y,\right.\\[-2mm]
 &\hspace{31mm}\left.
 {\operatorname*{ess\,sup}}_{X\in B_{3/2}}
 \int_{\R^n\setminus B_2}
 \frac{(1+|Q|)^{2-p}|\mathcal D_Qv(X,Y)|}
 {|X-Y|^{n+sp}}\dd Y\right\},\\
 \mathcal D_Qv(X,Y)
 &:=\Jp\bigl(Q\cdot(X-Y)+v(X)-v(Y)\bigr)
       -\Jp\bigl(Q\cdot(X-Y)\bigr).
\end{align*}
The number \(G\) is fixed before the bounded-slope compactness
argument and will not change during the iteration.  The first component
is unchanged by adding an affine function.  Thus
the bounded-slope and large-slope regimes see the same affine-invariant
excess and the same relative exterior flux.

Choose below a dyadic contraction \(\rho=2^{-N}\).  If
\begin{equation}
 w(X):=\frac{v(\rho X)-m(\rho X)}{\rho\eta},
 \qquad m\in\mathcal A,
 \label{eq:affine-invariant-rescaling}
\end{equation}
then, for \(j\ge N\), change of variables and invariance of
\(\mathcal A\) give the exact identity
\begin{equation}
 \mathbf c_\ell(w;j)=\eta^{-1}\mathbf c_\ell(v;j-N).
 \label{eq:exact-affine-excess-shift}
\end{equation}
Consequently, when \(\eta=\rho^\beta\),
\begin{equation}
 \sup_{j\ge N}2^{-j\beta}\mathbf c_\ell(w;j)
 =\mathscr C_{\ell,\beta}(v).
 \label{eq:exact-affine-weight-cancellation}
\end{equation}
If one affine map satisfies
\begin{equation}
 t^{-1}\|v-m\|_{L^\infty(B_{2t})}\le\varepsilon\eta
 \qquad(\rho\le t\le1/2),
 \label{eq:common-affine-inner-block}
\end{equation}
then the competitor zero in the definition of
\(\mathbf c_\ell(w;j)\), with \(t=\rho2^j\), also gives
\begin{equation}
 \sup_{0\le j<N}2^{-j\beta}\mathbf c_\ell(w;j)\le\varepsilon.
 \label{eq:common-affine-inner-levels}
\end{equation}
Thus the inner and outer levels are combined by a maximum, with no
multiplicative tail-renormalization constant.

We shall use the following physical intrinsic amplitude.  If
$\ell(x)=a+q\cdot(x-x_0)$, put
\begin{align*}
 \mathbf Z_r(u,\ell;x_0)
 &:={\fint}_{B_{2r}(x_0)}
 \left|\frac{u-\ell}{r}\right|^p\dd x,\\
 \mathbf E_{s,r}(u,\ell;x_0)
 &:=(1-s)r^{sp-p-n}
 \iint_{B_{2r}(x_0)\times B_{2r}(x_0)}
 \frac{D_\Phi\bigl(q\cdot(x-y);\delta(u-\ell)(x,y)\bigr)}
 {|x-y|^{n+sp}}\dd x\!\dd y.
\end{align*}
Then $\mathbf A_{s,r}(u,\ell;x_0)$ is the infimum of all $A>0$
such that
\begin{equation*}
 \mathbf Z_r(u,\ell;x_0)\le A^p,
 \qquad
 \mathbf E_{s,r}(u,\ell;x_0)
 \le(|q|+A)^{p-2}A^2.
\end{equation*}

We next prove the joint compactness statement needed by the
affine-invariant state.  For \(R\ge1\), put
\begin{equation}
 \mathscr E_\ell(U;R):=R^{-1}\inf_{m\in\mathcal A}
 \left(\fint_{B_{2R}}|U-m|^\ell\dd X\right)^{1/\ell}.
 \label{eq:continuous-affine-excess-main}
\end{equation}

\begin{lemma}[Coherent affine fits]
\label{lem:coherent-affine-fits-main}
Suppose that \(\sup_{R\ge1}\mathscr E_\ell(U;R)\le G\).  For dyadic
\(R\), let \(m_R=c_R+b_R\cdot X\) minimize
\eqref{eq:continuous-affine-excess-main}.  Then
\begin{align}
 |b_{2R}-b_R|+R^{-1}|c_{2R}-c_R|&\le CG,
 \label{eq:coherent-affine-coefficients-main}\\
 \left(\fint_{B_{2^{h+1}R}}|U-m_R|^\ell\dd X\right)^{1/\ell}
 &\le CG(h+1)2^hR\qquad(h\ge0).
 \label{eq:coherent-log-shells-main}
\end{align}
\end{lemma}

\begin{proof}
On \(B_{2R}\), the triangle inequality bounds the \(L^\ell\) norm of
\(m_{2R}-m_R\) by \(CGR\).  Finite-dimensional norm equivalence for
affine maps, after scaling by \(R\), gives
\eqref{eq:coherent-affine-coefficients-main}.  Sum that estimate from
\(R\) to \(2^hR\), estimate every affine difference on the largest
ball, and add the minimizing error at scale \(2^hR\).  This proves
\eqref{eq:coherent-log-shells-main}.
\end{proof}

\begin{lemma}[Sublinear normalization at a best-affine scale]
\label{lem:sublinear-best-affine-normalization-main}
Let \(U\) be an entire weak solution satisfying
\(\sup_{R\ge1}\mathscr E_\ell(U;R)\le G\), and set
\[
 v_R(X)=\frac{U(RX)-m_R(RX)}R,\qquad Q_R=b_R.
\]
With \(\theta=n/(n+p)<1\), there is
\begin{equation}
 1\le A_R\le C(1+|b_R|^\theta)
 \label{eq:sublinear-best-affine-amplitude-main}
\end{equation}
such that \(\widehat v_R=v_R/A_R\), with tilt
\(\widehat Q_R=Q_R/A_R\), satisfies
\begin{equation}
 \fint_{B_2}|\widehat v_R|^p\dd X
 +\mathcal E_{s,\widehat Q_R}(\widehat v_R;B_2)
 +\mathcal T^{\rm pt}_{s,\widehat Q_R}(\widehat v_R)\le1.
 \label{eq:sublinear-best-affine-unit-data-main}
\end{equation}
In particular, \(A_R/|b_R|\to0\) whenever \(|b_R|\to\infty\).
\end{lemma}

\begin{proof}
The minimizing property and \(\ell>p\) give a uniform local \(L^p\)
bound for \(v_R\).  Apply local boundedness and the Lipschitz estimate
\cite[Theorem~2.1]{BiswasTopp} to the constant-recentered total solution
\[
 W_R(X)=\frac{U(RX)-c_R}{R}=b_R\cdot X+v_R(X).
\]
Lemma~\ref{lem:coherent-affine-fits-main} controls its local \(L^p\)
norm and scalar tail by \(C(1+|b_R|)\).  The weak solution is locally
continuous and agrees with the viscosity solution used in
\cite[Proposition~1.5 and Theorem~2.1]{BiswasTopp}.  Applying that
estimate on a slightly larger ball and then using a finite covering
gives
\[
 \operatorname{Lip}(v_R;B_3)\le C(1+|b_R|).
\]
Lipschitz--\(L^p\) interpolation yields
\begin{equation}
 \|v_R\|_{L^\infty(B_3)}
 \le C(1+|b_R|^{n/(n+p)}).
 \label{eq:sublinear-best-affine-sup-main}
\end{equation}
Put \(S_R=C(1+|b_R|^\theta)\) and take \(A_R=C_0S_R\), where
\(C_0\) will be fixed at the end.  It remains to check the relative
tail.  If \(|\widehat Q_R|\le2\), use the global
\((p-1)\)-H\"older estimate for \(\Jp\).  If
\(\widehat Q_R=q e\), \(q>2\), the linearized scalar estimate gives
\begin{equation}
 (1+q)^{2-p}|\mathcal D_{\widehat Q_R}\widehat v_R(X,Y)|
 \le C|X-Y|^{p-2}|e\cdot\omega|^{p-2}
       |\delta\widehat v_R(X,Y)|,
 \label{eq:best-affine-angular-flux-main}
\end{equation}
where \(\omega=(X-Y)/|X-Y|\).  Since
\((p-2)\ell'>-1\), H\"older's inequality on the sphere and
\eqref{eq:coherent-log-shells-main} bound the \(h\)-th exterior shell,
after division by \(A_R\), by
\begin{equation}
 C A_R^{1-p}(h+1)^{p-1}2^{-h\gamma}
 +C A_R^{-1}(h+1)2^{-h\gamma}.
 \label{eq:best-affine-shell-series-main}
\end{equation}
The first term is the global \((p-1)\)-H\"older branch and the second
is the angular linearized branch.
The series is summable because \(\gamma>0\); the local value in the
pointwise difference is controlled by
\eqref{eq:sublinear-best-affine-sup-main}.  The same shell estimate
applies to both separated geometries in the definition of
\(\mathcal T^{\rm pt}\).

For completeness, choose \(\zeta\in C_c^\infty(B_3)\), with
\(\zeta=1\) on \(B_2\), and test the equation for \(v_R\) with
\(\zeta^2v_R\).  We give the two estimates entering this test.  If
\(|b_R|\le4S_R\), the global \(p\)-growth inequality, the bound
\eqref{eq:sublinear-best-affine-sup-main}, and the shell estimate
\eqref{eq:coherent-log-shells-main} control both the cutoff term and
the exterior interaction by
\[
 CS_R^p\le C(|b_R|+S_R)^{p-2}S_R^2.
\]
If \(|b_R|>4S_R\), use the angular modular bound for the local cutoff.
For the exterior interaction, apply the linearized angular flux bound
on every dyadic shell and then H\"older's inequality on the sphere.
The condition \((p-2)\ell'>-1\) and
\eqref{eq:coherent-log-shells-main} give
\[
 C|b_R|^{p-2}S_R^2
 \sum_{h\ge0}(h+1)2^{-h\gamma}
 \le C(|b_R|+S_R)^{p-2}S_R^2.
\]
Here the fixed factor \(G\) in the coherent shell estimate has been
absorbed in the constant defining \(S_R\).  Shifted monotonicity
therefore gives
\begin{equation}
 \begin{split}
 &(1-s)\iint_{B_2^2}
 \frac{D_\Phi(b_R\cdot(X-Y);\delta v_R(X,Y))}
 {|X-Y|^{n+sp}}\dd X\!\dd Y\\
 &\hspace{25mm}\le
 C(|b_R|+S_R)^{p-2}S_R^2.
 \end{split}
 \label{eq:best-affine-shifted-caccioppoli}
\end{equation}
Homogeneity gives
\[
 D_\Phi\left(\frac{b_R}{A_R}\cdot Z;
             \frac{t}{A_R}\right)
 =A_R^{-p}D_\Phi(b_R\cdot Z;t),
\qquad
 \lambda_{\widehat Q_R}^{2-p}
 =A_R^{p-2}(A_R+|b_R|)^{2-p}.
\]
Equations \eqref{eq:best-affine-shell-series-main} and
\eqref{eq:best-affine-shifted-caccioppoli} show, after increasing
\(C_0\), that the three quantities in
\eqref{eq:sublinear-best-affine-unit-data-main} have sum at most one.
The last assertion follows from
\(\theta<1\).
\end{proof}

\begin{lemma}[Large-slope state reproduction]
\label{lem:large-slope-escape-state-main}
Fix
\[
 0<\alpha<\bar\alpha<\gamma
\]
and set \(L_{\rm esc}=2\).  There are a dyadic
\(\rho_{\rm esc}\in(0,1/16)\), a threshold \(Q_{\rm esc}\), and
\(C_{\rm esc}<\infty\) with the following property.  Suppose that \(v\)
solves the equation shifted by \(Q\), \(|Q|\ge Q_{\rm esc}\), and
\[
 \fint_{B_2}|v|^p\dd X\le1,\qquad
 \mathcal E_{s,Q}(v;B_2)\le1,\qquad
 \mathcal T^{\rm pt}_{s,Q}(v)\le L_{\rm esc}.
\]
Then there is an affine map \(m=c+b\cdot X\), with
\(|c|+|b|\le C_{\rm esc}\), such that, with
\[
 \eta_{\rm esc}:=\rho_{\rm esc}^{\alpha},\qquad
 w(X):=\frac{v(\rho_{\rm esc}X)-m(\rho_{\rm esc}X)}
 {\rho_{\rm esc}\eta_{\rm esc}},\qquad
 P:=\frac{Q+b}{\eta_{\rm esc}},
\]
one has
\begin{equation}
 \fint_{B_2}|w|^p\dd X\le1,\qquad
 \mathcal E_{s,P}(w;B_2)\le1,\qquad
 \mathcal T^{\rm pt}_{s,P}(w)\le L_{\rm esc},\qquad
 |P|\ge Q_{\rm esc}.
 \label{eq:large-slope-escape-state-reproduction}
\end{equation}
\end{lemma}

\begin{proof}
Apply Corollary~\ref{cor:prescribed-large-slope-block-main} with the
fixed budget \(L=L_{\rm esc}\) and the two exponents
\(\alpha,\bar\alpha\).  Its stable estimate first fixes a constant
\(C_{\rm L}\), independently of the radius.  Let \(C_0\) dominate the
constants in the scalar and angular flux estimates, including the
coefficient bound for the affine map furnished by that corollary.
Fix \(C_{\rm esc}\) to dominate this affine coefficient bound.
Choose a dyadic \(\rho_{\rm esc}\) so small that
\begin{align}
& C_{\rm L}\rho_{\rm esc}^{\bar\alpha-\alpha}\le\frac12,
 \label{eq:escape-large-slope-endpoint}\\
& C_0\bigl[\rho_{\rm esc}^{\bar\alpha-\alpha}
 +\rho_{\rm esc}^{\gamma-\alpha}\bigr]+C_0L_{\rm esc}\rho_{\rm esc}^{\gamma-\alpha}
 \le L_{\rm esc}.
 \label{eq:escape-large-slope-tail-budget}
\end{align}
All exponents are positive.  With this radius fixed, the prescribed
scale corollary gives a threshold \(Q_{\rm L}\) and an affine map \(m\)
such that
\[
 \mathfrak a_{s,Q}(v,m;\rho_{\rm esc})
 \le C_{\rm L}\rho_{\rm esc}^{\bar\alpha}
 \le\frac12\eta_{\rm esc},
\qquad
 t^{-1}\|v-m\|_{L^\infty(B_{2t})}
 \le C_{\rm L}t^{\bar\alpha}
\]
for \(\rho_{\rm esc}\le t\le1/2\).  By the strict slack in the
preceding estimate, choose an admissible \(a<\eta_{\rm esc}\) in the
definition of \(\mathfrak a_{s,Q}(v,m;\rho_{\rm esc})\), and put
\(q=|Q+b|\).  The zero-order term scales as
\((a/\eta_{\rm esc})^p\).  Exact
homogeneity of the energy gives
\begin{equation*}
 \mathcal E_{s,P}(w;B_2)
 \le
 \left(\frac{q+\eta_{\rm esc}}{q+a}\right)^{2-p}
 \left(\frac a{\eta_{\rm esc}}\right)^2
 \le\left(\frac a{\eta_{\rm esc}}\right)^p\le1.
\end{equation*}
Here the middle inequality follows from
\((q+\eta_{\rm esc})/(q+a)\le\eta_{\rm esc}/a\).  This proves the
first two inequalities in
\eqref{eq:large-slope-escape-state-reproduction} without losing a
structural constant.

We record the pointwise-tail calculation because it fixes the order of
the parameters.  Put \(h=v-m\), \(Z=\rho_{\rm esc}X\), and
\(W=\rho_{\rm esc}Y\).  The change of variables gives
\begin{equation}
 \begin{split}
 &\lambda_P^{2-p}
 \frac{|\mathcal D_Pw(X,Y)|}{|X-Y|^{n+sp}}\dd Y\\
 &\qquad=
 \rho_{\rm esc}^{\gamma}\eta_{\rm esc}^{1-p}
 \lambda_P^{2-p}
 \frac{|\mathcal D_{Q+b}h(Z,W)|}
 {|Z-W|^{n+sp}}\dd W .
 \end{split}
 \label{eq:escape-pointwise-flux-scaling}
\end{equation}
We shall choose \(Q_{\rm esc}\) so large that both \(Q\) and
\(Q+b\) lie in the large-slope range.  The comparison
\begin{equation*}
 \frac{\lambda_P}{\lambda_{Q+b}}\le C\eta_{\rm esc}^{-1},
 \qquad
 \frac{\lambda_P}{\lambda_Q}\le C\eta_{\rm esc}^{-1}
\end{equation*}
shows why the factor \(\lambda_P^{2-p}\) in
\eqref{eq:escape-pointwise-flux-scaling} cannot be discarded.  On the
intervening annuli
\(\rho_{\rm esc}\lesssim|W|\lesssim1\), the angular linearized bound
gives
\begin{equation*}
 C\rho_{\rm esc}^{\gamma}\eta_{\rm esc}^{-1}
 \sum_{\rho_{\rm esc}\le t\le1}t^{\bar\alpha-\gamma}.
\end{equation*}
The finitely many shells above \(1/2\) are controlled by the local
bound and the coefficient bound for \(m\); for the second tail
geometry the old first tail component is used once
\(|W|\ge4/3\).  Thus the sum is bounded by the two terms in the
first line of \eqref{eq:escape-large-slope-tail-budget}.

On the remaining exterior region, decompose
\[
 \mathcal D_{Q+b}h
 =\mathcal D_Qv+
 \Jp(Q\cdot(Z-W))-\Jp((Q+b)\cdot(Z-W)).
\]
The old relative flux has the exact multiplier
\begin{equation*}
 \rho_{\rm esc}^{\gamma}\eta_{\rm esc}^{1-p}
 \left(\frac{\lambda_P}{\lambda_Q}\right)^{2-p}
 \le C\rho_{\rm esc}^{\gamma}\eta_{\rm esc}^{-1}.
\end{equation*}
The angular estimate applied to the affine-base correction gives the
same upper bound; more explicitly, its spherical and radial integrals
are bounded by
\[
 C|b|\rho_{\rm esc}^{\gamma}\eta_{\rm esc}^{-1}
 \int_1^\infty t^{-\gamma-1}\dd t
 \le C\rho_{\rm esc}^{\gamma-\alpha}.
\]
Consequently,
\begin{align}
 \mathcal T^{\rm pt}_{s,P}(w)
 \le C_0\bigl[
 \rho_{\rm esc}^{\bar\alpha-\alpha}
 +\rho_{\rm esc}^{\gamma-\alpha}\bigr]
 +C_0\rho_{\rm esc}^{\gamma-\alpha}
 \mathcal T^{\rm pt}_{s,Q}(v).
 \label{eq:escape-pointwise-tail-recurrence}
\end{align}
The choice \eqref{eq:escape-large-slope-tail-budget} closes the same
budget \(L_{\rm esc}\); in particular, there is no dependence
\(L\mapsto\rho(L)\mapsto L_{\rm new}\).

Finally choose
\[
 Q_{\rm esc}\ge
 \max\left\{Q_{\rm L}+1,
 2C_{\rm esc}+2,
 \frac{C_{\rm esc}}{1-\eta_{\rm esc}}\right\}.
\]
Then
\(|P|\ge\eta_{\rm esc}^{-1}(|Q|-|b|)\ge Q_{\rm esc}\).
This proves \eqref{eq:large-slope-escape-state-reproduction}.
\end{proof}

\begin{proposition}[Affine-Campanato Liouville theorem]
\label{prop:affine-campanato-liouville-main}
If \(U\) is an entire weak solution, \(U(0)=0\), and
\begin{equation*}
 \sup_{R\ge1}\mathscr E_\ell(U;R)<\infty,
\end{equation*}
then \(U\) is affine.
\end{proposition}

\begin{proof}
Let \(m_R=c_R+b_R\cdot X\) be the dyadic minimizing maps.  We claim
that \(\sup_R|b_R|<\infty\).  Otherwise choose \(R_j\to\infty\) with
\(|b_{R_j}|\to\infty\), and apply Lemma
\ref{lem:sublinear-best-affine-normalization-main}.  Put
\[
 \widetilde v_j(X)=
 \frac{U(R_jX)-m_{R_j}(R_jX)}{A_jR_j},
 \qquad \widetilde Q_j=\frac{b_{R_j}}{A_j},
 \qquad A_j=A_{R_j}.
\]
Then \(|\widetilde Q_j|\to\infty\).  Fix
\(0<\alpha<\bar\alpha<\gamma\), and fix
\(\rho_{\rm esc},Q_{\rm esc}\) by Lemma
\ref{lem:large-slope-escape-state-main}.  Discarding finitely many
indices, \(|\widetilde Q_j|\ge Q_{\rm esc}\).  Iterate that lemma while
\(R_j\rho_{\rm esc}^k\ge4\).  Writing
\(r_k=\rho_{\rm esc}^k\), \(H_k=r_k^\alpha\), one obtains accumulated
affine maps \(a_{j,k}\) such that
\begin{align}
 \left(\fint_{B_{2r_k}}|\widetilde v_j-a_{j,k}|^p\dd X\right)^{1/p}
 &\le Cr_kH_k,
 \label{eq:escaping-affine-error-main}\\
 |\nabla a_{j,k+1}-\nabla a_{j,k}|&\le CH_k.
 \label{eq:escaping-affine-increments-main}
\end{align}
At every step, \eqref{eq:large-slope-escape-state-reproduction}
preserves both the lower bound for the normalized tilt and the fixed
tail budget \(L_{\rm esc}\).  Thus the iteration is valid without a
parameter-dependence loop.

Choose \(k=k(j)\) so that
\(4\le R_jr_k<4\rho_{\rm esc}^{-1}\).  On this fixed physical ball,
\eqref{eq:escaping-affine-error-main} gives an affine approximation to
\(U\) whose slope is
\[
 b_{R_j}+A_j\nabla a_{j,k(j)}=b_{R_j}+O(A_j)
\]
and whose \(L^p\) error is at most \(CA_jR_j^{-\alpha}\).  By
\eqref{eq:coherent-affine-coefficients-main},
\(|b_R|\le C(1+\log R)\).  Hence
\(A_j\le C(1+(\log R_j)^\theta)\) by
\eqref{eq:sublinear-best-affine-amplitude-main}, so this error tends to
zero.  Norm
equivalence for affine maps on the fixed ball bounds the displayed
slope.  Since \(A_j/|b_{R_j}|\to0\), this is a contradiction.

Thus the slopes \(b_R\) are bounded.  Lemma
\ref{lem:coherent-affine-fits-main} now gives
\((\fint_{B_{2R}}|U|^\ell)^{1/\ell}\le CR\) and the corresponding
scale-invariant scalar-tail bound.  Local boundedness and
\cite[Proposition~1.5 and Theorem~2.1]{BiswasTopp}, applied to the
continuous weak solution \(U(R\,\cdot)/R\), show with a scale-independent
constant that \(U\) is globally Lipschitz.
Theorem~\ref{thm:fixed-order-centered-closure}
makes it affine.
\end{proof}

\begin{proposition}[Bounded-slope affine improvement]
\label{prop:bounded-slope-affine-invariant-block-main}
Fix \(Q_0,G_0<\infty\), \(0<\rho<1/16\), and
\(0<\alpha_1<\gamma\).  There are \(\beta_*>0\) and \(R_*\ge8\) such
that, for every \(0\le\beta\le\beta_*\) and every \(|Q|\le Q_0\), the
following holds.
If \(v\) solves the equation shifted by \(Q\) in \(B_{2R_*}\) and
\begin{equation}
 \mathfrak a_{s,Q}(v,0;1)\le1,
 \qquad\mathscr C_{\ell,\beta}(v)\le G_0,
 \qquad\mathcal T^{\rm pt}_{s,Q}(v)\le1,
 \label{eq:bounded-slope-affine-block-assumptions-main}
\end{equation}
then there is an affine map \(m\) satisfying
\begin{equation*}
 |m(0)|+|\nabla m|\le C_{\rm B},
\end{equation*}
where \(C_{\rm B}=C_{\rm B}(n,p,s,\alpha_1)\) is independent of
\(Q_0,G_0,\rho\), such that
\begin{equation}
\mathfrak a_{s,Q}(v,m;t)
 +t^{-1}\|v-m\|_{L^\infty(B_{2t})}\le t^{\alpha_1}
 \qquad(\rho\le t\le1/2).
 \label{eq:bounded-slope-affine-full-block-main}
\end{equation}
\end{proposition}

\begin{proof}
Let \(C_{\rm aff}\) be the structural norm-equivalence constant for
affine functions on \(B_1\).  Fix \(C_{\rm B}>4C_{\rm aff}\) before
selecting the counterexample sequence.

Suppose that the assertion is false.  Its precise negation says that,
for every \(\widehat\beta>0\) and every \(\widehat R\ge8\), there are
\(0\le\beta\le\widehat\beta\), \(|Q|\le Q_0\), and a solution on
\(B_{2\widehat R}\) satisfying
\eqref{eq:bounded-slope-affine-block-assumptions-main} but violating
\eqref{eq:bounded-slope-affine-full-block-main} for every affine map
\(m\) satisfying \(|m(0)|+|\nabla m|\le C_{\rm B}\).  Taking
\(\widehat\beta=1/j\) and
\(\widehat R=R_j\ge j\), choose corresponding counterexamples
\((v_j,Q_j,\beta_j)\).  Thus \(\beta_j\le1/j\), \(R_j\to\infty\),
and \(v_j\) solves the shifted equation in \(B_{2R_j}\).
After a subsequence, \(Q_j\to Q\).  Let
\(m_{j,h}=c_{j,h}+b_{j,h}\cdot X\) be a best affine fit on
\(B_{2^{h+1}}\).  Comparison on \(B_{2^{h+1}}\) and
finite-dimensional norm equivalence give
\begin{equation*}
 |b_{j,h+1}-b_{j,h}|+2^{-h}|c_{j,h+1}-c_{j,h}|
 \le C G_0 2^{h\beta_j}.
\end{equation*}
The unit-scale intrinsic bound and
\(\mathscr C_{\ell,\beta_j}(v_j)\le G_0\) give
\[
 |c_{j,0}|+|b_{j,0}|\le C(1+G_0)
\]
by finite-dimensional norm equivalence on \(B_2\).
Summing the preceding estimate yields
\begin{equation*}
 \left(\fint_{B_{2^{h+1}}}|v_j|^\ell\dd X\right)^{1/\ell}
 \le C(1+G_0)(h+1)2^{h(1+\beta_j)}.
\end{equation*}
In particular, \(v_j\) is uniformly bounded in \(L^\ell\) on every
fixed ball.  The same estimate gives the scalar tail required by local
boundedness.  Since the tilts \(Q_j\) are bounded, local boundedness and
the interior Lipschitz estimate give a locally uniform subsequential
limit \(v_j\to v_\infty\); no additional normalization by constants is
needed because the unit-scale \(L^p\) norm is bounded.

On the \(h\)-th exterior shell, the global \((p-1)\)-H\"older bound for
\(\Jp\) and the preceding estimate give
\begin{equation}
 C(h+1)^{p-1}2^{-h[sp-(p-1)(1+\beta_j)]}.
 \label{eq:bounded-slope-tail-tightness-main}
\end{equation}
Since \((1+\beta_j)(p-1)\le sp-\gamma/2\) for all large \(j\), the
series in \eqref{eq:bounded-slope-tail-tightness-main} is uniformly
summable.  We make the two truncations explicit.  On the near
diagonal, the local Lipschitz bound and the first-order difference of a
test function give
\[
 |\delta U_j|^{p-1}|\delta\varphi|
 |X-Y|^{-n-sp}
 \le C|X-Y|^{p-n-sp},
\]
whose radial integral is finite because \(p-sp>0\).  On the far
shells, \eqref{eq:bounded-slope-tail-tightness-main} is bounded by a
summable series with decay \(2^{-h\gamma/2}\).  The weak equation
therefore passes to the limit on every compact set, and
\(U_\infty=Q\cdot X+v_\infty\) is an entire weak solution.  Passing
first at dyadic radii and then
using radius comparability gives
\begin{equation*}
 \sup_{R\ge1}R^{-1}
 \inf_{m\in\mathcal A}
 \left(\fint_{B_{2R}}|U_\infty-m|^\ell\dd X\right)^{1/\ell}
 \le C(1+G_0).
\end{equation*}
After subtracting the constant \(U_\infty(0)\), Proposition
\ref{prop:affine-campanato-liouville-main} shows that \(U_\infty\), and
hence \(v_\infty\), is affine.  Write \(v_\infty=m_\infty\).  Its
coefficients are bounded by the unit-scale \(L^p\) bound and
finite-dimensional norm equivalence.  By the preceding choice,
\(|m_\infty(0)|+|\nabla m_\infty|\le C_{\rm B}/2\).

On each fixed ball the preceding Lipschitz bound supplies the integrable
majorant \(C|X-Y|^{p-n-sp}\) for every recentered Bregman integrand.
Local uniform convergence, dominated convergence, and compactness of the
interval of radii give
\begin{equation*}
 \sup_{\rho\le t\le1/2}
 \mathfrak a_{s,Q_j}(v_j,m_\infty;t)\longrightarrow0,
 \qquad
 \sup_{\rho\le t\le1/2}t^{-1}
 \|v_j-m_\infty\|_{L^\infty(B_{2t})}\longrightarrow0.
\end{equation*}
Since \(|m_\infty(0)|+|\nabla m_\infty|\le C_{\rm B}/2\), this
contradicts \eqref{eq:bounded-slope-affine-full-block-main}: for all
sufficiently large \(j\), the sum of the two displayed suprema is
smaller than \(\rho^{\alpha_1}\), and hence than \(t^{\alpha_1}\) for every
\(t\in[\rho,1/2]\).
\end{proof}


We fix the parameters in the following order.  First choose
\begin{equation}
 0<\alpha_1<\bar\alpha<\gamma.
 \label{eq:affine-route-intermediate-exponents-main}
\end{equation}
Let \(C_{\rm L}\) dominate both the constant and the affine coefficient
bound in Corollary~\ref{cor:prescribed-large-slope-block-main}, with
\(L=1\).  Replacing it by
\(\max\{C_{\rm L},2^{\alpha_1}C_{\rm B},1\}\), fix
\begin{equation}
 G_*:=2+C_{\rm L}.
 \label{eq:affine-route-campanato-budget-main}
\end{equation}
Next choose
\begin{equation}
 0<\bar\beta<\frac12\min\{\alpha_1,\gamma\}.
 \label{eq:affine-route-beta-bar-main}
\end{equation}
The scalar and angular flux inequalities, with constants enlarged by
the fixed number \(C_{\rm L}\), determine \(C_0\).  Choose a dyadic
\(\rho=2^{-N}<1/16\) so small that
\begin{align}
 C_{\rm L}\rho^{\bar\alpha-\alpha_1}&\le1,
 \label{eq:affine-route-large-endpoint-main}\\
 C_0\bigl[
 \rho^{(\alpha_1-\bar\beta)(p-1)}
 +\rho^{\gamma-(p-1)\bar\beta}
 +\rho^{\alpha_1-\bar\beta}
 +\rho^{\gamma-\bar\beta}\bigr]&\le\frac14.
 \label{eq:affine-route-tail-small-main}
\end{align}
All four exponents are positive.  Corollary
\ref{cor:prescribed-large-slope-block-main} now gives a threshold
\(Q_{\rm th}\) for this fixed \(\rho\).  Apply Proposition
\ref{prop:bounded-slope-affine-invariant-block-main} with
\[
 Q_0=Q_{\rm th},\qquad G_0=G_*,\qquad
 \rho\ \hbox{as above},\qquad
 \alpha_1\ \hbox{as in }\eqref{eq:affine-route-intermediate-exponents-main},
\]
and denote its compactness exponent by \(\beta_*\).  Finally set
\begin{equation}
 \beta:=\frac12\min\{\beta_*,\bar\beta\},
 \qquad\eta:=\rho^\beta.
 \label{eq:affine-route-final-beta-main}
\end{equation}
No radius, threshold, or structural constant is changed after
\(\beta_*\) has been obtained.

\begin{lemma}[Stability of the normalized state]
\label{lem:affine-invariant-state-reproduction-main}
Suppose that \(v\) solves the equation shifted by \(Q\) on the ball
required by the bounded-slope estimate and that
\begin{equation}
 \mathfrak a_{s,Q}(v,0;1)\le1,
 \qquad \mathcal T_{\beta,G_*}(v,Q)\le1.
 \label{eq:affine-route-unit-assumptions-main}
\end{equation}
If \(|Q|\le Q_{\rm th}\), let \(m=c+b\cdot X\) be the affine map in
Proposition~\ref{prop:bounded-slope-affine-invariant-block-main}; if
\(|Q|>Q_{\rm th}\), use the map in Corollary
\ref{cor:prescribed-large-slope-block-main}.  Put
\begin{equation}
 w(X)=\frac{v(\rho X)-m(\rho X)}{\rho\eta},
 \qquad P=\frac{Q+b}{\eta}.
 \label{eq:affine-route-renormalized-profile-main}
\end{equation}
Then
\begin{equation}
 \mathfrak a_{s,P}(w,0;1)\le1,\qquad
 \mathcal T_{\beta,G_*}(w,P)\le1.
 \label{eq:affine-route-unit-reproduction-main}
\end{equation}
\end{lemma}

\begin{proof}
The change of variables in the weak formulation and the homogeneity of
\(\Jp\) show that \(w\) solves the equation shifted by \(P\) on the
correspondingly rescaled domain.  In either regime the endpoint estimate is
\begin{equation}
 \mathfrak a_{s,Q}(v,m;\rho)\le\rho^{\alpha_1}<\eta.
 \label{eq:affine-route-common-endpoint-main}
\end{equation}
Indeed, this is the bounded-slope conclusion in the first regime and
follows from \eqref{eq:affine-route-large-endpoint-main} in the second.
Choose an admissible \(a<\eta\) in the definition of the left-hand
side of \eqref{eq:affine-route-common-endpoint-main}.  With
\(q=|Q+b|\), exact homogeneity gives
\[
 \fint_{B_2}|w|^p\dd X\le(a/\eta)^p,\qquad
 \mathcal E_{s,P}(w;B_2)
 \le
 \left(\frac{q+\eta}{q+a}\right)^{2-p}
 \left(\frac a\eta\right)^2
 \le\left(\frac a\eta\right)^p\le1.
\]
Thus the intrinsic estimate is preserved with constant one.

For \(0\le j<N\), put \(t=\rho2^j\).  The two finite-scale estimates give
\begin{equation}
 t^{-1}\|v-m\|_{L^\infty(B_{2t})}
 \le C_{\rm L}t^{\alpha_1}.
 \label{eq:affine-route-common-intermediate-main}
\end{equation}
Since \(\rho=2^{-N}\), one has \(\rho\le t\le1/2\); hence the displayed
range is exactly the range supplied by the bounded-slope proposition.
Here we used \(t^{\bar\alpha}\le t^{\alpha_1}\) in the large-slope
regime.  Consequently
\[
2^{-j\beta}\mathbf c_\ell(w;j)
 \le C_{\rm L}t^{\alpha_1-\beta}\le C_{\rm L}<G_*.
\]
For \(j\ge N\), the exact identity
\eqref{eq:exact-affine-weight-cancellation} gives the preceding state without
loss.  Hence \(\mathscr C_{\ell,\beta}(w)\le G_*\).

It remains to verify the pointwise flux.  After the change of variables,
split its domain into the intervening annuli
\(\rho\lesssim t\lesssim2\) and the remaining exterior region.  On the
intervening annuli put \(h=v-m\), \(Z=\rho X\), and \(W=\rho Y\).
The exact scaling identity is
\begin{equation}
 \begin{split}
 &\lambda_P^{2-p}
 \frac{|\mathcal D_Pw(X,Y)|}{|X-Y|^{n+sp}}\dd Y\\
 &\qquad=
 \rho^\gamma\eta^{1-p}\lambda_P^{2-p}
 \frac{|\mathcal D_{Q+b}h(Z,W)|}
 {|Z-W|^{n+sp}}\dd W .
 \end{split}
 \label{eq:affine-route-pointwise-flux-scaling}
\end{equation}
If \(|P|\le1\), the global \((p-1)\)-H\"older estimate gives the shell
bound
\begin{equation*}
 C\rho^\gamma\eta^{1-p}
 \int_{|W|\asymp t}
 \frac{|\delta h(Z,W)|^{p-1}}{|Z-W|^{n+sp}}\dd W.
\end{equation*}
If \(|P|>1\), put \(e=P/|P|\).  The angular linearized estimate gives
instead
\begin{equation*}
 C\rho^\gamma\eta^{-1}
 \int_{|W|\asymp t}
 \frac{|Z-W|^{p-2}|e\cdot\omega|^{p-2}
       |\delta h(Z,W)|}
 {|Z-W|^{n+sp}}\dd W,
 \qquad \omega=\frac{Z-W}{|Z-W|}.
\end{equation*}
By \eqref{eq:affine-route-common-intermediate-main},
\(|\delta h(Z,W)|\le Ct^{1+\alpha_1}\); angular H\"older and
Lemma~\ref{lem:large-slope-scalar-angular} therefore give
\begin{align}
 &C\rho^\gamma\eta^{1-p}
   \sum_{\rho\le t\le1/2}t^{\alpha_1(p-1)-\gamma}
 +C\rho^\gamma\eta^{-1}
   \sum_{\rho\le t\le1/2}t^{\alpha_1-\gamma}
 \notag\\
 &\quad\le C_0\bigl[
 \rho^{(\alpha_1-\beta)(p-1)}
 +\rho^{\gamma-(p-1)\beta}
 +\rho^{\alpha_1-\beta}
 +\rho^{\gamma-\beta}\bigr].
 \label{eq:affine-route-intervening-flux-main}
\end{align}
The two terms in each bracket are the lower- and upper-endpoint bounds
for the corresponding finite geometric sum; thus no sign of the sum
exponent is being assumed.  The finitely many remaining shells with
old radii between \(1/2\) and \(2\) are controlled in the scalar
branch by the local \(L^\ell\)-bound and in the angular branch by
H\"older's inequality, the same \(L^\ell\)-bound, and
\((p-2)\ell'>-1\).  The coefficient bound for \(m\) controls the
affine part.  More explicitly,
\(\mathscr C_{\ell,\beta}(v)\le G_*\) controls \(v\) modulo an affine
map on these fixed balls; the unit \(L^p(B_2)\)-bound and
finite-dimensional norm equivalence control that affine map, while the
estimate at \(t=\rho\) controls the inner endpoint \(Z=\rho X\).
Thus no estimate above the range \(t\le1/2\) is used.  These shells
contribute only
\(C\rho^{\gamma-(p-1)\beta}\) and
\(C\rho^{\gamma-\beta}\), which are already the second and fourth
upper-endpoint terms in \eqref{eq:affine-route-intervening-flux-main}.

On the remaining exterior region the H\"older and linearized alternatives
carry, respectively, the factors
\[
 \rho^\gamma\eta^{1-p}
 =\rho^{\gamma-(p-1)\beta},
 \qquad
 \rho^\gamma\eta^{-1}=\rho^{\gamma-\beta}.
\]
Here one uses the exact decomposition
\begin{equation*}
 \mathcal D_{Q+b}h(Z,W)
 =\mathcal D_Qv(Z,W)
 +\Jp(Q\cdot(Z-W))-\Jp((Q+b)\cdot(Z-W)).
\end{equation*}
For the first term, the exact coefficient relative to the old state is
\begin{equation*}
 \rho^\gamma\eta^{1-p}
 \left(\frac{\lambda_P}{\lambda_Q}\right)^{2-p}.
\end{equation*}
When \(|P|\le1\), this is bounded by
\(C\rho^\gamma\eta^{1-p}\); when \(|P|>1\), the comparison
\(\lambda_P\le C\eta^{-1}\lambda_Q\) bounds it by
\(C\rho^\gamma\eta^{-1}\).  Thus the two alternatives above are
precisely the two factors carried by the old pointwise state.

The change of affine base from \(Q\) to \(Q+b\) obeys the same
dichotomy.  More explicitly, homogeneity and the scalar/angular
estimates give, in spherical average,
\begin{equation*}
 \eta^{1-p}\lambda_P^{2-p}
 \int_{\mathbb S^{n-1}}
 \left|\Jp(Q\cdot(r\omega))
       -\Jp((Q+b)\cdot(r\omega))\right|\dd\omega
 \le Cr^{p-1}\bigl(\eta^{1-p}+\eta^{-1}\bigr).
\end{equation*}
This includes the transition region \(|Q+b|\asymp\eta\).  In both
finite-scale estimates
\(|b|+|c|\le C\) independently of \(Q,Q_{\rm th},\rho\); hence this
correction only changes \(C_0\).  Notice that, in the linearized case,
\[
 \mathcal D_{Q+b}h
 =\eta^{p-1}\mathcal D_P(h/\eta),
 \qquad Q+b=\eta P.
\]
Thus this branch also covers \(|Q+b|<1<|P|\).  Since
\(\beta\le\bar\beta\),
\eqref{eq:affine-route-tail-small-main}, with its constant enlarged at
the time it was fixed, makes the sum of
\eqref{eq:affine-route-intervening-flux-main} and the old far-field
contribution at most one.  This proves
\eqref{eq:affine-route-unit-reproduction-main}.
\end{proof}

\begin{proposition}[Intrinsic improvement of flatness]
\label{prop:unified-variable-block-recursion}
Let \(\rho,\beta,\eta,G_*\) and the two finite-scale estimates be fixed by
\eqref{eq:affine-route-intermediate-exponents-main}--\eqref{eq:affine-route-final-beta-main}.
Let \(r_0,H_0>0\), let \(\ell_0\) be affine, and define \(v_0,Q_0\) by
\eqref{eq:common-intrinsic-profile}.  Assume
\begin{equation*}
 \mathbf A_{s,r_0}(u,\ell_0;x_0)\le H_0,\qquad
 \mathcal T_{\beta,G_*}(v_0,Q_0)\le1.
\end{equation*}
Assume also that every rescaled profile is defined on the ball required
by the bounded-slope estimate and solves the corresponding shifted equation.
Then there are affine maps
\(\ell_k=a_k+q_k\cdot(x-x_0)\), radii \(r_k\), and amplitudes \(H_k\)
such that, with \(v_k,Q_k\) defined by
\eqref{eq:common-intrinsic-profile} using
\((r,H,\ell)=(r_k,H_k,\ell_k)\),
\begin{align}
 r_{k+1}&=\rho r_k,
 \notag\\
 H_{k+1}&=\rho^\beta H_k,
 \label{eq:unified-amplitude-recursion}\\
 \mathbf A_{s,r_k}(u,\ell_k;x_0)&\le H_k,
 \label{eq:unified-physical-amplitude}\\
 |a_{k+1}-a_k|+r_k|q_{k+1}-q_k|
 &\le C_\sharp r_kH_k,
 \label{eq:unified-affine-increments}\\
 \mathcal T_{\beta,G_*}(v_k,Q_k)&\le1.
 \label{eq:unified-tail-invariant}
\end{align}
Moreover,
\begin{equation}
 H_k\le C r_k^{\beta}.
 \label{eq:unified-amplitude-decay}
\end{equation}
The constants are independent of the normalized slopes \(Q_k\).
\end{proposition}

\begin{proof}
By exact homogeneity in \eqref{eq:common-intrinsic-profile}, the initial
condition \(\mathbf A_{s,r_0}(u,\ell_0;x_0)\le H_0\) is equivalent to
\(\mathfrak a_{s,Q_0}(v_0,0;1)\le1\).
Suppose the objects are constructed at level \(k\).  Apply Lemma
\ref{lem:affine-invariant-state-reproduction-main}.  Thus the
bounded-slope estimate is used when \(|Q_k|\le Q_{\rm th}\), and the
prescribed-scale large-slope estimate is used otherwise.  Write the
corresponding affine map as \(m_k=c_k+b_k\cdot X\), and set
\begin{equation*}
 \ell_{k+1}(x):=\ell_k(x)+r_kH_k
 m_k\left(\frac{x-x_0}{r_k}\right).
\end{equation*}
Since \(H_{k+1}/H_k=\eta=\rho^\beta\), Lemma
\ref{lem:affine-invariant-state-reproduction-main} gives simultaneously
\eqref{eq:unified-physical-amplitude} and
\eqref{eq:unified-tail-invariant}.  The common coefficient bound in the
two finite-scale estimates gives
\eqref{eq:unified-affine-increments}.

For clarity, the exact homogeneity behind the physical amplitude is
\begin{align*}
 \mathbf Z_{r_{k+1}}(u,\ell_{k+1};x_0)
 &=H_k^p\mathfrak Z(v_k,m_k;\rho),\\
 \mathbf E_{s,r_{k+1}}(u,\ell_{k+1};x_0)
 &=H_k^p\mathfrak B_{s,Q_k}(v_k,m_k;\rho),\\
 q_{k+1}&=H_k(Q_k+b_k),\qquad
 H_{k+1}=H_k\eta .
\end{align*}
Thus
\(\mathfrak a_{s,Q_k}(v_k,m_k;\rho)\le\eta\) is exactly
\(\mathbf A_{s,r_{k+1}}(u,\ell_{k+1};x_0)\le H_{k+1}\).

The new normalized remainder is
\begin{equation*}
 v_{k+1}(X)=\frac{1}{\rho\eta}
 [v_k(\rho X)-m_k(\rho X)].
\end{equation*}
Thus the same amplitude is used in both regimes.  Iterating the two
exact recursions gives
\[
 r_k=\rho^kr_0,\qquad H_k=\rho^{k\beta}H_0
      =H_0r_0^{-\beta}r_k^\beta,
\]
which proves \eqref{eq:unified-amplitude-decay}.  The affine increments
are summable because \(\sum_k r_kH_k<\infty\) and
\(\sum_kH_k<\infty\).
\end{proof}
\begin{corollary}[Campanato conclusion]
\label{cor:common-recursion-campanato}
Under the hypotheses of Proposition
\ref{prop:unified-variable-block-recursion}, the affine slopes converge
to a vector $q(x_0)$ and
\begin{equation}
 \left(\fint_{B_r(x_0)}
 |u-u(x_0)-q(x_0)\cdot(x-x_0)|^p\dd x\right)^{1/p}
 \le Cr^{1+\beta}
 \label{eq:common-recursion-campanato-decay}
\end{equation}
for all sufficiently small $r$.  Consequently
if the hypotheses hold with the same constants at every center in an
open set, then \(\nabla u\in C^{\beta}_{\rm loc}\) there.
\end{corollary}

\begin{proof}
By \eqref{eq:unified-affine-increments} and
\eqref{eq:unified-amplitude-decay},
\[
 |q_{k+1}-q_k|\le Cr_k^{\beta},
 \qquad
 |a_{k+1}-a_k|\le Cr_k^{1+\beta}.
\]
Both series converge because \(r_k=r_0\rho^k\).  The zero-order part
of \eqref{eq:unified-physical-amplitude}, together with the continuity
of \(u\), identifies \(\lim_{k\to\infty}a_k=u(x_0)\) and then gives
\eqref{eq:common-recursion-campanato-decay} first at the radii $r_k$ and
then at every intermediate radius, using
\(r_{k+1}=\rho r_k\).  The standard overlapping-ball
argument identifies the limiting slope with $\nabla u$.  More
explicitly, if \(d=|x_0-x_1|\) and a ball of radius \(4d\) stays in the
interior, compare the two limiting affine maps on the overlap of
\(B_{4d}(x_0)\) and \(B_{4d}(x_1)\).  Finite-dimensional norm
equivalence and \eqref{eq:common-recursion-campanato-decay} give
\[
 |q(x_0)-q(x_1)|\le Cd^\beta.
\]
This is the asserted \(C^\beta\) modulus.
\end{proof}

\begin{proof}[Proof of Theorem~\ref{thm:main}]
By homogeneity suppose that
$\|u\|_{L^\infty(\R^n)}\le1$.  The local H\"older theory first gives a
continuous representative.  For the homogeneous equation, the
weak--viscosity equivalence in
\cite[Proposition~1.5]{BiswasTopp} therefore permits the use of its
viscosity estimate.  Since \(sp>p-1\), Theorem~2.1 there, applied on
interior balls and combined by a finite covering of \(B_{3/2}\),
supplies
\begin{equation}
 \Lip(u;B_{3/2})\le L_0=L_0(n,p,s).
 \label{eq:main-initial-lipschitz}
\end{equation}
Use the dyadic radius and exponent \(\beta>0\) fixed in
\eqref{eq:affine-route-beta-bar-main}--\eqref{eq:affine-route-final-beta-main}.  Take \(r_0>0\) so small that
\(2R_*r_0\le1\).  Then every ball required by the bounded-slope
estimate, at the initial or any subsequent scale and with center in
\(B_{1/2}\), is contained in \(B_{3/2}\).

Fix \(x_0\in B_{1/2}\) and take \(\ell_0\equiv u(x_0)\).  Choose
\(H_0=C(1+L_0)\), with the structural factor
large enough to normalize both the intrinsic amplitude and the exterior
state.  From \eqref{eq:main-initial-lipschitz},
\begin{align*}
& \fint_{B_{2r_0}(x_0)}
 \left|\frac{u-u(x_0)}{r_0}\right|^p\dd x
 \le C L_0^p,\\
 &(1-s)r_0^{sp-p-n}
 \iint_{B_{2r_0}(x_0)^2}
 \frac{|\delta u(x,y)|^p}{|x-y|^{n+sp}}\dd x\!\dd y
 \le C L_0^p.
\end{align*}
Thus \(\mathbf A_{s,r_0}(u,u(x_0);x_0)\le H_0\).

For the initial normalized remainder
\[
 v_0(X)=\frac{u(x_0+r_0X)-u(x_0)}{r_0H_0},
\]
the competitor zero in the affine excess and the Lipschitz estimate
give \(\mathbf c_\ell(v_0;j)\le CL_0/H_0\) as long as the
corresponding physical ball stays in \(B_{3/2}\).  On the remaining
balls, boundedness of \(u\) gives
\(\mathbf c_\ell(v_0;j)\le C/(r_0H_0 2^j)\).  Consequently
\(\mathscr C_{\ell,\beta}(v_0)\le C(1+L_0)/H_0\).
Since the initial normalized tilt is zero, the relative flux is estimated by
\(|\delta v_0|^{p-1}\).  The annuli inside the Lipschitz region give the
convergent series
\[
 C\left(\frac{L_0}{H_0}\right)^{p-1}
 \sum_{m\ge0}2^{-m[sp-(p-1)]},
\]
whereas the part beyond that region is bounded by
\[
 C(r_0H_0)^{1-p}r_0^{sp}
 =C H_0^{1-p}r_0^{sp-p+1}.
\]
These estimates apply to both separated components in
\(\mathcal T^{\rm pt}_{s,0}(v_0)\); changing the fixed inner and outer
radii only changes the structural constant.
The affine-excess component scales as \(H_0^{-1}\), whereas the
relative-flux component scales as \(H_0^{1-p}\).  Increasing the
structural factor in \(H_0\), and then
decreasing \(r_0\) if necessary, gives
\[
 \mathcal T_{\beta,G_*}(v_0,0)\le1.
\]
At every subsequent scale, the alternatives
\(|Q_k|\le Q_{\rm th}\) and \(|Q_k|>Q_{\rm th}\) exhaust the normalized
slopes.  Lemma~\ref{lem:affine-invariant-state-reproduction-main}
uses the corresponding finite-scale estimate without changing the common
amplitude or exterior state.  Proposition
\ref{prop:unified-variable-block-recursion} therefore constructs the
affine maps and gives \(H_k\le Cr_k^\beta\), with constants independent
of the center.
Corollary~\ref{cor:common-recursion-campanato} yields the affine
Campanato expansion at every point of $B_{1/2}$.  Its overlapping-ball
argument gives the gradient modulus when
\(|x_0-x_1|\le c r_0\).  For larger distances, the initial Lipschitz
bound absorbs the estimate:
\[
 |\nabla u(x_0)-\nabla u(x_1)|
 \le2L_0\le C r_0^{-\beta}|x_0-x_1|^\beta.
\]
This proves \eqref{eq:fixed-order-gradient-holder-estimate}, with
\(\alpha_0=\beta\).
\end{proof}

\begin{remark}[Uniformity with respect to the normalized slope]
\label{rem:logical-content-common-splice}
The affine-invariant excess absorbs every accumulated affine map and is
transported exactly by \eqref{eq:exact-affine-weight-cancellation}.
The bounded-slope estimate follows from the nonlinear affine-Campanato
Liouville theorem, whereas the large-slope estimate follows from
finite-scale convergence to the anisotropic stable equation.  After
fixing a common contraction, the amplitude update
\eqref{eq:unified-amplitude-recursion} and the state
\eqref{eq:common-tail-state} hold in both regimes.  The
exponent \(\beta\) is limited by the nonexplicit compactness number
\(\beta_*\) and by the strict tail margin \(sp-p+1\).  Its choice does
not require an additional annular renormalization constant.
\end{remark}


\subsection*{Acknowledgments}
This work was supported by the National Natural Science Foundation of China (No. 12471128).

\subsection*{Conflict of interest}
The author declares that there is no conflict of interest.

\subsection*{Data availability}
No datasets were generated or analyzed for this work.


\begin{thebibliography}{99}

\bibitem{AdamsTrace}
D.~R. Adams,
\emph{A trace inequality for generalized potentials},
Studia Math. \textbf{48} (1973), 99--105.

\bibitem{AlbericoCianchiPickSlavikova}
A.~Alberico, A.~Cianchi, L.~Pick and L.~Slav\'ikov\'a,
\emph{Fractional Orlicz--Sobolev embeddings},
J. Math. Pures Appl. (9) \textbf{149} (2021), 216--253.

\bibitem{BiswasTopp}
A.~Biswas and E.~Topp,
\emph{Lipschitz regularity of fractional $p$-Laplacian},
Ann. PDE \textbf{11} (2025), Paper No.~27, 43~pp.

\bibitem{BoegeleinDuzaarLiaoMolicaBisciServadei}
V.~B\"ogelein, F.~Duzaar, N.~Liao, G.~Molica Bisci and R.~Servadei,
\emph{Gradient regularity for $(s,p)$-harmonic functions},
Calc. Var. Partial Differential Equations \textbf{64} (2025),
Paper No.~253, 55~pp.

\bibitem{BoegeleinDuzaarLiaoMolicaBisciServadeiJFA}
V.~B\"ogelein, F.~Duzaar, N.~Liao, G.~Molica Bisci and R.~Servadei,
\emph{Regularity for the fractional $p$-Laplace equation},
J. Funct. Anal. \textbf{289} (2025), Paper No.~111078, 69~pp.

\bibitem{BoegeleinDuzaarLiaoMolicaBisciServadeiHigher}
V.~B\"ogelein, F.~Duzaar, N.~Liao, G.~Molica Bisci and R.~Servadei,
\emph{Higher regularity theory for $(s,p)$-harmonic functions},
Rend. Lincei Mat. Appl. \textbf{35} (2024), 311--321.

\bibitem{BrascoLindgren}
L.~Brasco and E.~Lindgren,
\emph{Higher Sobolev regularity for the fractional $p$-Laplace equation
in the superquadratic case},
Adv. Math. \textbf{304} (2017), 300--354.

\bibitem{BrascoLindgrenSchikorra}
L.~Brasco, E.~Lindgren and A.~Schikorra,
\emph{Higher H\"older regularity for the fractional $p$-Laplacian in the
superquadratic case},
Adv. Math. \textbf{338} (2018), 782--846.

\bibitem{DeFilippisMingione}
C.~De Filippis and G.~Mingione,
\emph{Gradient regularity in mixed local and nonlocal problems},
Math. Ann. \textbf{388} (2024), 261--328.

\bibitem{DiCastroKuusiPalatucci}
A.~Di Castro, T.~Kuusi and G.~Palatucci,
\emph{Local behavior of fractional $p$-minimizers},
Ann. Inst. H. Poincar\'e C Anal. Non Lin\'eaire \textbf{33} (2016),
1279--1299.

\bibitem{DiCastroKuusiPalatucciHarnack}
A.~Di Castro, T.~Kuusi and G.~Palatucci,
\emph{Nonlocal Harnack inequalities},
J. Funct. Anal. \textbf{267} (2014), 1807--1836.

\bibitem{DieningKimLeeNowak}
L.~Diening, K.~Kim, H.-S.~Lee and S.~Nowak,
\emph{Higher differentiability for the fractional $p$-Laplacian},
Math. Ann. \textbf{391} (2025), 5631--5693.

\bibitem{DieningKimLeeNowakGradient}
L.~Diening, K.~Kim, H.-S.~Lee and S.~Nowak,
\emph{Nonlinear nonlocal potential theory at the gradient level},
J. Eur. Math. Soc. (2025), published online first.

\bibitem{DieningNowak}
L.~Diening and S.~Nowak,
\emph{Calder\'on--Zygmund estimates for the fractional $p$-Laplacian},
Ann. PDE \textbf{11} (2025), Paper No.~6, 69~pp.

\bibitem{GarainLindgren}
P.~Garain and E.~Lindgren,
\emph{Higher H\"older regularity for the fractional $p$-Laplace equation
in the subquadratic case},
Math. Ann. \textbf{390} (2024), 5753--5792.

\bibitem{GJS}
\begingroup\emergencystretch=1em
D.~Giovagnoli, D.~Jesus and L.~Silvestre,
\emph{$C^{1+\alpha}$ regularity for fractional $p$-harmonic functions},
Preprint, arXiv:2509.26565, 2025.\par
\endgroup

\bibitem{IannizzottoMosconi}
A.~Iannizzotto and S.~Mosconi,
\emph{Fine boundary regularity for the singular fractional $p$-Laplacian},
J. Differential Equations \textbf{412} (2024), 322--379.

\bibitem{IannizzottoMosconiSquassina}
A.~Iannizzotto, S.~Mosconi and M.~Squassina,
\emph{Global H\"older regularity for the fractional $p$-Laplacian},
Rev. Mat. Iberoam. \textbf{32} (2016), 1353--1392.

\bibitem{JarohsWethSMP}
S.~Jarohs and T.~Weth,
\emph{On the strong maximum principle for nonlocal operators},
Math. Z. \textbf{293} (2019), 81--111.

\bibitem{KKL}
J.~Korvenp\"a\"a, T.~Kuusi and E.~Lindgren,
\emph{Equivalence of solutions to fractional $p$-Laplace type
equations}, J. Math. Pures Appl. (9) \textbf{132} (2019), 1--26.

\bibitem{KuusiMingioneSireMeasure}
T.~Kuusi, G.~Mingione and Y.~Sire,
\emph{Nonlocal equations with measure data},
Comm. Math. Phys. \textbf{337} (2015), 1317--1368.

\bibitem{KuusiMingioneSireSelf}
T.~Kuusi, G.~Mingione and Y.~Sire,
\emph{Nonlocal self-improving properties},
Anal. PDE \textbf{8} (2015), 57--114.

\bibitem{KuusiMingioneSireSurvey}
T.~Kuusi, G.~Mingione and Y.~Sire,
\emph{Regularity issues involving the fractional $p$-Laplacian},
in: \emph{Recent Developments in Nonlocal Theory},
De Gruyter, Berlin, 2018, 303--334.

\bibitem{KuusiNowakSire}
T.~Kuusi, S.~Nowak and Y.~Sire,
\emph{Gradient regularity and first-order potential estimates for a class
of nonlocal equations},
to appear in Amer. J. Math., arXiv:2212.01950.

\bibitem{LindgrenViscosity}
E.~Lindgren,
\emph{H\"older estimates for viscosity solutions of equations of
fractional $p$-Laplace type},
NoDEA Nonlinear Differential Equations Appl. \textbf{23} (2016),
Paper No.~55, 18~pp.

\bibitem{RosOtonSerraStable}
X.~Ros-Oton and J.~Serra,
\emph{Regularity theory for general stable operators},
J. Differential Equations \textbf{260} (2016), 8675--8715.

\bibitem{ZhangPartial}
C.~Zhang,
\emph{Towards gradient H\"older regularity for singular fractional
$p$-Laplace equations},
arXiv:2608.16243, 2026.

\end{thebibliography}

\end{document}
